% Generated mechanically from frozen input p01/ko/topology.tex.
% Do not edit this generated file; edit the bound locale source instead.

% OK, start here.
%

\setcounter{chapter}{4}
\chapter{위상수학}



\phantomsection
\label{topology-section-phantom}


\section{서론}
\label{topology-section-introduction}

\noindent
이 문서에서는 기초 위상수학을 설명한다.
참고문헌으로 \cite{Engelking}이 있다.

\section{기본 개념}
\label{topology-section-topology-basic}

\noindent
다음은 위상수학의 기본 개념들을 나열한 것이다. 이 가운데 일부는 뒤의
본문에서 더 자세히 논의하고, 일부는 이 목록에서 정의하지만, 나머지는
기본적인 것으로 간주하여 정의하지 않는다. 이탤릭체로 표시한 개념 대부분이
익숙하지 않다면, 계속 읽기 전에 위상수학 입문서를 살펴볼 것을 권한다.

\begin{enumerate}
\item
\label{topology-item-space}
$X$는 {\it 위상공간}이다.
\item
\label{topology-item-point}
$x\in X$는 {\it 점}이다.
\item
\label{topology-item-locally-closed}
$E \subset X$는 {\it 국소닫힌} 부분집합이다.
\item
\label{topology-item-closed-point}
$x\in X$는 {\it 닫힌점}이다.
\item
\label{topology-item-dense}
$E \subset X$는 {\it 조밀한} 부분집합이다.
\item
\label{topology-item-continuous}
$f : X_1 \to X_2$는 {\it 연속}이다.
\item
\label{topology-item-upper-semi-continuous}
확장실수 함수 $f : X \to \mathbf{R} \cup \{\infty, -\infty\}$가
모든 $a \in \mathbf{R}$에 대하여 $\{x \in X \mid f(x) < a\}$가
열린집합이면 {\it 위 준연속}이라고 한다.
\item
\label{topology-item-lower-semi-continuous}
확장실수 함수 $f : X \to \mathbf{R} \cup \{\infty, -\infty\}$가
모든 $a \in \mathbf{R}$에 대하여 $\{x \in X \mid f(x) > a\}$가
열린집합이면 {\it 아래 준연속}이라고 한다.
\item 공간들의 연속사상 $f : X \to Y$가 임의의 열린집합
$U \subset X$에 대하여 $f(U)$를 $Y$의 열린집합으로 보내면
{\it 열린사상}이라고 한다.
\item 공간들의 연속사상 $f : X \to Y$가 임의의 닫힌집합
$Z \subset X$에 대하여 $f(Z)$를 $Y$의 닫힌집합으로 보내면
{\it 닫힌사상}이라고 한다.
\item
\label{topology-item-neighbourhood}
{\it $x \in X$의 근방}은 $x$를 포함하는 열린 부분집합을 포함하는
임의의 부분집합 $E \subset X$이다.
\item
\label{topology-item-induced-topology}
부분집합 $E \subset X$ 위의 {\it 유도위상}.
\item
\label{topology-item-covering}
$\mathcal{U} : U = \bigcup_{i \in I} U_i$는 $U$의
{\it 열린 덮개}이다(주의: 임의의 $U_i$가 공집합인 것을 허용하며,
$U$가 공집합인 경우에는 $I$ 자체가 공집합인 것도 허용한다).
\item
\label{topology-item-subcover}
(\ref{topology-item-covering})과 같은 덮개의 {\it 부분덮개}란
$I' \subset I$인 열린 덮개
$\mathcal{U}' : U = \bigcup_{i \in I'} U_i$를 말한다.
\item
\label{topology-item-refinement}
열린 덮개 $\mathcal{V}$가 열린 덮개 $\mathcal{U}$의
{\it 세분}이라는 것은($\mathcal{V} : U = \bigcup_{j \in J} V_j$이고
$\mathcal{U} : U = \bigcup_{i \in I} U_i$일 때), 각 $V_j$가 어느 한
$U_i$ 안에 완전히 포함된다는 뜻이다.
\item
\label{topology-item-fundamental-system}
{\it $\{ E_i \}_{i \in I}$는 $X$에서 $x$의 근방들의 기본계이다}.
\item
\label{topology-item-Hausdorff}
위상공간 $X$에서 서로 다른 임의의 두 점 $x, y \in X$마다
$x \in U$, $y \in V$인 서로소 열린집합 $U, V \subset X$가 존재할 때,
그리고 오직 그때에만 그 공간을 {\it 하우스도르프} 또는 {\it 분리된}
공간이라고 한다.
\item 두 위상공간의 {\it 곱}.
\label{topology-item-product}
\item
\label{topology-item-fibre-product}
연속사상 $f : X \to Y$와 $g : Z \to Y$의 쌍의
{\it 섬유곱 $X \times_Y Z$}.
\item
\label{topology-item-discrete-indiscrete}
집합 위의 {\it 이산위상}과 {\it 비이산위상}.
\item 기타.
\end{enumerate}



\section{하우스도르프 공간}
\label{topology-section-Hausdorff}

\noindent
위상공간들의 범주는 유한 곱을 갖는다.

\begin{lemma}
\label{topology-lemma-Hausdorff}
$X$를 위상공간이라 하자. 다음은 서로 동치이다.
\begin{enumerate}
\item $X$는 하우스도르프 공간이다.
\item 대각선 $\Delta(X) \subset X \times X$는 닫혀 있다.
\end{enumerate}
\end{lemma}

\begin{proof}
$X$가 하우스도르프라고 가정하자. $(x, y) \not \in \Delta(X)$, 즉
$x \neq y$라 하자. $x \in U$이고 $y \in V$인 $X$의 서로소 열린집합
$U$와 $V$가 존재한다. 따라서 $U \times V \subset
(X \times X) \setminus \Delta(X) $이다. 이는
$ (X \times X) \setminus \Delta(X)$가 $ X \times X$의 열린집합임을
보이며, 이는 대각선 $\Delta(X) \subset X \times X$가 $X \times X$에서
닫혀 있다는 것과 동치이다. 역도 비슷하다. 여집합
$(X \times X) \setminus \Delta(X)$는 정확히 $ x \neq y$인
$(x, y) \in X\times X$들로 이루어진다. 따라서 $ \Delta(X)
\subset X \times X$가 닫혀 있으면, 곱위상의 정의에 의해 그러한 모든
$ (x, y)$에 대하여 $(x, y) \in U \times V$이고
$(U \times V)\cap \Delta(X) = \emptyset$인 열린집합 $U, V \subset X$가
존재한다. 다시 말해 $x \in U$, $y \in V$이고
$U \cap V = \emptyset$이다.
\end{proof}







\begin{lemma}
\label{topology-lemma-graph-closed}
\begin{slogan}
하우스도르프 공간으로 가는 사상의 그래프는 닫혀 있다.
\end{slogan}
$f : X \to Y$를 위상공간들의 연속사상이라 하자.
$Y$가 하우스도르프이면 $f$의 그래프는 $X \times Y$에서 닫혀 있다.
\end{lemma}

\begin{proof}
그래프는 사상 $X \times Y \to Y \times Y$에 의한 대각선의 역상이다.
따라서 보조정리는 Lemma \ref{topology-lemma-Hausdorff}에서 따른다.
\end{proof}

\begin{lemma}
\label{topology-lemma-section-closed}
$f : X \to Y$를 위상공간들의 연속사상이라 하자.
$s : Y \to X$를 $f \circ s = \text{id}_Y$를 만족하는 연속사상이라 하자.
$X$가 하우스도르프이면 $s(Y)$는 닫혀 있다.
\end{lemma}

\begin{proof}
$s(Y) = \{x \in X \mid x = s(f(x))\}$이므로
Lemma \ref{topology-lemma-Hausdorff}에서 따른다.
\end{proof}

\begin{lemma}
\label{topology-lemma-fibre-product-closed}
$X \to Z$와 $Y \to Z$를 위상공간들의 연속사상이라 하자.
$Z$가 하우스도르프이면 $X \times_Z Y$는 $X \times Y$에서 닫혀 있다.
\end{lemma}

\begin{proof}
$X \times_Z Y$는 $X \times Y \to Z \times Z$에 의한
$\Delta(Z)$의 역상이므로 Lemma \ref{topology-lemma-Hausdorff}에서 따른다.
\end{proof}




\section{분리 사상}
\label{topology-section-separated}

\noindent
정의와 몇 가지 간단한 보조정리만 제시한다.

\begin{definition}
\label{topology-definition-separated}
위상공간들의 연속사상 $f : X \to Y$에 대하여 대각사상
$\Delta : X \to X \times_Y X$가 닫힌사상일 때, 그리고 오직 그때에만
그 사상을 {\it 분리되었다}고 한다.
\end{definition}

\begin{lemma}
\label{topology-lemma-separated}
$f : X \to Y$를 위상공간들의 연속사상이라 하자.
다음은 서로 동치이다.
\begin{enumerate}
\item $f$는 분리되어 있다.
\item $\Delta(X) \subset X \times_Y X$는 닫힌 부분집합이다.
\item $Y$의 같은 점으로 가는 서로 다른 점 $x, x' \in X$가 주어지면,
$x$와 $x'$의 서로소 열린 근방들이 존재한다.
\end{enumerate}
\end{lemma}

\begin{proof}
$f$가 분리되어 있으면 Definition \ref{topology-definition-separated}에 의해
$\Delta$는 닫힌사상이다. $X$가 $X$에서 닫혀 있으므로
$\Delta(X)$는 $ X\times_Y X $에서 닫혀 있다. 따라서 (1)은 (2)를 함의한다.

\medskip\noindent
$\Delta(X) \subset X \times_Y X$가 닫힌 부분집합이라고 가정하고 그 열린
여집합을 $U$라 하자. 그러면 $W \cap (X \times_Y X) = U$인 열린집합
$W \subset X \times X$가 존재한다. 그런데 곱위상의 정의에 의해
$(x, x') \in W \cap (X \times_Y X)$이면, $x \in V$, $x' \in V'$이고
$V \times V' \subset W$인 $X$의 열린집합 $V$와 $V'$가 존재한다.
만일 $V \cap V' \neq \emptyset$이면 $z \in V \cap V'$인 점이 존재한다.
그러나 $(z, z) \in X \times_Y X$이므로
$(z,z) \in (V \times V') \cap (X\times_Y X) \subset U$가 되어 모순이다.
따라서 $V \cap V' = \emptyset $이고, $x$와 $x'$의 서로소 열린 근방 두 개를
얻는다. 이로써 (2)가 (3)을 함의함을 보였다.

\medskip\noindent
마지막으로 $Y$의 같은 점으로 가는 서로 다른 점 $x, x' \in X$가 주어질
때마다 $x$와 $x'$의 서로소 열린 근방들이 존재한다고 가정하자.
$F$를 $X$의 닫힌집합이라 하자. $\Delta(F)$가 $X \times_Y X$의 닫힌
부분집합임을 보이겠다. $(x, x') \in X \times_Y X$가 $\Delta(F)$에
속하지 않는 점이라 하자. 그러면 $x \not = x'$이거나 $x \not \in F$이다.
첫째 경우에는 $x, x'$의 서로소 열린 근방 $V, V' \subset X$를 택한다.
그러면 $(V \times V') \cap X \times_Y X$는 $\Delta(F)$와 만나지 않는
$(x, x')$의 열린 근방이다. 둘째 경우에는
$((X \setminus F) \times X) \cap X \times_Y X$가 $\Delta(F)$와 만나지
않는 $(x, x')$의 열린 근방이다. 따라서 (3)이 (1)을 함의함을 보였다.
\end{proof}

\begin{lemma}
\label{topology-lemma-from-hausdorff}
$f : X \to Y$를 위상공간들의 연속사상이라 하자.
$X$가 하우스도르프이면 $f$는 분리되어 있다.
\end{lemma}

\begin{proof}
$X \times X$에서 $\Delta(X)$가 닫혀 있으면 $X \times_Y X$에서도
$\Delta(X)$가 닫혀 있으므로 Lemmas \ref{topology-lemma-separated}와
\ref{topology-lemma-Hausdorff}에서 바로 따른다.
\end{proof}

\begin{lemma}
\label{topology-lemma-base-change-separated}
$f : X \to Y$와 $Y' \to Y$를 위상공간들의 연속사상이라 하자.
$f$가 분리되어 있으면 $f' : Y' \times_Y X \to Y'$도 분리되어 있다.
\end{lemma}

\begin{proof}
Lemma \ref{topology-lemma-separated}의 조건 (2)에서 따른다. 실제로
$X' = Y' \times_Y X$라 하면 섬유곱 $X' \times_{Y'} X'$ 안의 대각선
$\Delta(X')$는 $X \times_Y X$ 안의 $\Delta(X)$의 역상이다.
\end{proof}




\section{기저}
\label{topology-section-bases}

\noindent
위상공간의 기저에 관한 기본 내용을 다룬다.

\begin{definition}
\label{topology-definition-base}
$X$를 위상공간이라 하자. $X$의 부분집합들의 모음 $\mathcal{B}$가 다음
조건들을 만족하면 이를 {\it $X$ 위의 위상기저} 또는
{\it $X$ 위의 위상의 기저}라고 한다.
\begin{enumerate}
\item 모든 원소 $B \in \mathcal{B}$는 $X$에서 열려 있다.
\item 임의의 열린집합 $U \subset X$와 임의의 $x \in U$에 대하여,
$x \in B \subset U$를 만족하는 원소 $B \in \mathcal{B}$가 존재한다.
\end{enumerate}
\end{definition}

\noindent
다음 보조정리를 위상을 정의하는 데 사용하기도 한다.

\begin{lemma}
\label{topology-lemma-make-base}
$X$를 집합, $\mathcal{B}$를 부분집합들의 모음이라 하자.
$X = \bigcup_{B \in \mathcal{B}} B$라고 가정하자. 또한
$B_1, B_2 \in \mathcal{B}$이고 $x \in B_1 \cap B_2$이면
$x \in B_3 \subset B_1 \cap B_2$인 $B_3 \in \mathcal{B}$가 존재한다고
가정하자. 그러면 $\mathcal{B}$가 기저가 되는 $X$ 위의 위상은 유일하게
존재한다.
\end{lemma}

\begin{proof}
$\mathcal B$의 원소들의 합집합으로 쓸 수 있는 $X$의 부분집합들의 집합을
$\sigma(\mathcal B)$라 하자. $\sigma(\mathcal{B})$가 위상임을 보이겠다.
실제로 공집합은 (빈 합집합으로서) $\sigma(\mathcal{B})$의 원소이고,
$X$는 ($\mathcal{B}$의 모든 원소의 합집합으로서)
$\sigma(\mathcal{B})$의 원소이다. $\sigma(\mathcal{B})$가 합집합에 대해
닫혀 있음은 분명하다. 마지막으로 $U, V \in \sigma(\mathcal{B})$라 하고,
모든 $i \in I$와 $j \in J$에 대하여 $U_i, V_j \in \mathcal{B}$인 표현
$U = \bigcup_{i \in I} U_i$와 $V = \bigcup_{j \in J} V_j$를 택하자. 그러면
$$
U \cap V = \bigcup\nolimits_{i \in I, j \in J}
U_i \cap V_j
$$
이다. 보조정리의 가정에 의해 $U_i \cap V_j \in \sigma(\mathcal{B})$이므로
$U \cap V$도 그러하다. 따라서 $\sigma(\mathcal{B})$는 위상이다.
Definition \ref{topology-definition-base}의 성질 (1), (2)는 이 위상에 대해 바로
성립한다. 유일성을 보이기 위해 $\mathcal B$가 기저인 $X$ 위의 위상을
$\mathcal T$라 하자. Definition \ref{topology-definition-base}의 (1)에 의해
$\mathcal{B}$의 모든 원소가 $\mathcal{T}$에 속하므로
$\sigma(\mathcal{B}) \subset \mathcal{T}$이다. 반대로 Definition
\ref{topology-definition-base}의 (2)에 의해 $\mathcal{T}$의 모든 원소는
$\mathcal{B}$의 원소들의 합집합이다. 즉
$\mathcal{T} \subset \sigma(\mathcal{B})$이다. 이로써 증명이 끝난다.
\end{proof}

\begin{lemma}
\label{topology-lemma-refine-covering-basis}
$X$를 위상공간이라 하자.
$\mathcal{B}$를 $X$ 위의 위상의 기저라 하자.
$\mathcal{U} : U = \bigcup_i U_i$를 $U \subset X$의 열린 덮개라 하자.
그러면 각 $V_j$가 기저 $\mathcal{B}$의 원소이고 $\mathcal{U}$를 세분하는
열린 덮개 $U = \bigcup V_j$가 존재한다.
\end{lemma}

\begin{proof}
$ x \in U = \bigcup_{i\in I} U_i $이면
$ x \in U_{i_x} $인 $ i_x \in I $가 존재한다. 따라서
$ x \in B_{i_x} \subset U_{i_x}$를 만족하는
$ B_{i_x} \in \mathcal{B}$가 존재한다. 이제
$J = \{i_x | x \in U\}$라 두고, $j = i_x \in J$에 대하여
$V_j = B_{i_x}$로 두자. 그러면 $\{V_j\}_{j \in J}$는 원하는 $U$의
열린 덮개이다.
\end{proof}

\begin{definition}
\label{topology-definition-subbase}
$X$를 위상공간이라 하자. $X$의 부분집합들의 모음 $\mathcal{B}$의
원소들의 유한 교집합들이 $X$ 위의 위상의 기저를 이룰 때,
$\mathcal{B}$를 {\it $X$ 위의 위상부분기저} 또는
{\it $X$ 위의 위상의 부분기저}라고 한다.
\end{definition}

\noindent
특히 $\mathcal{B}$의 모든 원소는 열려 있다.

\begin{lemma}
\label{topology-lemma-subbase}
$X$를 집합이라 하자. $X$의 부분집합들의 임의의 모음 $\mathcal{B}$가
주어지면, $\mathcal{B}$가 부분기저가 되는 $X$ 위의 위상은 유일하게
존재한다.
\end{lemma}

\begin{proof}
관례에 따라 $\bigcap_\emptyset B = X $이다. 따라서 $\mathcal B$의
원소들의 유한 교집합들의 집합에 Lemma \ref{topology-lemma-make-base}를 적용하면 된다.
\end{proof}

\begin{lemma}
\label{topology-lemma-create-map-from-subcollection}
$X$를 위상공간이라 하고, $\mathcal{B}$를 $X$의 열린집합들의 모음이라 하자.
$X = \bigcup_{U \in \mathcal{B}} U$라고 가정하고,
$U, V \in \mathcal{B}$에 대하여
$U \cap V = \bigcup_{W \in \mathcal{B}, W \subset U \cap V} W$라고 하자.
그러면 다음을 만족하는 위상공간들의 연속사상 $f : X \to Y$가 존재한다.
\begin{enumerate}
\item $U \in \mathcal{B}$이면 상 $f(U)$는 열려 있다.
\item $U \in \mathcal{B}$이면 $f^{-1}(f(U)) = U$이다.
\item 열린집합들 $f(U)$, $U \in \mathcal{B}$는 $Y$ 위의 위상의 기저를
이룬다.
\end{enumerate}
\end{lemma}

\begin{proof}
$X$의 점들 위의 동치관계 $\sim$을 다음 규칙으로 정의한다.
$$
x \sim y \Leftrightarrow
(\forall U \in \mathcal{B} : x \in U \Leftrightarrow y \in U)
$$
$Y$를 동치류들의 집합, $f : X \to Y$를 자연스러운 사상이라 하자.
구성에 의해 (2)가 성립한다. $\mathcal{B}$에 관한 가정들은
부분집합들 $f(U)$, $U \in \mathcal{B}$에 대하여 Lemma
\ref{topology-lemma-make-base}의 가정들과 정확히 일치한다. 따라서 (3)을 만족하는
$Y$ 위의 위상이 유일하게 존재한다. 그러면 (1)도 분명하다.
\end{proof}


\section{몫위상 사상}
\label{topology-section-submersive}

\noindent
$X$가 위상공간이고 $E \subset X$가 부분집합이면 보통 $E$에
{\it 유도위상}을 준다.

\begin{lemma}
\label{topology-lemma-induced}
$X$를 위상공간이라 하자. $Y$를 집합이라 하고
$f : Y \to X$를 집합들의 단사사상이라 하자. $Y$ 위의 유도위상은
다음 각각의 조건으로 특징지어지는 위상이다.
\begin{enumerate}
\item $f$가 연속이 되게 하는 $Y$ 위의 가장 약한 위상이다.
\item $Y$의 열린 부분집합들은 열린집합 $U \subset X$에 대한
$f^{-1}(U)$들이다.
\item $Y$의 닫힌 부분집합들은 닫힌집합 $Z \subset X$에 대한
$f^{-1}(Z)$들이다.
\end{enumerate}
\end{lemma}

\begin{proof}
집합 $\mathcal{T} = \{f^{-1}(U) | U \subset X \text{ open}\}$은
$Y$ 위의 위상이다. 먼저 $\emptyset = f^{-1}(\emptyset)$이고
$f^{-1}(X) = Y$이다. 따라서 $\mathcal{T}$는 $\emptyset$과 $Y$를 포함한다.

\medskip\noindent
이제 $V_i \in \mathcal{T}$인 열린 부분집합들의 모음을
$\{V_i\}_{i \in I}$라 하고, $U_i$가 $X$의 열린 부분집합이 되도록
$V_i = f^{-1}(U_i)$로 쓰자. 그러면
$$
\bigcup\nolimits_{i\in I} V_i =
\bigcup\nolimits_{i\in I} f^{-1}(U_i) =
f^{-1}\left(\bigcup\nolimits_{i\in I} U_i\right)
$$
이다. $\bigcup_{i\in I} U_i$가 $X$에서 열려 있으므로
$\bigcup_{i\in I} V_i \in \mathcal{T}$이다. 이제
$V_1, V_2 \in \mathcal{T}$라 하자. $V_1 = f^{-1}(U_1)$이고
$V_2 = f^{-1}(U_2)$인 $X$의 열린집합 $U_1, U_2$가 존재한다. 그러면
$$
V_1 \cap V_2 = f^{-1}(U_1) \cap f^{-1}(U_2) = f^{-1}(U_1 \cap U_2)
$$
이다. $U_1 \cap U_2$가 $X$에서 열려 있으므로
$V_1 \cap V_2 \in \mathcal{T}$이다.

\medskip\noindent
연속사상의 정의에 따라 $f$가 연속이 되는 $Y$ 위의 임의의 위상은
$\mathcal{T}$를 포함한다. 따라서 $\mathcal{T}$는 실제로 $f$가 연속이
되게 하는 $Y$ 위의 가장 약한 위상이다. 이로써 (1)과 (2)가 동치임을
보였다.

\medskip\noindent
(2)와 (3)의 동치는 모든 부분집합 $E \subset X$에 대한 등식
$f^{-1}(X \setminus E) = Y \setminus f^{-1}(E)$에서 따른다.
\end{proof}

\noindent
쌍대적으로 $X$가 위상공간이고 $X \to Y$가 집합들의 전사이면,
$Y$에는 {\it 몫위상}을 줄 수 있다.

\begin{lemma}
\label{topology-lemma-quotient}
$X$를 위상공간이라 하자. $Y$를 집합이라 하고 $f : X \to Y$를
집합들의 전사사상이라 하자. $Y$ 위의 몫위상은 다음 각각의 조건으로
특징지어지는 위상이다.
\begin{enumerate}
\item $f$가 연속이 되게 하는 $Y$ 위의 가장 강한 위상이다.
\item $Y$의 부분집합 $V$가 열려 있을 필요충분조건은
$f^{-1}(V)$가 열려 있는 것이다.
\item $Y$의 부분집합 $Z$가 닫혀 있을 필요충분조건은
$f^{-1}(Z)$가 닫혀 있는 것이다.
\end{enumerate}
\end{lemma}

\begin{proof}
집합 $\mathcal{T} = \{V\subset Y | f^{-1}(V) \text{ is open}\}$은
$Y$ 위의 위상이다. 먼저 $\emptyset = f^{-1}(\emptyset)$이고
$f^{-1}(Y) = X$이다. 따라서 $\mathcal{T}$는 $\emptyset$과 $Y$를 포함한다.

\medskip\noindent
$(V_i)_{i \in I}$를 $V_i \in \mathcal{T}$인 원소족이라 하자. 그러면
$$
\bigcup\nolimits_{i\in I} f^{-1}(V_i) =
f^{-1}\left(\bigcup\nolimits_{i\in I} V_i\right)
$$
이다. $\bigcup_{i\in I}f^{-1}(V_i)$가 $X$에서 열려 있으므로
$\bigcup_{i\in I} V_i \in \mathcal{T}$이다. 또한
$V_1, V_2 \in \mathcal{T}$이면
$$
f^{-1}(V_1) \cap f^{-1}(V_2) = f^{-1}(V_1 \cap V_2)
$$
이다. $f^{-1}(V_1) \cap f^{-1}(V_2)$가 $X$에서 열려 있으므로
$V_1 \cap V_2 \in \mathcal{T}$이다.

\medskip\noindent
마지막으로 연속함수의 정의에 따라 $f$가 연속이 되는 $Y$ 위의 위상은
$\mathcal{T}$에 포함된다. 따라서 $\mathcal{T}$는 $f$가 연속이 되게 하는
$Y$ 위의 가장 강한 위상이다. 이로써 (1)과 (2)가 동치임을 보였다.

\medskip\noindent
마지막으로 (2)와 (3)의 동치는 모든 부분집합 $E \subset X$에 대한
% P01-E1315: frozen quotient-complement formula has domain/codomain symbols reversed.
$f^{-1}(X\setminus E) = Y
\setminus f^{-1}(E)$에서 따른다.
\end{proof}

\noindent
$f : X \to Y$를 위상공간들의 연속사상이라 하자. 그러면 집합들의 사상으로서
$X \to f(X) \to Y$라는 분해를 얻는다. $f(X)$에는 전사사상
$X \to f(X)$에서 오는 몫위상이나 단사사상 $f(X) \to Y$에서 오는
유도위상을 줄 수 있다. 사상
$$
(f(X), \text{quotient topology})
\longrightarrow
(f(X), \text{induced topology})
$$
은 연속이다.

\begin{definition}
\label{topology-definition-submersive}
$f : X \to Y$를 위상공간들의 연속사상이라 하자.
\begin{enumerate}
\item $f(X)$ 위의 유도위상과 몫위상이 일치하면 $f$를
{\it 위상공간들의 엄밀한 사상}이라고 한다(위의 논의를 보라).
\item $f$가 전사이고 엄밀하면 $f$를 {\it 몫위상 사상}
이라고 한다.\footnote{이는 미분다양체 사이의 침강(submersion) 개념과는
매우 다르다! “몫위상 사상” 대신 “엄밀하고 전사인 사상”이라고 쓰는 편이
아마도 좋다.}
\end{enumerate}
\end{definition}

\noindent
따라서 연속사상 $f : X \to Y$가 몫위상 사상일 필요충분조건은 $f$가
전사이고, 임의의 $T \subset Y$에 대하여 $T$가 열리거나 닫힐
필요충분조건이 $f^{-1}(T)$도 그러한 것이다. 다시 말해 $Y$에는 전사사상
$X \to Y$에 대한 몫위상이 주어진다.

\begin{lemma}
\label{topology-lemma-open-morphism-quotient-topology}
$f : X \to Y$를 위상공간들의 전사이고 열린 연속사상이라 하자.
$T \subset Y$를 부분집합이라 하자. 그러면
\begin{enumerate}
\item $f^{-1}(\overline{T}) = \overline{f^{-1}(T)}$.
\item $T \subset Y$가 닫혀 있을 필요충분조건은 $f^{-1}(T)$가 닫혀 있는
것이다.
\item $T \subset Y$가 열려 있을 필요충분조건은 $f^{-1}(T)$가 열려 있는
것이다.
\item $T \subset Y$가 국소닫혀 있을 필요충분조건은 $f^{-1}(T)$가
국소닫혀 있는 것이다.
\end{enumerate}
특히 $f$는 몫위상 사상이다.
\end{lemma}

\begin{proof}
$\overline{f^{-1}(T)} \subset f^{-1}(\overline{T})$임은 분명하다.
$x \in X$이고 $x \not \in \overline{f^{-1}(T)}$라 하자. 그러면
$U \cap f^{-1}(T) = \emptyset$인 열린 근방 $x \in U \subset X$가
존재한다. $f$가 열린사상이므로 $f(U)$는 $T$와 만나지 않는 $f(x)$의
열린 근방이다. 따라서 $x \not \in f^{-1}(\overline{T})$이다. 이로써
(1)을 보였다. (2)는 (1)의 쉬운 귀결이다. (3)은 $f$가 열린 전사라는
사실에서 바로 따른다. (4)를 보이자. $f^{-1}(T)$가 국소닫혀 있으면
$f^{-1}(T) \subset \overline{f^{-1}(T)} = f^{-1}(\overline{T})$는 열려 있다.
그러므로 사상 $f^{-1}(\overline{T}) \to \overline{T}$에 (3)을 적용하면
$T$가 $\overline{T}$에서 열려 있음을 알 수 있다. 즉 $T$는 국소닫혀 있다.
\end{proof}

\begin{lemma}
\label{topology-lemma-closed-morphism-quotient-topology}
$f : X \to Y$를 위상공간들의 전사이고 닫힌 연속사상이라 하자.
$T \subset Y$를 부분집합이라 하자. 그러면
\begin{enumerate}
\item $\overline{T} = f(\overline{f^{-1}(T)})$.
\item $T \subset Y$가 닫혀 있을 필요충분조건은 $f^{-1}(T)$가 닫혀 있는
것이다.
\item $T \subset Y$가 열려 있을 필요충분조건은 $f^{-1}(T)$가 열려 있는
것이다.
\item $T \subset Y$가 국소닫혀 있을 필요충분조건은
$f^{-1}(T)$가 국소닫혀 있는 것이다.
\end{enumerate}
특히 $f$는 몫위상 사상이다.
\end{lemma}

\begin{proof}
$\overline{f^{-1}(T)} \subset f^{-1}(\overline{T})$임은 분명하다.
그러면 $T \subset f(\overline{f^{-1}(T)}) \subset \overline{T}$는 닫힌
부분집합이므로 (1)을 얻는다. (2)는 $f$가 닫힌 전사라는 사실에서 바로
따른다. (3)은 $T$의 여집합에 (2)를 적용하면 얻는다. (4)를 보이자.
$f^{-1}(T)$가 국소닫혀 있으면
$f^{-1}(T) \subset \overline{f^{-1}(T)}$는 열려 있다. (1)에 의해 사상
$\overline{f^{-1}(T)} \to \overline{T}$가 전사이므로, $f$가 유도하는
사상 $\overline{f^{-1}(T)} \to \overline{T}$에 (3)을 적용하여 $T$가
$\overline{T}$에서 열려 있음을 얻는다. 즉 $T$는 국소닫혀 있다.
\end{proof}






\section{연결성분}
\label{topology-section-connected-components}

\begin{definition}
\label{topology-definition-connected-components}
$X$를 위상공간이라 하자.
\begin{enumerate}
\item $X$가 공집합이 아니고, $T_i \subset X$가 열린 동시에 닫혀 있으며
$X = T_1 \amalg T_2$일 때마다 $T_1 = \emptyset$이거나
$T_2 = \emptyset$이면 $X$를 {\it 연결되었다}고 한다.
\item $T \subset X$가 연결되어 있고 $T$가 $X$의 극대 연결 부분집합이면 이를 $X$의
{\it 연결성분}이라고 한다.
\end{enumerate}
\end{definition}

\noindent
공집합은 연결되어 있지 않다.

\begin{lemma}
\label{topology-lemma-image-connected-space}
$f : X \to Y$를 위상공간들의 연속사상이라 하자.
$E \subset X$가 연결 부분집합이면 $f(E) \subset Y$도 연결되어 있다.
\end{lemma}

\begin{proof}
$A \subset f(E)$를 $f(E)$의 열린 동시에 닫힌 부분집합이라 하자.
$f$가 연속이므로 $f^{-1}(A)$는 $E$의 열린 동시에 닫힌 부분집합이다.
$E$가 연결되어 있으므로 $f^{-1}(A) = \emptyset$이거나
$f^{-1}(A) = E$이다. 한편 $A\subset f(E)$이므로
$A = f(f^{-1}(A))$이다. 실제로 $x\in f(f^{-1}(A))$이면
$f(y) = x$인 $y\in f^{-1}(A)$가 존재하고, $y\in f^{-1}(A)$이므로
$f(y) \in A$, 즉 $x\in A$이다. 반대로 $x\in A$이면
$A\subset f(E)$이므로 $f(y) = x$인 $y\in E$가 존재한다. 따라서
$y\in f^{-1}(A)$이고 $x\in f(f^{-1}(A))$이다. 그러므로
$A = \emptyset$이거나 $A = f(E)$이며, 이는 $f(E)$가 연결되어 있음을 보인다.
\end{proof}

\begin{lemma}
\label{topology-lemma-connected-components}
$X$를 위상공간이라 하자.
\begin{enumerate}
\item $T \subset X$가 연결되어 있으면 그 폐포도 연결되어 있다.
\item $X$의 모든 연결성분은 닫혀 있다(반드시 열려 있는 것은 아니다).
\item $X$의 모든 연결 부분집합은 $X$의 유일한 연결성분에 포함된다.
\item $X$의 모든 점은 유일한 연결성분에 포함된다. 다시 말해 $X$는
연결성분들의 서로소 합이다.
\end{enumerate}
\end{lemma}

\begin{proof}
연결 부분집합 $T$의 폐포를 $\overline{T}$라 하자.
$T_i \subset \overline{T}$가 열린 동시에 닫혀 있고
$\overline{T} = T_1 \amalg T_2$라고 하자. 그러면
$T = (T\cap T_1) \amalg (T \cap T_2)$이다. 따라서 $T$는 둘 중 하나와
같으며, 이를 $T = T_1 \cap T$라 하자. 그러면
$\overline{T} \subset T_1$이다. 이로써 (1)과 (2)를 얻는다.

\medskip\noindent
$A$를 $X$의 연결 부분집합들로 이루어진 공집합이 아닌 집합이라 하고,
$\Omega = \bigcap_{T \in A} T$가 공집합이 아니라고 하자.
$E = \bigcup_{T \in A} T$가 연결되어 있음을 보이겠다. 실제로 $E$는
$\Omega$를 포함하므로 공집합이 아니다. $E_i$가 $E$에서 닫혀 있고
$E = E_1 \amalg E_2$라고 하자. 번호를 바꾸어서 $E_1$이 $\Omega$와
만난다고 해도 된다. 그러면 각 $T \in A$는 $E_1$과 만나며, $T$가
연결되어 있으므로 $E_1$에 포함되어야 한다. 따라서 $E \subset E_1$이며
주장이 증명된다.

\medskip\noindent
$W \subset X$를 공집합이 아닌 연결 부분집합이라 하자. $W$를 포함하는
$X$의 모든 연결 부분집합들의 집합에 앞 문단의 결과를 적용하면 $E$가
$X$의 연결성분임을 알 수 있다. 이로부터 (3)의 존재성과 유일성이 따른다.

\medskip\noindent
$x \in X$라 하자. 앞 문단에서 $W = \{x\}$로 두면 $x$가 $X$의 유일한
연결성분에 포함됨을 알 수 있다. 서로 다른 두 연결성분은 서로소여야 한다
(둘째 문단의 결과에 의한다).

\medskip\noindent
연결성분이 열려 있지 않은 예로 곱위상을 준 무한 곱
$\prod_{n \in \mathbf{N}} \{0, 1\}$을 택하면 된다. 그 연결성분들은
한원소집합이며 열려 있지 않다.
\end{proof}

\begin{remark}
\label{topology-remark-quasi-components}
\begin{reference}
\cite[Example 6.1.24]{Engelking}
\end{reference}
$X$를 위상공간, $x \in X$를 점이라 하자. $Z \subset X$를 $x$를 지나는
$X$의 연결성분이라 하자. $x$를 포함하는 $X$의 모든 열린 동시에 닫힌
부분집합들의 교집합을 $E$라 하자. $Z \subset E$임은 분명하다.
일반적으로 $Z \not = E$이다. 예를 들어
$X = \{x, y, z_1, z_2, \ldots\}$에 모든 $n$에 대한
$\{z_n\}$, $\{x, z_n, z_{n + 1}, \ldots\}$,
$\{y, z_n, z_{n + 1}, \ldots\}$를 열린집합의 기저로 갖는 위상을 주자.
그러면 $Z = \{x\}$이고 $E = \{x, y\}$이다. 자세한 내용은 생략한다.
\end{remark}

\begin{lemma}
\label{topology-lemma-connected-fibres-quotient-topology-connected-components}
$f : X \to Y$를 위상공간들의 연속사상이라 하자. 다음을 가정한다.
\begin{enumerate}
\item $f$의 모든 섬유는 연결되어 있다.
\item 집합 $T \subset Y$가 닫혀 있을 필요충분조건은 $f^{-1}(T)$가
닫혀 있는 것이다.
\end{enumerate}
그러면 $f$는 $X$와 $Y$의 연결성분들의 집합 사이의 전단사를 유도한다.
\end{lemma}

\begin{proof}
$T \subset Y$를 연결성분이라 하자. Lemma
\ref{topology-lemma-connected-components}에 의해 $T$는 닫혀 있다.
$f^{-1}(T)$가 연결되어 있음을 보이면 보조정리가 따른다. 실제로 Lemma
\ref{topology-lemma-image-connected-space}에 의해 $X$의 모든 연결 부분집합은 $Y$의
한 연결성분 안으로 사상된다. $Z_1$, $Z_2$가 닫혀 있고
$f^{-1}(T) = Z_1 \amalg Z_2$라고 하자. 임의의 $t \in T$에 대하여
$f^{-1}(\{t\}) = Z_1 \cap f^{-1}(\{t\}) \amalg Z_2 \cap f^{-1}(\{t\})$이다.
(1)에 의해 $f^{-1}(\{t\})$가 연결되어 있으므로
$f^{-1}(\{t\}) \subset Z_1$이거나 $f^{-1}(\{t\}) \subset Z_2$이다.
다시 말해 $f^{-1}(T_i) = Z_i$인 분해 $T = T_1 \amalg T_2$가 있다.
(2)에 의해 $T_i$는 $Y$에서 닫혀 있다. 따라서 원하는 대로
$T_1 = \emptyset$이거나 $T_2 = \emptyset$이다.
\end{proof}

\begin{lemma}
\label{topology-lemma-connected-fibres-connected-components}
$f : X \to Y$를 위상공간들의 연속사상이라 하자. 다음을 가정한다.
(a) $f$는 열린사상이다.
(b) $f$의 모든 섬유는 연결되어 있다.
그러면 $f$는 $X$와 $Y$의 연결성분들의 집합 사이의 전단사를 유도한다.
\end{lemma}

\begin{proof}
이는 Lemma
\ref{topology-lemma-connected-fibres-quotient-topology-connected-components}의
특수한 경우이다.
\end{proof}

\begin{lemma}
\label{topology-lemma-finite-fibre-connected-components}
$f : X \to Y$를 공집합이 아닌 위상공간들의 연속사상이라 하자. 다음을
가정한다. (a) $Y$는 연결되어 있다. (b) $f$는 열린 동시에 닫힌사상이다.
(c) 섬유 $f^{-1}(y)$가 유한집합인 점 $y\in Y$가 존재한다. 그러면 $X$는
많아야 $|f^{-1}(y)|$개의 연결성분을 갖는다. 따라서 $X$의 모든 연결성분
$T$는 열린 동시에 닫혀 있고, $f(T)$는 $Y$의 공집합이 아닌 열린 동시에
닫힌 부분집합이므로 $Y$와 같다.
\end{lemma}

\begin{proof}
어떤 $N \in \mathbf{N}$에 대하여 위상공간 $X$가 적어도 $N$개의
연결성분을 갖는다고 하자. 귀납법으로 $X$를 $N$개의 공집합이 아닌 열린 동시에
닫힌 부분집합들 $X_1, \ldots , X_N$의 서로소 합으로 나타내는 분해
$X = X_1 \amalg \ldots \amalg X_N$을 얻는다. $f$가 열린 동시에 닫힌
사상이므로 각 $f(X_i)$는 $Y$의 공집합이 아닌 열린 동시에 닫힌 부분집합이고,
따라서 $Y$와 같다. 특히 각 $1 \leq i \leq N$에 대하여 교집합
$X_i \cap f^{-1}(y)$는 공집합이 아니다. 그러므로 $f^{-1}(y)$는 적어도
$N$개의 원소를 갖는다.
\end{proof}

\begin{definition}
\label{topology-definition-totally-disconnected}
모든 연결성분이 한원소집합인 위상공간을 {\it 완전불연결}이라고 한다.
\end{definition}

\noindent
이산공간은 완전불연결이다. 완전불연결 공간이 반드시 이산공간인 것은
아니다. 예를 들어 $\mathbf{Q} \subset \mathbf{R}$는 완전불연결이지만
이산공간은 아니다.

\begin{lemma}
\label{topology-lemma-space-connected-components}
$X$를 위상공간이라 하자. $\pi_0(X)$를 $X$의 연결성분들의 집합이라 하자.
$X \to \pi_0(X)$를 $x \in X$를 $x$를 지나는 $X$의 연결성분으로 보내는
사상이라 하자. $\pi_0(X)$에 몫위상을 주자. 그러면 $\pi_0(X)$는
완전불연결 공간이고, $X$에서 완전불연결 공간 $Y$로 가는 임의의 연속사상
$X \to Y$는 $\pi_0(X)$를 통해 분해된다.
\end{lemma}

\begin{proof}
Lemma \ref{topology-lemma-connected-fibres-quotient-topology-connected-components}에
의해 $\pi_0(X)$의 연결성분들은 한원소집합이다. 둘째 명제의 증명은
생략한다.
\end{proof}

\begin{definition}
\label{topology-definition-locally-connected}
위상공간 $X$의 모든 점 $x \in X$가 연결 근방들의 기본계를 가지면
$X$를 {\it 국소연결}이라고 한다.
\end{definition}

\begin{lemma}
\label{topology-lemma-locally-connected}
$X$를 위상공간이라 하자. $X$가 국소연결이면 다음이 성립한다.
\begin{enumerate}
\item $X$의 모든 열린 부분집합은 국소연결이다.
\item $X$의 연결성분들은 열려 있다.
\end{enumerate}
따라서 $X$의 열린 부분집합들의 연결성분들도 열려 있다. 특히 모든 점은
열린 연결 근방들의 기본계를 갖는다.
\end{lemma}

\begin{proof}
모든 $x\in X$에 대하여 $x$의 연결 근방들의 기본계를
$\mathcal{N}(x)$라 쓰고, $U\subset X$를 $X$의 열린 부분집합이라 하자.
모든 $x\in U$에 대하여 $U$는 $x$의 근방이므로 집합
$\{V \in \mathcal{N}(x) | V\subset U\}$는 공집합이 아니며 $U$에서
$x$의 연결 근방들의 기본계이다. 따라서 $U$는 국소연결이고 (1)이
증명된다.

\medskip\noindent
$x \in \mathcal{C} \subset X$이고 $\mathcal{C}$가 $x$의 연결성분이라고
하자. $X$가 국소연결이므로 $x$의 연결 근방 $\mathcal{N}$이 존재한다.
그러므로 연결성분의 정의에 의해 $\mathcal{N} \subset \mathcal{C}$이고,
$\mathcal{C}$는 $x$의 근방이다. 따라서 $\mathcal{C}$는 그 각 점의
근방이다. 다시 말해 $\mathcal{C}$는 열려 있고 (2)가 증명된다.
\end{proof}

\section{기약성분}
\label{topology-section-irreducible-components}

\begin{definition}
\label{topology-definition-irreducible-components}
$X$를 위상공간이라 하자.
\begin{enumerate}
\item $X$가 공집합이 아니고, $Z_i$가 닫혀 있으며
$X = Z_1 \cup Z_2$일 때마다 $X = Z_1$이거나 $X = Z_2$이면
$X$를 {\it 기약}이라고 한다.
\item $Z \subset X$가 $X$의 극대 기약 부분집합이면 $Z$를 $X$의
{\it 기약성분}이라고 한다.
\end{enumerate}
\end{definition}

\noindent
기약공간은 분명히 연결되어 있다.

\begin{lemma}
\label{topology-lemma-image-irreducible-space}
$f : X \to Y$를 위상공간들의 연속사상이라 하자.
$E \subset X$가 기약 부분집합이면 $f(E) \subset Y$도 기약이다.
\end{lemma}

\begin{proof}
$E = X$이고(즉 $X$가 기약이고) $f(E) = Y$라고(즉 $f$가 전사라고)
가정해도 됨은 분명하다. 먼저 $X \not = \emptyset$이므로
$Y \not = \emptyset$이다. 다음으로 $Y_1$, $Y_2$가 닫혀 있고
$Y = Y_1 \cup Y_2$라고 하자. 이때 $X_i = f^{-1}(Y_i)$는 $X$에서 닫혀 있고
$X = X_1 \cup X_2$이다. $X$에 대한 가정에 의해
$X = X_1$이거나 $X = X_2$이다. $f$가 전사이므로
$Y = Y_1$이거나 $Y = Y_2$이다.
\end{proof}

\begin{lemma}
\label{topology-lemma-irreducible}
$X$를 위상공간이라 하자.
\begin{enumerate}
\item $T \subset X$가 기약이면 $X$에서의 그 폐포도 기약이다.
\item $X$의 모든 기약성분은 닫혀 있다.
\item $X$의 모든 기약 부분집합은 $X$의 한 기약성분에 포함된다.
\item $X$의 모든 점은 $X$의 어떤 기약성분에 포함된다. 다시 말해
$X$는 기약성분들의 합집합이다.
\end{enumerate}
\end{lemma}

\begin{proof}
기약 부분집합 $T$의 폐포를 $\overline{T}$라 하자.
$Z_i \subset \overline{T}$가 닫혀 있고
$\overline{T} = Z_1 \cup Z_2$이면
$T = (T\cap Z_1) \cup (T \cap Z_2)$이다. 따라서 $T$는 둘 중 하나와
같으며, 이를 $T = Z_1 \cap T$라 하자. 그러면 분명히
$\overline{T} \subset Z_1$이다. 이로써 (1)이 증명된다. (2)는 (1)과
기약성분의 정의에서 바로 따른다.

\medskip\noindent
$T \subset X$를 기약이라 하자. $T \subset T' \subset X$인 기약
부분집합들의 집합을 $A$라 하자. $T \in A$이므로 $A$는 공집합이 아니다.
포함관계에서 오는 부분순서 $T' \leq T'' \Leftrightarrow T' \subset T''$를
$A$ 위에 준다. 극대 전순서 부분집합 $A' \subset A$를 택하고
$E = \bigcup_{T' \in A'} T'$라 하자. $E$가 기약임을 보이겠다.
이제 $E =  Z_1 \cup Z_2$가 $E$의 두 닫힌 부분집합의 합집합이라고
하자. 각 $T' \in A'$에 대하여 $T'$의 기약성에 의해
$T' \subset Z_1$이거나 $T' \subset Z_2$이다. 어떤 $T'_0 \in A'$에
대하여 $T'_0 \not\subset Z_1$이라고 하자(그렇지 않으면 이미 끝난다).
$A'$는 전순서 집합이므로 모든 $T' \in A'$에 대하여 곧바로
$T' \subset Z_2$임을 알 수 있다. 따라서 $E = Z_2$이다. 이로써 (3)을
증명했다. 한원소공간은 기약이므로 (4)는 (3)의 바로 귀결이다.
\end{proof}

\begin{lemma}
\label{topology-lemma-pick-irreducible-components}
$X$를 위상공간이라 하고
$X = \bigcup_{i = 1, \ldots, n} X_i$라고 하자. 각 $X_i$가 $X$의 기약
닫힌 부분집합이고 어느 $X_i$도 나머지 원소들의 합집합에 포함되지 않는다고
가정하자. 그러면 각 $X_i$는 $X$의 기약성분이고, $X$의 각 기약성분은
$X_i$ 가운데 하나이다.
\end{lemma}

\begin{proof}
$Y$를 $X$의 기약성분이라 하자.
$Y = \bigcup_{i = 1, \ldots, n} (Y \cap X_i)$로 쓰자. 각 $X_i$가 $X$에서
닫혀 있으므로 각 $Y \cap X_i$는 $Y$에서 닫혀 있다. $Y$의 기약성에 의해
$Y$는 $Y \cap X_i$ 중 하나와 같다. 즉 $Y \subset X_i$이다.
기약성분의 극대성에 의해 $Y = X_i$를 얻는다.

\medskip\noindent
반대로 $X_i$ 하나를 택하자. Lemma \ref{topology-lemma-irreducible}에 의해 이는
$X$의 어떤 기약성분 $Y$에 포함된다. 앞 문단에 의해 어떤 $j$에 대하여
$Y = X_j$이다. 따라서 $X_i \subset X_j$이고, 원래의 합집합에는 불필요한
원소가 없으므로 $X_i = X_j$이다. 이는 기약성분이다.
\end{proof}

\begin{lemma}
\label{topology-lemma-surjective-continuous-irreducible-components}
$f : X \to Y$를 위상공간들의 전사 연속사상이라 하자.
$X$가 유한 개, 이를 $n$개라 하자, 의 기약성분을 가지면 $Y$는
$\leq n$개의 기약성분을 갖는다.
\end{lemma}

\begin{proof}
$X_1, \ldots, X_n$이 $X$의 기약성분들이라고 하자. Lemmas
\ref{topology-lemma-image-irreducible-space}와 \ref{topology-lemma-irreducible}에 의해
$f(X_i)$의 폐포 $Y_i \subset Y$는 기약이다. $f$가 전사이므로 $Y$는
$Y_i$들의 합집합이다. $Y = \bigcup_{i \in I} Y_i$가 되도록 하는 극소
부분집합 $I \subset \{1, \ldots, n\}$을 택할 수 있다. 그러면 Lemma
\ref{topology-lemma-pick-irreducible-components}를 적용하여 $i \in I$에 대한
$Y_i$들이 $Y$의 기약성분들임을 알 수 있다.
\end{proof}

\noindent
한원소집합은 기약이다. 따라서 $x \in X$가 한 점이면
폐포 $\overline{\{x\}}$는 $X$의 기약 닫힌 부분집합이다.

\begin{definition}
\label{topology-definition-generic-point}
$X$를 위상공간이라 하자.
\begin{enumerate}
\item $Z \subset X$를 기약 닫힌 부분집합이라 하자. $Z$의
{\it 일반점}은 $Z = \overline{\{\xi\}}$를 만족하는 점 $\xi \in Z$이다.
\item 모든 $x, x' \in X$와 $x \not = x'$에 대하여 두 점 중 정확히 하나만
포함하는 $X$의 닫힌 부분집합이 존재하면 $X$를 {\it 콜모고로프}라고 한다.
\item 모든 기약 닫힌 부분집합이 일반점을 가지면 $X$를
{\it 준소버}라고 한다.
\item 모든 기약 닫힌 부분집합이 유일한 일반점을 가지면 $X$를
{\it 소버}라고 한다.
\end{enumerate}
\end{definition}

\noindent
위상공간 $X$가 각각 콜모고로프, 준소버, 소버일 필요충분조건은
$X$에서 $X$의 기약 닫힌 부분집합들의 집합으로 가는 사상
$x\mapsto\overline{\{x\}}$가 각각 단사, 전사, 전단사인 것이다.
따라서 위상공간이 소버일 필요충분조건은 준소버이면서 콜모고로프인 것이다.

\begin{lemma}
\label{topology-lemma-sober-subspace}
$X$를 위상공간이라 하고 $Y\subset X$라 하자.
\begin{enumerate}
\item $X$가 콜모고로프이면 $Y$도 콜모고로프이다.
\item $Y$가 $X$에서 국소닫힌이라고 하자. $X$가 준소버이면 $Y$도
준소버이다.
\item $Y$가 $X$에서 국소닫힌이라고 하자. $X$가 소버이면 $Y$도
소버이다.
\end{enumerate}
\end{lemma}

\begin{proof}
(1)의 증명. $X$가 콜모고로프라고 하자. $x,y\in Y$이고 $x\neq y$라 하자.
그러면
$\overline{\overline{\{x\}}\cap Y}=\overline{\{x\}}\neq\overline{\{y\}}=
\overline{\overline{\{y\}}\cap Y}$이다. 따라서
$\overline{\{x\}}\cap Y\neq\overline{\{y\}}\cap Y$이다. 이는 $Y$가
콜모고로프임을 보인다.

\medskip\noindent
(2)의 증명. $X$가 준소버라고 하자. $Y$가 닫힌 경우와 $Y$가 열린 경우만
생각하면 충분하다. 먼저 $Y$가 닫혀 있다고 하자. $Z$를 $Y$의 기약 닫힌
부분집합이라 하자. 그러면 $Z$는 $X$의 기약 닫힌 부분집합이다. 따라서
$\overline{\{x\}} = Z$인 $x \in Z$가 존재한다. 그러면
$\overline{\{x\}} \cap Y = Z$이다. 이는 $Y$가 준소버임을 보인다.
다음으로 $Y$가 열려 있다고 하자. $Z$를 $Y$의 기약 닫힌 부분집합이라
하자. 그러면 $\overline{Z}$는 $X$의 기약 닫힌 부분집합이다. 따라서
$\overline{\{x\}}=\overline{Z}$인 $x \in \overline{Z}$가 존재한다.
$x\notin Y$이면
$Z=Z\cap Y\subset\overline{Z}\cap Y=\overline{\{x\}}\cap Y=\emptyset$이라는
모순을 얻는다. 따라서 $x\in Y$이다. 그러면
$Z=\overline{Z}\cap Y=\overline{\{x\}}\cap Y$이다. 이는 $Y$가
준소버임을 보인다.

\medskip\noindent
(3)의 증명. (1)과 (2)에서 바로 따른다.
\end{proof}

\begin{lemma}
\label{topology-lemma-sober-local}
$X$를 위상공간이라 하고 $(X_i)_{i\in I}$를 $X$의 덮개라 하자.
\begin{enumerate}
\item 모든 $i\in I$에 대하여 $X_i$가 $X$에서 국소닫힌이라고 하자.
그러면 $X$가 콜모고로프일 필요충분조건은 모든 $i\in I$에 대하여
$X_i$가 콜모고로프인 것이다.
\item 모든 $i\in I$에 대하여 $X_i$가 $X$에서 열려 있다고 하자.
그러면 $X$가 준소버일 필요충분조건은 모든 $i\in I$에 대하여
$X_i$가 준소버인 것이다.
\item 모든 $i\in I$에 대하여 $X_i$가 $X$에서 열려 있다고 하자.
그러면 $X$가 소버일 필요충분조건은 모든 $i\in I$에 대하여
$X_i$가 소버인 것이다.
\end{enumerate}
\end{lemma}

\begin{proof}
(1)의 증명. $X$가 콜모고로프이면 Lemma \ref{topology-lemma-sober-subspace}에 의해
모든 $i\in I$에 대하여 $X_i$도 콜모고로프이다. 모든 $i\in I$에 대하여
$X_i$가 콜모고로프라고 하자. $x,y\in X$이고
$\overline{\{x\}}=\overline{\{y\}}$라 하자. $x\in X_i$인 $i\in I$가
존재한다. $X_i$가 $U$의 닫힌 부분집합이 되는 열린 부분집합
$U\subset X$가 존재한다. $y\notin U$이면
$x\in\overline{\{x\}}\cap U=\overline{\{y\}}\cap U=\emptyset$이라는 모순을
얻는다. 따라서 $y\in U$이다. 그러므로
$y\in\overline{\{y\}}\cap U=\overline{\{x\}}\cap U\subset X_i$이다.
이는 $y\in X_i$임을 보인다. 따라서
$\overline{\{x\}}\cap X_i=\overline{\{y\}}\cap X_i$이다. $X_i$가
콜모고로프이므로 $x=y$이다. 이는 $X$가 콜모고로프임을 보인다.

\medskip\noindent
(2)의 증명. $X$가 준소버이면 Lemma \ref{topology-lemma-sober-subspace}에 의해
모든 $i\in I$에 대하여 $X_i$도 준소버이다. 모든 $i\in I$에 대하여
$X_i$가 준소버라고 하자. $Y$를 $X$의 기약 닫힌 부분집합이라 하자.
$Y\neq\emptyset$이므로 $X_i\cap Y\neq\emptyset$인 $i\in I$가 존재한다.
$X_i$가 $X$에서 열려 있으므로 $X_i\cap Y$는 $Y$에서 공집합이 아닌
열린 부분집합이며, 따라서 기약이고 $Y$에서 조밀하다. 그러므로
$X_i\cap Y$는 $X_i$의 기약 닫힌 부분집합이다. $X_i$가 준소버이므로
$X_i\cap Y=\overline{\{x\}}\cap X_i\subset\overline{\{x\}}$인
$x\in X_i\cap Y$가 존재한다. $X_i\cap Y$가 $Y$에서 조밀하고 $Y$가
$X$에서 닫혀 있으므로
$Y=\overline{X_i\cap Y}\cap Y\subset\overline{X_i\cap Y}\subset
\overline{\{x\}}\subset Y$이다. 따라서 $Y=\overline{\{x\}}$이다.
이는 $X$가 준소버임을 보인다.

\medskip\noindent
(3)의 증명. (1)과 (2)에서 바로 따른다.
\end{proof}

\begin{example}
\label{topology-example-quasi-sober-not-kolmogorov}
$X$를 기수가 적어도 $2$인 비이산공간이라 하자. 그러면 $X$는
준소버이지만 콜모고로프가 아니다. 더욱이 그 한원소집합들의 모임은
$X$를 이산공간, 따라서 콜모고로프 공간들로 덮는다.
\end{example}

\begin{example}
\label{topology-example-kolmogorov-not-quasi-sober}
$Y$를 무한집합이라 하고, 닫힌집합들이 $Y$와 $Y$의 유한 부분집합들인
위상을 주자. 그러면 $Y$는 콜모고로프이지만 준소버가 아니다. 그러나 그
한원소집합들의 모임은 이산공간, 따라서 소버 공간들로 이루어진 덮개이다.
\end{example}

\begin{example}
\label{topology-example-not-kolmogorov-not-quasi-sober}
$X$와 $Y$를 각각 Example \ref{topology-example-quasi-sober-not-kolmogorov}와
Example \ref{topology-example-kolmogorov-not-quasi-sober}에서와 같이 택하자.
그러면 $X\amalg Y$는 콜모고로프도 준소버도 아니다.
\end{example}

\begin{example}
\label{topology-example-sober-subspace}
$Z$를 무한집합이라 하고 $z\in Z$라 하자. 닫힌집합들이 $Z$와
$Z\setminus\{z\}$의 유한 부분집합들인 위상을 $Z$에 주자. 그러면 $Z$는
소버이지만 그 부분공간 $Z\setminus\{z\}$는 준소버가 아니다.
\end{example}

\begin{example}
\label{topology-example-Hausdorff}
위상공간 $X$가 하우스도르프일 필요충분조건은 서로 다른 임의의 두 점
$x, y \in X$에 대하여 $x \in U$, $y \in V$를 만족하는 서로소 열린
부분집합들 $U, V \subset X$가 존재하는 것임을 상기하자. 이 경우
$X$가 기약일 필요충분조건은 $X$가 한원소집합인 것이다. 마찬가지로
$X$의 임의의 부분집합이 기약일 필요충분조건은 그것이 한원소집합인 것이다.
따라서 하우스도르프 공간은 소버이다.
\end{example}

\begin{lemma}
\label{topology-lemma-irreducible-on-top}
$f : X \to Y$를 위상공간들의 연속사상이라 하자. 다음을 가정한다.
(a) $Y$는 기약이다.
(b) $f$는 열린사상이다.
(c) $f^{-1}(y)$가 기약인 점들 $y \in Y$의 조밀한 모임이 존재한다.
그러면 $X$는 기약이다.
\end{lemma}

\begin{proof}
$Z_i$가 닫혀 있고 $X = Z_1 \cup Z_2$라고 하자.
열린집합들 $U_1 = Z_1 \setminus Z_2 = X \setminus Z_2$와
$U_2 = Z_2 \setminus Z_1 = X \setminus Z_1$을 생각하자. 모순을 얻기
위해 $U_1$과 $U_2$가 모두 공집합이 아니라고 가정하자. (b)에 의해
$f(U_i)$는 열려 있다. (a)에 의해 $Y$는 기약이고, 따라서
$f(U_1) \cap f(U_2) \not = \emptyset$이다. (c)에 의해 이 교집합의 한 점에
대응하고 섬유 $X_y = f^{-1}(y)$가 기약인 점 $y$가 존재한다. 그러면
$X_y \cap U_1$과 $X_y \cap U_2$는 $X_y$의 서로소인 공집합이 아닌 열린
부분집합들이므로 모순이다.
\end{proof}

\begin{lemma}
\label{topology-lemma-irreducible-fibres-irreducible-components}
$f : X \to Y$를 위상공간들의 연속사상이라 하자. 다음을 가정한다.
(a) $f$는 열린사상이다.
(b) 모든 $y \in Y$에 대하여 섬유 $f^{-1}(y)$는 기약이다.
그러면 $f$는 기약성분들 사이의 전단사를 유도한다.
\end{lemma}

\begin{proof}
가정 (b)는 $f$가 전사임을 함의함을 지적해 둔다(Definition
\ref{topology-definition-irreducible-components} 참조).
$T \subset Y$를 기약성분이라 하자. Lemma \ref{topology-lemma-irreducible}에서
보았듯이 $T$는 닫혀 있다. $X$의 모든 기약 부분집합은 Lemma
\ref{topology-lemma-image-irreducible-space}에 의해 $Y$의 한 기약성분으로
사상되므로, $f^{-1}(T)$가 기약임을 보이면 보조정리가 따른다.
$f^{-1}(T) \to T$는 Lemma \ref{topology-lemma-irreducible-on-top}의 가정들을
만족한다. 따라서 결론을 얻는다.
\end{proof}

\noindent
다음 보조정리의 구성을 때때로 ‘소버화’라고 한다.

\begin{lemma}
\label{topology-lemma-make-sober}
\begin{reference}
\cite[Page 68 ff]{EGA1-second}
\end{reference}
$X$를 위상공간이라 하자. $X$에서 소버 위상공간 $X'$으로 가는 표준
연속사상
$$
c : X \longrightarrow X'
$$
가 존재하며, 이는 $X$에서 소버 위상공간으로 가는 연속사상들 가운데
보편적이다. 더욱이 대응 $U' \mapsto c^{-1}(U')$은 $X'$의 열린집합들과
$X$의 열린집합들 사이의 전단사이며 유한 교집합과 임의의 합집합과
가환한다. 상 $c(X)$는 콜모고로프 위상공간이고 사상
$c : X \to c(X)$는 $X$에서 콜모고로프 공간으로 가는 사상들에 대해
보편적이다.
\end{lemma}

\begin{proof}
$X'$을 $X$의 기약 닫힌 부분집합들의 집합이라 하고
$$
c : X \to X', \quad x \mapsto \overline{\{x\}}.
$$
라 하자. 열린집합 $U \subset X$에 대하여 $U' \subset X'$을 $U$와 만나는
$X$의 기약 닫힌 부분집합들의 집합이라 하자. 그러면
$c^{-1}(U') = U$이다. 특히 $U_1 \not = U_2$가 $X$의 열린집합들이면
$U'_1 \not = U_2'$이다. 따라서 $c$는 $U'$ 꼴인 $X'$의 부분집합들과
$X$의 열린집합들 사이의 전단사를 유도한다.

\medskip\noindent
$U_1, U_2$를 $X$에서 열린집합이라 하자. $Z \in U'_1$이고
$Z \in U'_2$라고 하자. 그러면 $Z \cap U_1$과 $Z \cap U_2$는 기약공간
$Z$의 공집합이 아닌 열린 부분집합들이므로 $Z \cap U_1 \cap U_2$는
공집합이 아니다. 따라서 $(U_1 \cap U_2)' = U'_1 \cap U'_2$이다.
규칙 $U \mapsto U'$은 임의의 합집합과도 양립한다(세부사항은 생략한다).
따라서 $U'$들의 모임은 $X'$ 위의 위상을 이루고 보조정리에서 말한
전단사를 얻음이 분명하다.

\medskip\noindent
이제 $X'$이 소버임을 보이자. $T \subset X'$을 기약 닫힌 부분집합이라
하자. $X' \setminus T = U'$가 되는 열린집합 $U \subset X$를 택하자.
그러면 보조정리의 전단사가 갖는 성질들에 의해 $Z = X \setminus U$는
기약이다. $Z \in T$가 유일한 일반점임을 주장한다. 실제로 $Z$를 포함하지
않는 $V' \subset X'$ 꼴의 임의의 열린집합은 $Z$와 만나지 않는, 즉
$U$에 포함되는 열린집합 $V \subset X$에서 와야 한다.

\medskip\noindent
마지막으로 보편성질을 확인하자. $f : X \to Y$를 소버 위상공간으로 가는
연속사상이라 하자. $f' : X' \to Y$를 기약 닫힌 부분집합
$Z \subset X$를 $\overline{f(Z)}$의 유일한 일반점으로 보내는 사상으로
정한다. 집합들의 사상으로서 $f' \circ c = f$임은 곧바로 따르고, $c$의
성질들에 의해 $f'$은 연속이다. 연속사상 $f'$의 유일성에 대한 확인은
생략한다. 콜모고로프 공간들에 관한 명제들의 증명도 생략한다.
\end{proof}

\begin{lemma}
\label{topology-lemma-remove-irreducible-connected}
$X$를 유한 개의 기약성분 $X_1, \ldots, X_n$을 갖는 연결 위상공간이라
하자. $n > 1$이면
$X' = \bigcup_{i \not = j} X_i$가 연결이 되는 $1 \leq j \leq n$이
존재한다.
\end{lemma}

\begin{proof}
이는 그래프 이론 문제이다. 꼭짓점들이 $V = \{1, \ldots, n\}$이고
$X_i \cap X_j$가 공집합이 아닐 필요충분조건으로 $i$와 $j$ 사이에 변이
있는 그래프를 $\Gamma$라 하자. $X$의 연결성은 $\Gamma$가 연결임을
뜻한다. 우리의 문제는 $\Gamma \setminus \{j\}$가 여전히 연결이 되는
$1 \leq j \leq n$을 찾는 것이다.
% P01-E0015: the frozen source quantifies graph vertices j,j' in E rather than V.
거리가 최대인 $j, j' \in E$를 택하면 $j$가 조건을 만족한다(잎을 택하라!).
세부사항은 생략한다.
\end{proof}

\section{뇌터 위상공간}
\label{topology-section-noetherian}

\begin{definition}
\label{topology-definition-noetherian}
% P01-E1159: the frozen definition uses X without binding it.
위상공간의 닫힌 부분집합들에 대한 내림사슬 조건이 $X$에서 성립하면 그
공간을 {\it 뇌터}라고 한다. 모든 점이 뇌터 근방을 가지는 위상공간을
{\it 국소뇌터}라고 한다.
\end{definition}

\begin{lemma}
\label{topology-lemma-Noetherian}
$X$를 뇌터 위상공간이라 하자.
\begin{enumerate}
\item $X$의 임의의 부분집합에 유도위상을 주면 뇌터이다.
\item 공간 $X$는 유한 개의 기약성분을 갖는다.
\item $X$의 각 기약성분은 $X$의 공집합이 아닌 열린집합을 포함한다.
\end{enumerate}
\end{lemma}

\begin{proof}
$T \subset X$를 $X$의 부분집합이라 하자. $T$의 닫힌 부분집합들의
내림사슬 $T_1 \supset T_2 \supset \ldots$를 택하자.
$Z_i \subset X$가 닫혀 있고 $T_i =  T \cap Z_i$라고 쓰자.
닫힌 부분집합들의 내림사슬
$Z_1 \supset Z_1\cap Z_2 \supset Z_1 \cap Z_2 \cap Z_3 \ldots$를 생각하자.
이는 가정에 의해 안정하므로 원래 수열 $T_i$도 안정한다. 따라서 $T$는
뇌터이다.

\medskip\noindent
유한 개의 기약성분을 갖지 않는 $X$의 닫힌 부분집합들의 집합을 $A$라
하자. 모순을 얻기 위해 $A$가 공집합이 아니라고 가정하자.
% P01-E1160: the frozen ordering formula compares bare indices although A contains sets.
집합 $A$에는 포함관계에서 오는 부분순서
$\alpha \leq \alpha' \Leftrightarrow Z_{\alpha} \subset Z_{\alpha'}$가 있다.
내림사슬 조건에 의해 $A$의 최소 원소 $Z$를 찾을 수 있다.
% P01-E1161: descending-chain condition gives a minimal, not necessarily least, element.
$Z$는 기약성분들의 유한 합집합이 아니므로 기약이 아니다. 따라서
$Z = Z' \cup Z''$로 쓸 수 있고 둘 다 진정으로 더 작은 닫힌 부분집합이다.
구성에 의해 $Z' = \bigcup Z'_i$와 $Z'' = \bigcup Z''_j$는 각각 그
기약성분들의 유한 합집합이다(Lemma \ref{topology-lemma-irreducible}).
따라서 $Z = \bigcup Z'_i \cup \bigcup Z''_j$는 기약 닫힌 부분집합들의
유한 합집합이다. 이 표현에서 불필요한 원소들을 제거하면 $Z$를 그
기약성분들로 분해한 것이 된다(Lemma
\ref{topology-lemma-pick-irreducible-components}). 이는 모순이다.

\medskip\noindent
$Z \subset X$를 $X$의 기약성분이라 하자. $Z_1, \ldots, Z_n$을 $X$의
나머지 기약성분들이라 하자.
$U = Z \setminus (Z_1\cup\ldots\cup Z_n)$을 생각하자. 그렇지 않으면
기약공간 $Z$가 다른 $Z_i$ 중 하나에 포함될 것이므로 이는 공집합이 아니다.
$X = Z \cup Z_1 \cup \ldots Z_n$이므로(Lemma
\ref{topology-lemma-irreducible} 참조)
$U = X \setminus (Z_1\cup\ldots\cup Z_n)$이기도 하며, 따라서 $X$에서
열려 있다. 그러므로 $Z$는 $X$의 공집합이 아닌 열린집합을 포함한다.
\end{proof}

\begin{lemma}
\label{topology-lemma-image-Noetherian}
$f : X \to Y$를 위상공간들의 연속사상이라 하자.
\begin{enumerate}
\item $X$가 뇌터이면 $f(X)$도 뇌터이다.
\item $X$가 국소뇌터이고 $f$가 열린사상이면 $f(X)$는 국소뇌터이다.
\end{enumerate}
\end{lemma}

\begin{proof}
(1)의 경우 $Z_1 \supset Z_2 \supset Z_3 \supset \ldots$가 $f(X)$의 닫힌
부분집합들의 내림사슬이라고 하자(늘 그렇듯 $Y$의 부분집합으로서의
유도위상을 준다). 그러면
$f^{-1}(Z_1) \supset f^{-1}(Z_2) \supset f^{-1}(Z_3) \supset \ldots$는
$X$의 닫힌 부분집합들의 내림사슬이다. 따라서 이 사슬은 안정한다.
$f(f^{-1}(Z_i)) = Z_i$이므로
$Z_1 \supset Z_2 \supset Z_3 \supset \ldots$도 안정한다고 결론 내린다.
(2)의 경우 $y \in f(X)$라 하자. $f(x) = y$인 $x \in X$를 택하자.
가정에 의해 $x$의 뇌터 근방 $E \subset X$가 존재한다. 그러면 (1)에 의해
$f(E) \subset f(X)$는 뇌터 근방이다.
\end{proof}

\begin{lemma}
\label{topology-lemma-finite-union-Noetherian}
$X$를 위상공간이라 하자. $X_i \subset X$, $i = 1, \ldots, n$을 유한한
부분집합들의 모임이라 하자. 각 $X_i$가 (유도위상에 관하여) 뇌터이면
$\bigcup_{i = 1, \ldots, n}  X_i$도 (유도위상에 관하여) 뇌터이다.
\end{lemma}

\begin{proof}
$X' = \bigcup_{i = 1, \ldots, n}  X_i$에 유도위상을 주고, 그 닫힌
부분집합들의 감소수열을 $\{F_m\}_{m \in \mathbf{N}}$이라 하자.
$X$의 닫힌 부분집합들의 감소수열 $\{G_m\}_{m \in \mathbf{N}}$ 중
모든 $m$에 대하여 $F_m = G_m \cap X'$를 만족하는 것을 찾을 수
있다(작은 세부사항은 생략한다). $X_i$가 뇌터이고
$\{G_m \cap X_i\}_{m \in \mathbf{N}}$가 $X_i$의 닫힌 부분집합들의
감소수열이므로, 모든 $m \geq m_i$에 대하여
$G_m \cap X_i = G_{m_i} \cap X_i$가 되는 $m_i \in \mathbf{N}$가
존재한다. $m_0 = \max_{i = 1, \ldots, n} m_i$라 하자. 그러면 분명히
% P01-E1164: the frozen display omits a union operator after the ellipsis.
$$
F_m = G_m \cap X' = G_m \cap (X_1 \cup \ldots \cup X_n) =
(G_m \cap X_1) \cup \ldots (G_m \cap X_n)
$$
는 $m \geq m_0$에 대하여 안정하며 증명이 끝난다.
\end{proof}

\begin{example}
\label{topology-example-locally-Noetherian-no-closed-point}
공집합이 아닌 콜모고로프 뇌터 위상공간은 모두 닫힌점을 갖는다(Lemmas
\ref{topology-lemma-quasi-compact-closed-point}와
\ref{topology-lemma-Noetherian-quasi-compact}를 결합하라).
$X = \{1, 2, 3, \ldots \}$라 하자. 열린집합들이 $\emptyset$,
$\{1, 2, \ldots, n\}$, $n \geq 1$, 그리고 $X$인 위상을 $X$ 위에
정의하자. 그러면 $X$는 닫힌점이 전혀 없는 국소뇌터 위상공간이다.
이 공간은 국소뇌터 스킴의 바탕 위상공간일 수 없다. Properties,
Lemma \ref{properties-lemma-locally-Noetherian-closed-point}를 보라.
\end{example}

\begin{lemma}
\label{topology-lemma-locally-Noetherian-locally-connected}
$X$를 국소뇌터 위상공간이라 하자. 그러면 $X$는 국소연결이다.
\end{lemma}

\begin{proof}
$x \in X$라 하자. $E$를 $x$의 근방이라 하자. $E$에 포함되는 $x$의
연결 근방을 찾아야 한다. 가정에 의해 $x$의 뇌터 근방 $E'$이 존재한다.
그러면 Lemma \ref{topology-lemma-Noetherian}에 의해 $E \cap E'$은 뇌터이다.
$E \cap E' = Y_1 \cup \ldots \cup Y_n$을 기약성분들로의 분해라 하자.
Lemma \ref{topology-lemma-Noetherian}를 보라.
$E'' = \bigcup_{x \in Y_i} Y_i$라 하자. 이는 $x$를 포함하는 $E \cap E'$의
연결 부분집합이다. 이는 $E \cap E'$의 열린집합
$E \cap E' \setminus (\bigcup_{x \not \in Y_i} Y_i)$를 포함하므로 $X$에서
$x$의 근방이다. 이로써 보조정리가 증명된다.
\end{proof}

\section{크룰 차원}
\label{topology-section-krull-dimension}

\begin{definition}
\label{topology-definition-Krull}
$X$를 위상공간이라 하자.
\begin{enumerate}
\item $X$의 {\it 기약 닫힌 부분집합들의 사슬}은
$Z_0 \subset Z_1 \subset \ldots \subset Z_n \subset X$ 꼴의 수열로서,
각 $Z_i$는 닫힌 기약 부분집합이고 $i = 0, \ldots, n - 1$에 대하여
$Z_i \not = Z_{i + 1}$인 것이다.
\item $X$의 기약 닫힌 부분집합들의 사슬
$Z_0 \subset Z_1 \subset \ldots \subset Z_n \subset X$의
{\it 길이}는 정수 $n$이다.
\item $X$의 {\it 차원}, 더 정확히는 {\it 크룰 차원} $\dim(X)$은 다음
공식으로 정의되는 $\{-\infty, 0, 1, 2, 3, \ldots, \infty\}$의 원소이다.
$$
\dim(X) =
\sup \{\text{lengths of chains of irreducible closed subsets}\}
$$
따라서 $\dim(X) = -\infty$일 필요충분조건은 $X$가 공집합인 것이다.
\item $x \in X$라 하자. {\it $x$에서 $X$의 크룰 차원}을
$$
\dim_x(X) = \min \{\dim(U), x\in U\subset X\text{ open}\}
$$
으로 정의한다. 이는 $U$가 $X$에서 $x$의 열린 근방들을 달릴 때의
$\dim(U)$의 최솟값이다.
\end{enumerate}
\end{definition}

\noindent
$U' \subset U \subset X$가 열려 있으면
$\dim(U') \leq \dim(U)$임에 유의하라. 따라서 $\dim_x(X) = d$이면
$x$는 $\dim(U) = \dim_x(X)$인 열린 근방들 $U$의 기본계를 갖는다.

\begin{lemma}
\label{topology-lemma-dimension-supremum-local-dimensions}
$X$를 위상공간이라 하자. 그러면 $\dim(X) = \sup \dim_x(X)$이며,
상한은 $X$의 점들 $x$를 달린다.
\end{lemma}

\begin{proof}
모든 $x \in X$에 대하여 $\dim(X) \geq \dim_x(X)$임은 분명하다
(Definition \ref{topology-definition-Krull} 뒤의 논의를 보라).
% P01-E1165: the frozen English transition is a sentence fragment.
따라서 한 방향의 부등식을 얻었다. 역방향을 위하여 $n \geq 0$이고
$\dim(X) \geq n$이라고 하자. 그러면 기약 닫힌 부분집합들의 사슬
$Z_0 \subset Z_1 \subset \ldots \subset Z_n \subset X$을 찾을 수 있다.
$x \in Z_0$를 택하자. $x$의 모든 열린 근방 $U$에 대하여 $U$ 안의
기약 닫힌 부분집합들의 사슬
$$
Z_0 \cap U \subset Z_1 \cap U \subset \ldots \subset Z_n \cap U
$$
를 얻는다. 실제로 집합 $U \cap Z_i$는 $U$에서 기약이고 닫혀 있으며
포함관계들은 엄격하다(세부사항은 생략한다. 힌트: $U \cap Z_i$의 폐포는
$Z_i$이다). 이로부터 $\dim_x(X) \geq n$임을 알 수 있고, 이는 다른
방향의 부등식을 증명한다.
\end{proof}

\begin{example}
\label{topology-example-Krull-Rn}
통상적인 유클리드 공간 $\mathbf{R}^n$의 크룰 차원은 $0$이다.
\end{example}

\begin{example}
\label{topology-example-krull-2set}
$X = \{s, \eta\}$라 하고 열린집합들이
$\{\emptyset, \{\eta\}, \{s, \eta\}\}$라고 하자. 이 경우 기약 닫힌
부분집합들의 한 극대 사슬은 $\{s\} \subset \{s, \eta\}$이다.
따라서 $\dim(X) = 1$이다. 이 예를 일반화하여 크룰 차원이 $n$인
$(n + 1)$원소 위상공간을 얻기는 쉽다.
\end{example}

\begin{definition}
\label{topology-definition-equidimensional}
$X$를 위상공간이라 하자. $X$의 모든 기약성분이 같은 차원을 가지면
$X$를 {\it 등차원}이라고 한다.
\end{definition}

\section{여차원과 카테너리 공간}
\label{topology-section-catenary-spaces}

\noindent
여기서는 기약 닫힌 부분집합들의 여차원만 정의한다.

\begin{definition}
\label{topology-definition-codimension}
$X$를 위상공간이라 하자. $Y \subset X$를 기약 닫힌 부분집합이라 하자.
$X$에서 $Y$의 {\it 여차원}은 $Y$에서 시작하는 $X$의 기약 닫힌
부분집합들의 사슬
$$
Y = Y_0 \subset Y_1 \subset \ldots \subset Y_e \subset X
$$
의 길이들 $e$의 상한이다. 이를 $\text{codim}(Y, X)$로 나타낸다.
\end{definition}

\noindent
여차원은 $\{0, 1, 2, \ldots\} \cup \{\infty\}$의 원소이다.
$\text{codim}(Y, X) < \infty$이면 모든 사슬은 극대 사슬로 확장될 수
있다(그러나 이 사슬들이 모두 같은 길이를 가져야 하는 것은 아니다).

\begin{lemma}
\label{topology-lemma-codimension-at-generic-point}
$X$를 위상공간이라 하자. $Y \subset X$를 기약 닫힌 부분집합이라 하자.
$Y \cap U$가 공집합이 아닌 열린 부분집합 $U \subset X$를 택하자. 그러면
$$
\text{codim}(Y, X) = \text{codim}(Y \cap U, U)
$$
이다.
\end{lemma}

\begin{proof}
대응 $T \mapsto \overline{T}$는 $U$의 기약 닫힌 부분집합들과 $U$와 만나는
$X$의 기약 닫힌 부분집합들 사이에서 포함관계를 보존하는 전단사 사상을
정의한다. 이로부터 보조정리가 쉽게 따른다. 세부사항은 생략한다.
\end{proof}

\begin{example}
\label{topology-example-Noetherian-infinite-codimension}
$X = [0, 1]$을 단위구간이라 하고 다음 위상을 주자. 집합들
$[0, 1]$, $(1 - 1/n, 1]$, $n \in \mathbf{N}$, 그리고 $\emptyset$이
열려 있다. 따라서 닫힌집합들은 $\emptyset$, $\{0\}$,
$[0, 1 - 1/n]$, $n > 1$, 그리고 $[0, 1]$이다. 이는 분명히 뇌터
위상공간이다. 그러나 기약 닫힌 부분집합 $Y = \{0\}$의 여차원은 무한하다.
즉 $\text{codim}(Y, X) = \infty$이다. 이를 보려면 닫힌집합들
$[0, 1 - 1/n]$이 모두 기약임을 관찰하기만 하면 된다.
\end{example}

\begin{definition}
\label{topology-definition-catenary}
$X$를 위상공간이라 하자. 기약 닫힌 부분집합들의 모든 쌍
$T \subset T'$에 대하여 $\text{codim}(T, T') < \infty$이고 기약 닫힌
부분집합들의 모든 극대 사슬
$$
T = T_0 \subset T_1 \subset \ldots \subset T_e = T'
$$
이 같은 길이(여차원과 같은 길이)를 가지면 $X$를 {\it 카테너리}라고 한다.
\end{definition}

\begin{lemma}
\label{topology-lemma-catenary}
$X$를 위상공간이라 하자. 다음은 서로 동치이다.
\begin{enumerate}
\item $X$는 카테너리이다.
\item $X$는 카테너리 공간들로 이루어진 열린 덮개를 갖는다.
\end{enumerate}
더욱이 이 경우 $X$의 임의의 국소닫힌 부분공간은 카테너리이다.
\end{lemma}

\begin{proof}
$X$가 카테너리이고 $U \subset X$가 열린 부분집합이라고 하자.
대응 $T \mapsto \overline{T}$는 $U$의 기약 닫힌 부분집합들과 $U$와
만나는 $X$의 기약 닫힌 부분집합들 사이에서 포함관계를 보존하는 전단사
사상을 정의한다. 이로부터 보조정리가 쉽게 따른다. 세부사항은 생략한다.
\end{proof}

\begin{lemma}
\label{topology-lemma-catenary-in-codimension}
$X$를 위상공간이라 하자. 다음은 서로 동치이다.
\begin{enumerate}
\item $X$는 카테너리이다.
\item 기약 닫힌 부분집합들의 모든 쌍 $Y \subset Y'$에 대하여
$\text{codim}(Y, Y') < \infty$이고, 기약 닫힌 부분집합들의 모든 삼중항
$Y \subset Y' \subset Y''$에 대하여
$$
\text{codim}(Y, Y'') = \text{codim}(Y, Y') + \text{codim}(Y', Y'').
$$
\end{enumerate}
\end{lemma}

\begin{proof}
% P01-E1166: the frozen English imperative is missing "us".
$X$가 카테너리라고 가정하자. Definition \ref{topology-definition-catenary}에 의해
기약 닫힌 부분집합들의 모든 쌍 $Y \subset Y'$에 대하여
$\text{codim}(Y,Y') < \infty$이다. $Y \subset Y' \subset Y''$을
$X$의 기약 닫힌 부분집합들의 삼중항이라 하자.
$$
Y = Y_0 \subset Y_1 \subset ... \subset Y_{e_1} = Y'
$$
을 $Y$와 $Y'$ 사이의 기약 닫힌 부분집합들의 극대 사슬이라 하되
$e_1 = \text{codim}(Y,Y')$라고 하자. 또한
$$
Y' = Y_{e_1} \subset Y_{e_1 + 1}\subset ... \subset Y_{e_1 + e_2} = Y''
$$
을 $Y'$과 $Y''$ 사이의 기약 닫힌 부분집합들의 극대 사슬이라 하되
$e_2 = \text{codim}(Y',Y'')$라고 하자. 두 사슬이 극대이므로 이어 붙인
사슬
$$
Y = Y_0\subset Y_1 \subset ... \subset Y_{e_1} = Y' =
Y_{e_1} \subset Y_{e_1+1}\subset ... \subset Y_{e_1+e_2}=Y''
$$
도 $Y$와 $Y''$ 사이에서 극대이고, 그 길이는 $e_1 + e_2$와 같다.
% P01-E1167: the frozen English uses "equals to" twice.
$X$가 카테너리이므로 각 극대 사슬은 여차원과 같은 길이를 갖는다.
따라서 (2)의 등식
$\text{codim}(Y,Y'') = e_1 + e_2 = \text{codim}(Y,Y') + \text{codim}(Y',Y'')$
이 성립한다.

\medskip\noindent
역으로, 귀납법으로 다음을 보인다. $Y = Y_1 \subset ... \subset
Y_n = Y'$이면 $ \text{codim}(Y,Y') = \text{codim}(Y_1,Y_2) + ... +
\text{codim}(Y_{n-1},Y_n)$이다. 따라서 극대 사슬들은 같은 길이를
가질 수밖에 없다.
\end{proof}

\section{준콤팩트 공간과 사상}
\label{topology-section-quasi-compact}

\noindent
‘콤팩트’라는 말은 하우스도르프 위상공간에만 쓰기로 한다. 대수기하학에서
나타나는 많은 공간은 하우스도르프가 아니기 때문이다.

\begin{definition}
\label{topology-definition-quasi-compact}
준콤팩트성.
\begin{enumerate}
\item 위상공간 $X$의 모든 열린 덮개가 유한 부분덮개를 가지면 $X$를
{\it 준콤팩트}라고 한다.
\item 임의의 준콤팩트 열린집합 $V \subset Y$의 역상 $f^{-1}(V)$가
준콤팩트이면 연속사상 $f : X \to Y$를 {\it 준콤팩트}라고 한다.
\item 포함사상 $Z \to X$가 준콤팩트이면 부분집합 $Z \subset X$를
{\it 레트로콤팩트}라고 한다.
\end{enumerate}
\end{definition}

\noindent
위상수학에 관한 많은 문헌은 준콤팩트이면서 하우스도르프인 공간을
{\it 콤팩트}라고 하지만, 다른 문헌은 하우스도르프 조건을 생략한다.
대수기하학에서의 혼동을 피하기 위해 우리는 준콤팩트라는 용어를 쓴다.
사상의 준콤팩트성 개념은 ‘고유사상’의 개념과 매우 다르다. 후자에서는
(닫힘성과 분리성 외에도) 표적의 임의의 준콤팩트 부분집합의 역상이
준콤팩트일 것을 요구하는 반면, 위 정의에서는 준콤팩트 {\it 열린}집합만
생각하기 때문이다.

\begin{lemma}
\label{topology-lemma-composition-quasi-compact}
준콤팩트 사상들의 합성은 준콤팩트이다.
\end{lemma}

\begin{proof}
정의에서 바로 따른다.
\end{proof}

\begin{lemma}
\label{topology-lemma-closed-in-quasi-compact}
준콤팩트 위상공간의 닫힌 부분집합은 준콤팩트이다.
\end{lemma}

\begin{proof}
$E \subset X$를 준콤팩트 공간 $X$의 닫힌 부분집합이라 하자.
$E = \bigcup V_j$를 열린 덮개라 하자. $V_j = E \cap U_j$를 만족하는
열린집합 $U_j \subset X$를 택하자. 그러면
$X = (X \setminus E) \cup \bigcup U_j$는 $X$의 열린 덮개이다. 따라서
어떤 $n$과 첨자 $j_i$에 대하여
$X = (X \setminus E) \cup U_{j_1} \cup \ldots \cup U_{j_n}$이다.
그러므로 원하는 대로 $E = V_{j_1} \cup \ldots \cup V_{j_n}$이다.
\end{proof}

\begin{lemma}
\label{topology-lemma-quasi-compact-in-Hausdorff}
$X$를 하우스도르프 위상공간이라 하자.
\begin{enumerate}
\item $E \subset X$가 준콤팩트이면 닫혀 있다.
\item $E_1, E_2 \subset X$가 서로소 준콤팩트 부분집합들이면
$U_1 \cap U_2 = \emptyset$이고 $E_i \subset U_i$인 열린집합들이
존재한다.
\end{enumerate}
\end{lemma}

\begin{proof}
(1)의 증명. $x \in X$, $x \not \in E$라 하자. 모든 $e \in E$에 대하여
$e \in V_e$, $x \in U_e$인 서로소 열린집합 $V_e$와 $U_e$를 찾을 수 있다.
$E \subset \bigcup V_e$이므로
$E \subset V_{e_1} \cup \ldots \cup V_{e_n}$인 유한 개의
$e_1, \ldots, e_n$을 찾을 수 있다. 그러면
$U = U_{e_1} \cap \ldots \cap U_{e_n}$은 $x$의 열린 근방이며
$V_{e_1} \cup \ldots \cup V_{e_n}$와 만나지 않는다. 특히 이는 $E$와
만나지 않는다. 따라서 $E$는 닫혀 있다.

\medskip\noindent
(2)의 증명. (1)의 증명에서 $x \in E_1$이 주어지면
$U_x \cap V_x = \emptyset$을 만족하는 열린 근방 $x \in U_x$와 열린집합
$E_2 \subset V_x$를 찾을 수 있음을 보았다. $E_1$은 준콤팩트이므로
$E_1 \subset U = U_{x_1} \cup \ldots \cup U_{x_n}$인 유한 개의
$x_i \in E_1$을 찾을 수 있다. $V = V_{x_1} \cap \ldots \cap V_{x_n}$로
택하면 증명이 끝난다.
\end{proof}

\begin{lemma}
\label{topology-lemma-closed-in-compact}
$X$를 준콤팩트 하우스도르프 공간이라 하고 $E \subset X$라 하자.
다음은 서로 동치이다. (a) $E$는 $X$에서 닫혀 있다. (b) $E$는
준콤팩트이다.
\end{lemma}

\begin{proof}
(a) $\Rightarrow$ (b)는 Lemma \ref{topology-lemma-closed-in-quasi-compact}이고,
(b) $\Rightarrow$ (a)는 Lemma
\ref{topology-lemma-quasi-compact-in-Hausdorff}이다.
\end{proof}

\noindent
다음은 실제로 준콤팩트 성질을 바꾸어 쓴 것이다.

\begin{lemma}
\label{topology-lemma-intersection-closed-in-quasi-compact}
$X$를 준콤팩트 위상공간이라 하자. $\{Z_\alpha\}_{\alpha \in A}$가
닫힌 부분집합들의 모임이고 각 유한 부분모임의 교집합이 공집합이 아니면
$\bigcap_{\alpha \in A} Z_\alpha$도 공집합이 아니다.
\end{lemma}

\begin{proof}
$\bigcap_{\alpha \in A} Z_{\alpha} = \emptyset$이라고 가정하자.
여집합을 취하면 $\bigcup_{\alpha \in A} (X\setminus Z_{\alpha})=X$이다.
부분집합들 $Z_\alpha$가 닫혀 있으므로
$\bigcup_{\alpha \in A} (X \setminus Z_{\alpha})$는 준콤팩트 공간 $X$의
열린 덮개이다. 따라서
$X = \bigcup_{\alpha \in J} (X\setminus Z_{\alpha})$가 되는 유한 부분집합
$J \subset A$가 존재한다. 그 여집합은 공집합이다. 즉
$\bigcap_{\alpha \in J} Z_{\alpha} = \emptyset$이다. 이는
$\{Z_{\alpha}\}_{\alpha \in J}$의 유한 부분모임 중
$\bigcap_{\alpha \in J} Z_{\alpha}=\emptyset$인 것이 존재함을 보이므로,
대우에 의해 결론이 따른다.
\end{proof}

\begin{lemma}
\label{topology-lemma-image-quasi-compact}
$f : X \to Y$를 위상공간들의 연속사상이라 하자.
\begin{enumerate}
\item $X$가 준콤팩트이면 $f(X)$도 준콤팩트이다.
\item $f$가 준콤팩트이면 $f(X)$는 레트로콤팩트이다.
\end{enumerate}
\end{lemma}

\begin{proof}
$f(X) = \bigcup V_i$가 열린 덮개이면
$X = \bigcup f^{-1}(V_i)$는 열린 덮개이다. 따라서 $X$가 준콤팩트이면
어떤 $i_1, \ldots, i_n \in I$에 대하여
$X = f^{-1}(V_{i_1}) \cup \ldots \cup f^{-1}(V_{i_n})$이고, 따라서
$f(X) = V_{i_1} \cup \ldots \cup V_{i_n}$이다. 이는 (1)을 증명한다.
$f$가 준콤팩트라고 하고 $V \subset Y$를 준콤팩트 열린집합이라 하자.
그러면 $f^{-1}(V)$는 준콤팩트이므로 (1)에 의해
$f(f^{-1}(V)) = f(X) \cap V$는 준콤팩트이다. 따라서 $f(X)$는
레트로콤팩트이다.
\end{proof}

\begin{lemma}
\label{topology-lemma-quasi-compact-closed-point}
$X$를 위상공간이라 하자. 다음을 가정한다.
\begin{enumerate}
\item $X$는 공집합이 아니다.
\item $X$는 준콤팩트이다.
\item $X$는 콜모고로프이다.
\end{enumerate}
그러면 $X$는 닫힌점을 갖는다.
\end{lemma}

\begin{proof}
$X$에서 한원소집합들의 모든 폐포로 이루어진 집합
$$
\mathcal{T} =
\{Z \subset X \mid Z = \overline{\{x\}} \text{ for some }x \in X\}
$$
을 생각하자. $X$가 공집합이 아니므로 이는 공집합이 아니다.
포함관계를 써서 $\mathcal{T}$를 부분순서 집합으로 만들자.
$Z_\alpha$, $\alpha \in A$가 $\mathcal{T}$의 전순서 부분집합이라고 하자.
Lemma \ref{topology-lemma-intersection-closed-in-quasi-compact}에 의해
$\bigcap_{\alpha \in A} Z_\alpha \not = \emptyset$이다. 따라서
$x \in \bigcap_{\alpha \in A} Z_\alpha$인 것이 존재하며,
$Z = \overline{\{x\}}\in \mathcal{T}$는 이 모임의 하계이다. 초른
보조정리에 의해 극소 원소 $Z \in \mathcal{T}$가 존재한다.
$X$가 콜모고로프이므로 어떤 $x$에 대하여 $Z = \{x\}$이고
$x \in X$는 닫힌점이라고 결론 내린다.
\end{proof}

\begin{lemma}
\label{topology-lemma-closed-points-quasi-compact}
$X$를 준콤팩트 콜모고로프 공간이라 하자. 그러면 $X$의 닫힌점들의 집합
$X_0$는 준콤팩트이다.
\end{lemma}

\begin{proof}
$X_0 = \bigcup U_{i, 0}$을 열린 덮개라 하자. 어떤 열린집합
$U_i \subset X$에 대하여 $U_{i, 0} = X_0 \cap U_i$라고 쓰자.
$\bigcup U_i$의 여집합을 $Z$라 하자. 이는 $X$의 닫힌 부분집합이므로
준콤팩트이고(Lemma \ref{topology-lemma-closed-in-quasi-compact}) 콜모고로프이다.
$Z$가 공집합이 아니면 Lemma \ref{topology-lemma-quasi-compact-closed-point}에
의해 닫힌점을 가지는데, 이는 $X_0 \subset \bigcup U_i$에 모순이다.
따라서 $Z = \emptyset$이고 $X = \bigcup U_i$이다. $X$가 준콤팩트이므로
이 덮개는 유한 부분덮개를 가지며 결론이 따른다.
\end{proof}

\begin{lemma}
\label{topology-lemma-connected-component-intersection}
$X$를 위상공간이라 하자. 다음을 가정한다.
\begin{enumerate}
\item $X$는 준콤팩트이다.
\item $X$는 준콤팩트 열린집합들로 이루어진 위상기저를 갖는다.
\item 두 준콤팩트 열린집합의 교집합은 준콤팩트이다.
\end{enumerate}
임의의 $x \in X$에 대하여 $x$를 포함하는 $X$의 연결성분은 $x$를
포함하는 $X$의 모든 열린 동시에 닫힌 부분집합들의 교집합이다.
\end{lemma}

\begin{proof}
$T$를 $x$를 포함하는 연결성분이라 하자. $S = \bigcap_{\alpha \in A}
Z_\alpha$를 $x$를 포함하는 $X$의 모든 열린 동시에 닫힌 부분집합
$Z_\alpha$의 교집합이라 하자. $S$가 $X$에서 닫혀 있음에 유의하라.
$Z_\alpha$들의 임의의 유한 교집합은 어떤 $Z_\alpha$이다.
$T$가 연결이고 $x \in T$이므로 $T \subset S$이다. $S$가 연결임을 보이면
충분하다. 그렇지 않으면 $S$에서 열리고 닫힌 $B$와 $C$에 대한 서로소 합
분해 $S = B \amalg C$가 존재한다. 특히 $B$와 $C$는 $X$에서 닫혀 있고,
Lemma \ref{topology-lemma-closed-in-quasi-compact}와 가정 (1)에 의해
준콤팩트이다. 가정 (2)에 의해
$B = S \cap U$와 $C = S \cap V$를 만족하는 준콤팩트 열린집합들
$U, V \subset X$가 존재한다(세부사항은 생략한다). 그러면
$U \cap V \cap S = \emptyset$이다. 따라서
$\bigcap_\alpha U \cap V \cap Z_\alpha = \emptyset$이다. 가정 (3)에 의해
교집합 $U \cap V$는 준콤팩트이다. Lemma
\ref{topology-lemma-intersection-closed-in-quasi-compact}에 의해 어떤
$\alpha' \in A$에 대하여
$U \cap V \cap Z_{\alpha'} = \emptyset$이다.
$X \setminus (U \cup V)$는 $S$와 서로소이고 $X$에서 닫혀 있으므로
준콤팩트이다. 따라서 같은 보조정리를 써서 어떤 $\alpha'' \in A$에 대하여
$Z_{\alpha''} \subset U \cup V$임을 알 수 있다.
% P01-E1172: the frozen proof does not explicitly choose the index alpha used next.
그러면 $Z_\alpha = Z_{\alpha'} \cap Z_{\alpha''}$는 $U \cup V$에 포함되고
$U \cap V$와 서로소이다. 따라서
$Z_\alpha = U \cap Z_\alpha \amalg V \cap Z_\alpha$는 두 열린 조각으로의
분해이고, 따라서 $U \cap Z_\alpha$와 $V \cap Z_\alpha$는 $X$에서 열리고
닫혀 있다. 그러므로 가령 $x \in B$이면 $S \subset U \cap Z_\alpha$이고,
$C = \emptyset$이라고 결론 내린다.
\end{proof}

\begin{lemma}
\label{topology-lemma-connected-component-intersection-compact-Hausdorff}
$X$를 위상공간이라 하고 준콤팩트이며 하우스도르프라고 가정하자.
임의의 $x \in X$에 대하여 $x$를 포함하는 $X$의 연결성분은 $x$를
포함하는 $X$의 모든 열린 동시에 닫힌 부분집합들의 교집합이다.
\end{lemma}

\begin{proof}
$T$를 $x$를 포함하는 연결성분이라 하자. $S = \bigcap_{\alpha \in A}
Z_\alpha$를 $x$를 포함하는 $X$의 모든 열린 동시에 닫힌 부분집합
$Z_\alpha$의 교집합이라 하자. $S$가 $X$에서 닫혀 있음에 유의하라.
$Z_\alpha$들의 임의의 유한 교집합은 어떤 $Z_\alpha$이다.
$T$가 연결이고 $x \in T$이므로 $T \subset S$이다. $S$가 연결임을 보이면
충분하다. 그렇지 않으면 $S$에서 열리고 닫힌 $B$와 $C$에 대한 서로소 합
분해 $S = B \amalg C$가 존재한다. 특히 $B$와 $C$는 $X$에서 닫혀 있고
Lemma \ref{topology-lemma-closed-in-quasi-compact}에 의해 준콤팩트이다.
Lemma \ref{topology-lemma-quasi-compact-in-Hausdorff}에 의해
$B \subset U$, $C \subset V$인 서로소 열린집합들 $U, V \subset X$가
존재한다.
% P01-E1173: the frozen complement formulas omit parentheses around U union V.
그러면 $X \setminus U \cup V$는 $X$에서 닫혀 있으므로
준콤팩트이다(Lemma \ref{topology-lemma-closed-in-quasi-compact}). Lemma
\ref{topology-lemma-intersection-closed-in-quasi-compact}에 의해 어떤 $\alpha$에
대하여 $(X \setminus U \cup V) \cap Z_\alpha = \emptyset$이다.
다시 말해 $Z_\alpha \subset U \cup V$이다. 따라서
$Z_\alpha = Z_\alpha \cap V \amalg Z_\alpha \cap U$는 두 열린 조각으로의
분해이고, 따라서 $U \cap Z_\alpha$와 $V \cap Z_\alpha$는 $X$에서 열리고
닫혀 있다. 그러므로 가령 $x \in B$이면 $S \subset U \cap Z_\alpha$이고,
$C = \emptyset$이라고 결론 내린다.
\end{proof}

\begin{lemma}
\label{topology-lemma-closed-union-connected-components}
$X$를 위상공간이라 하자. 다음을 가정한다.
\begin{enumerate}
\item $X$는 준콤팩트이다.
\item $X$는 준콤팩트 열린집합들로 이루어진 위상기저를 갖는다.
\item 두 준콤팩트 열린집합의 교집합은 준콤팩트이다.
\end{enumerate}
부분집합 $T \subset X$에 대하여 다음은 서로 동치이다.
\begin{enumerate}
\item[(a)] $T$는 $X$의 열린 동시에 닫힌 부분집합들의 교집합이다.
\item[(b)] $T$는 $X$에서 닫혀 있고 $X$의 연결성분들의 합집합이다.
\end{enumerate}
\end{lemma}

\begin{proof}
(a)가 (b)를 함의함은 분명하다. (b)를 가정하자. $x \in X$,
$x \not \in T$라 하자. $x \in C \subset X$를 $x$를 포함하는 $X$의
연결성분이라 하자. Lemma \ref{topology-lemma-connected-component-intersection}에
의해 $C = \bigcap V_\alpha$는 $C$를 포함하는 $X$의 모든 열린 동시에
닫힌 부분집합 $V_\alpha$의 교집합이다. 특히 임의의 쌍별 교집합
$V_\alpha \cap V_\beta$도 어떤 $V_\alpha$로 나타난다.
$T$가 $X$의 연결성분들의 합집합이므로 $C \cap T = \emptyset$이다.
따라서 $T \cap \bigcap V_\alpha = \emptyset$이다. $T$는 준콤팩트 공간의
닫힌 부분집합으로서 준콤팩트이므로(Lemma
\ref{topology-lemma-closed-in-quasi-compact} 참조), Lemma
\ref{topology-lemma-intersection-closed-in-quasi-compact}에 의해 어떤 $\alpha$에
대하여 $T \cap V_\alpha = \emptyset$이다. 이 $\alpha$에 대하여
$U_\alpha = X \setminus V_\alpha$는 $T$를 포함하고 $x$는 포함하지 않는
$X$의 열린 동시에 닫힌 부분집합이다. 보조정리가 따른다.
\end{proof}

\begin{lemma}
\label{topology-lemma-Noetherian-quasi-compact}
$X$를 뇌터 위상공간이라 하자.
\begin{enumerate}
\item 공간 $X$는 준콤팩트이다.
\item $X$의 임의의 부분집합은 레트로콤팩트이다.
\end{enumerate}
\end{lemma}

\begin{proof}
$X = \bigcup U_i$를 $X$의 열린 덮개라 하자. 첨자집합은 $I$이고 이 덮개는
유한 열린 덮개에 의한 세분을 갖지 않는다고 가정하자. $I$의 원소들
$i_1, i_2, \ldots $을 다음과 같이 귀납적으로 택하자.
$U_{i_{n + 1}}$이 $U_{i_1} \cup \ldots \cup U_{i_n}$에 포함되지 않도록
$i_{n + 1}$을 택한다. 그러면
% P01-E1174: the frozen chain omits parentheses around the accumulated union.
$X \supset (X \setminus U_{i_1}) \supset
(X \setminus U_{i_1} \cup U_{i_2}) \supset \ldots$는 닫힌 부분집합들의
엄격히 감소하는 무한수열이다. 이는 $X$가 뇌터라는 사실과 모순이다.
이로써 첫 명제가 증명된다. 이제 모든 $X$의 부분집합은 Lemma
\ref{topology-lemma-Noetherian}에 의해 뇌터이므로 둘째 명제도 분명하다.
\end{proof}

\begin{lemma}
\label{topology-lemma-quasi-compact-locally-Noetherian-Noetherian}
준콤팩트 국소뇌터 공간은 뇌터이다.
\end{lemma}

\begin{proof}
조건들에서 $X$가 뇌터 부분집합들로 이루어진 유한 덮개를 가짐이 바로
따르고, Lemma \ref{topology-lemma-finite-union-Noetherian}에 의해 $X$는 뇌터이다.
\end{proof}

\begin{lemma}[알렉산더 부분기저 정리]
\label{topology-lemma-subbase-theorem}
$X$를 위상공간이라 하고 $\mathcal{B}$를 $X$의 부분기저라 하자.
$\mathcal{B}$의 원소들로 이루어진 모든 $X$의 덮개가 유한 세분을 가지면
$X$는 준콤팩트이다.
\end{lemma}

\begin{proof}
$X$의 어떤 열린 덮개가 유한 세분을 갖지 않는다고 가정하자. 초른
보조정리를 사용하면 유한 세분을 갖지 않는 극대 열린 덮개
$X = \bigcup_{i \in I} U_i$를 택할 수 있다(세부사항은 생략한다).
다시 말해, 열린 부분집합 $U \subset X$가 $U_i$ 중 어느 것과도 같지
않으면 덮개 $X = U \cup \bigcup_{i \in I} U_i$는 유한 세분을 갖는다.
$U_i \in \mathcal{B}$인 첨자들의 집합을 $I' \subset I$라 하자.
$\bigcup_{i \in I'} U_i \not = X$이다. 그렇지 않으면 $\mathcal{B}$에
대한 가정으로부터 $X$를 덮는 유한 세분을 얻기 때문이다.
$x \in X$, $x \not \in \bigcup_{i \in I'} U_i$를 택하자.
$x \in U_i$인 $i \in I$를 택하자.
$x \in V_1 \cap \ldots \cap V_n \subset U_i$인
$V_1, \ldots, V_n \in \mathcal{B}$를 택하자. $\mathcal{B}$가 $X$의
부분기저이므로 이는 가능하다. $V_j$는 어느 $U_i$와도 같지 않음에
유의하라. 택한 덮개의 극대성에 의해 각 $j$에 대하여
$X = V_j \cup U_{i_{j, 1}} \cup \ldots \cup U_{i_{j, n_j}}$가 되는
$i_{j, 1}, \ldots, i_{j, n_j} \in I$가 존재한다.
$V_1 \cap \ldots \cap V_n \subset U_i$이므로
$X = U_i \cup \bigcup U_{i_{j, l}}$라고 결론 내리며, 이는 모순이다.
\end{proof}

\section{국소 준콤팩트 공간}
\label{topology-section-locally-quasi-compact}

\noindent
한 점의 근방이 반드시 열린집합일 필요는 없음을 상기하자.

\begin{definition}
\label{topology-definition-locally-quasi-compact}
위상공간 $X$의 모든 점이 준콤팩트 근방들의 기본계를 가지면 이를
{\it 국소 준콤팩트}라고 한다.\footnote{이 용어법은 표준적이지 않을 수
있다. 문헌에서 쓰이는 다른 개념으로는 다음이 있다. (1) 모든 점이 어떤
준콤팩트 근방을 갖는다. (2) 모든 점이 닫힌 준콤팩트 근방을 갖는다.
스킴에서는 모든 점이 열린 준콤팩트 근방들의 기본계를 갖는다는 성질이
성립한다.}
\end{definition}

\noindent
문헌에서 {\it 국소 콤팩트 공간}이라는 말은 흔히 다음 보조정리에
나오는 공간을 가리킨다.

\begin{lemma}
\label{topology-lemma-locally-quasi-compact-Hausdorff}
하우스도르프 공간이 국소 준콤팩트일 필요충분조건은 모든 점이 준콤팩트
근방을 갖는 것이다.
\end{lemma}

\begin{proof}
$X$를 하우스도르프 공간이라 하자. $x \in X$이고
$x \in E \subset X$가 준콤팩트 근방이라 하자. 그러면 Lemma
\ref{topology-lemma-quasi-compact-in-Hausdorff}에 의해 $E$는 닫혀 있다.
$x \in U \subset X$가 $x$의 열린 근방이라고 하자. 그러면
$Z = E \setminus U$는 $x$를 포함하지 않는 $E$의 닫힌 부분집합이다.
Lemma \ref{topology-lemma-quasi-compact-in-Hausdorff}에 의해 $E$의 서로소인 두
열린 부분집합 $W, V \subset E$를 $x \in V$이고 $Z \subset W$가 되도록 찾을 수
있다. 따라서 $\overline{V} \subset E$는 $E \cap U$에 포함되는 $x$의
닫힌 근방이다. 또한 $\overline{V}$는 $E$의 닫힌 부분집합이므로
준콤팩트이다(Lemma \ref{topology-lemma-closed-in-quasi-compact}). 이런 식으로
$x$의 준콤팩트 근방들로 이루어진 기본계를 얻는다.
\end{proof}

\begin{lemma}[베르 범주 정리]
\label{topology-lemma-baire-category-locally-compact}
$X$를 국소 준콤팩트 하우스도르프 공간이라 하자.
$U_n \subset X$, $n \geq 1$을 조밀한 열린 부분집합들이라 하자. 그러면
$\bigcap_{n \geq 1} U_n$은 $X$에서 조밀하다.
\end{lemma}

\begin{proof}
$U_n$을 $\bigcap_{i = 1, \ldots, n} U_i$로 바꾸어
$U_1 \supset U_2 \supset \ldots$라고 가정해도 좋다.
$x \in X$라 하자. $x$가 $\bigcap_{n \geq 1} U_n$의 폐포에 속함을 보일
것이다. 이를 위해 $E$를 $x$의 근방이라 하자.
$E \cap \bigcap_{n \geq 1} U_n$이 공집합이 아님을 보일 때 $E$를 더 작은
근방으로 바꾸어도 좋다. $E$를 더 작은 근방으로 바꾸어 $E$가
준콤팩트라고 가정할 수 있다.

\medskip\noindent
$x_0 = x$ 및 $E_0 = E$로 놓자. 아래에서 점
$x_i \in E_{i - 1} \cap U_i$와 $E_i \subset E_{i - 1} \cap U_i$인
$x_i$의 준콤팩트 근방 $E_i$를 귀납적으로 택한다. $X$가 하우스도르프이므로
부분집합들 $E_i \subset X$는 닫혀 있다(Lemma
\ref{topology-lemma-quasi-compact-in-Hausdorff}). $E_i$들은 공집합도 아니므로
$\bigcap_{i \geq 1} E_i$는 공집합이 아니다(Lemma
\ref{topology-lemma-intersection-closed-in-quasi-compact}).
$\bigcap_{i \geq 1} E_i \subset E \cap \bigcap_{n \geq 1} U_n$이므로
보조정리가 증명된다.

\medskip\noindent
기초 단계 $i = 0$은 위에서 처리하였다. 이제 귀납 단계를 보자.
$E_{i - 1}$은 $x_{i - 1}$의 근방이므로
$x_{i - 1} \in W \subset E_{i - 1}$인 열린 부분집합을 찾을 수 있다.
$U_i$가 $X$에서 조밀하므로 $W \cap U_i$는 공집합이 아니다.
임의의 $x_i \in W \cap U_i$를 택하자. 국소 준콤팩트 공간의 정의에 의해
$W \cap U_i$에 포함되는 $x_i$의 준콤팩트 근방 $E_i$를 찾을 수 있다.
그러면 원하는 대로 $E_i \subset E_{i - 1} \cap U_i$이다.
\end{proof}

\begin{lemma}
\label{topology-lemma-relatively-compact-refinement}
$X$를 하우스도르프 준콤팩트 공간이라 하자.
$X = \bigcup_{i \in I} U_i$를 열린 덮개라 하자. 그러면 모든 $i$에 대하여
$\overline{V_i} \subset U_i$인 열린 덮개
$X = \bigcup_{i \in I} V_i$가 존재한다.
\end{lemma}

\begin{proof}
$x \in X$라 하자. $x \in U_{i(x)}$인 $i(x) \in I$를 택하자.
$X \setminus U_{i(x)}$와 $\{x\}$는 $X$의 서로소 닫힌 부분집합들이므로,
Lemmas \ref{topology-lemma-closed-in-quasi-compact}와
\ref{topology-lemma-quasi-compact-in-Hausdorff}에 의해 폐포가
$X \setminus U_{i(x)}$와 서로소인 $x$의 열린 근방 $U_x$가 존재한다.
따라서 $\overline{U_x} \subset U_{i(x)}$이다. $X$가 준콤팩트이므로
$X = U_{x_1} \cup \ldots \cup U_{x_m}$인 유한한 점들의 목록
$x_1, \ldots, x_m$이 존재한다. 이제
% P01-E1180: the frozen union subscript binds i instead of the varying index j.
$V_i = \bigcup_{i = i(x_j)} U_{x_j}$로 놓으면 증명이 끝난다.
\end{proof}

\begin{lemma}
\label{topology-lemma-refine-covering}
$X$를 하우스도르프 준콤팩트 공간이라 하자.
$X = \bigcup_{i \in I} U_i$를 열린 덮개라 하자. 정수 $p \geq 0$이 주어지고,
$I$의 모든 $(p + 1)$-튜플 $i_0, \ldots, i_p$에 대하여 열린 덮개
$U_{i_0} \cap \ldots \cap U_{i_p} = \bigcup W_{i_0 \ldots i_p, k}$가
주어졌다고 하자. 그러면 열린 덮개 $X = \bigcup_{j \in J} V_j$와 사상
$\alpha : J \to I$가 존재하여 $\overline{V_j} \subset U_{\alpha(j)}$이고,
각 $V_{j_0} \cap \ldots \cap V_{j_p}$는 어떤 $k$에 대하여
$W_{\alpha(j_0) \ldots \alpha(j_p), k}$에 포함된다.
\end{lemma}

\begin{proof}
$X$가 준콤팩트이므로 $I$가 유한한 경우로 환원된다(세부사항은 생략한다).
$I$가 유한한 경우를 $p$에 대한 귀납법으로 증명한다. 기초 단계
$p = 0$에서는 열린 덮개 $X = \bigcup W_{i_0, k}$를 세분하는 Lemma
\ref{topology-lemma-relatively-compact-refinement}의 덮개를 택하면 되므로 자명하다.

\medskip\noindent
귀납 단계를 보자. 보조정리가 $p - 1$에 대하여 증명되었다고 가정하자.
$I$의 모든 $p$-튜플 $i'_0, \ldots, i'_{p - 1}$에 대하여
$U_{i'_0} \cap \ldots \cap U_{i'_{p - 1}} =
\bigcup W_{i'_0 \ldots i'_{p - 1}, k}$를
$\{i'_0, \ldots, i'_{p - 1}\} = \{i_0, \ldots, i_p\}$인
$(p + 1)$-튜플들에 대한 덮개
$U_{i_0} \cap \ldots \cap U_{i_p} = \bigcup W_{i_0 \ldots i_p, k}$의
공통 세분이라 하자(여기서 등식은 집합의 등식이다).
$I$가 유한하므로 이러한 덮개는 유한 개뿐이다. 귀납가정에 의해 이
열린집합들에 대한 해, 즉 $X = \bigcup V_j$와
$\alpha : J \to I$가 존재한다. 이 단계에서 덮개
$X = \bigcup_{j \in J} V_j$와 $\alpha$는
$\overline{V_j} \subset U_{\alpha(j)}$를 만족하며,
$\alpha(j_0), \ldots, \alpha(j_p)$에 중복이 있으면 각
$V_{j_0} \cap \ldots \cap V_{j_p}$가 어떤 $k$에 대하여
$W_{\alpha(j_0) \ldots \alpha(j_p), k}$에 포함된다. 물론 $J$도
유한하다고 가정할 수 있고 그렇게 하겠다.

\medskip\noindent
쌍마다 서로 다른 $i_0, \ldots, i_p \in I$를 고정하자.
$i_0 = \alpha(j_0), \ldots, i_p = \alpha(j_p)$이고
$V_{j_0} \cap \ldots \cap V_{j_p}$가 어떤 $k$에 대해서도
$W_{\alpha(j_0) \ldots \alpha(j_p), k}$에 포함되지 {\bf 않는}
$j_0, \ldots, j_p \in J$의 $(p + 1)$-튜플들을 생각하자. 그러한
$(p + 1)$-튜플의 개수를 $N$이라 하자. $N$을 줄이는 방법을 보이겠다.
다음 포함관계가 성립한다.
$$
\overline{V_{j_0}} \cap \ldots \cap \overline{V_{j_p}} \subset
U_{i_0} \cap \ldots \cap U_{i_p} = \bigcup W_{i_0 \ldots i_p, k}
$$
따라서 좌변이 $\bigcup_{k \in K} W_{i_0 \ldots i_p, k}$에 포함되도록 하는
유한집합 $K = \{k_1, \ldots, k_t\}$를 찾을 수 있다. 이제 열린 덮개
$$
V_{j_0} =
(V_{j_0} \setminus (\overline{V_{j_1}} \cap \ldots \cap \overline{V_{j_p}}))
\cup (\bigcup\nolimits_{k \in K} V_{j_0} \cap W_{i_0 \ldots i_p, k})
$$
를 생각한다. 우변의 첫 열린집합과
$V_{j_1} \cap \ldots \cap V_{j_p}$의 교집합은 공집합이고, 우변의 다른
열린집합들 $V_{j_0, k}$는
% P01-E1181: the frozen proof uses V_{j_0,k} without defining that symbol.
% P01-E1182: the frozen formula omits an intersection operator after V_{j_1}.
$V_{j_0, k} \cap V_{j_1} \ldots \cap V_{j_p} \subset
W_{\alpha(j_0) \ldots \alpha(j_p), k}$를 만족한다.
$J' = J \amalg K$로 놓자. $j \in J$에 대하여 $j \not = j_0$이면
$V'_j = V_j$로 놓고
$V'_{j_0} =
V_{j_0} \setminus (\overline{V_{j_1}} \cap \ldots \cap \overline{V_{j_p}})$로
놓는다. $k \in K$에 대해서는 $V'_k = V_{j_0, k}$로 놓는다. 마지막으로
사상 $\alpha' : J' \to I$는 $J$ 위에서는 $\alpha$로 주어지고 $K$의 모든
원소를 $i_0$로 보낸다. 간단히 확인하면 이 치환에 의해 $N$이 하나
감소한다. 이 절차를 $N$번 되풀이하면 $N = 0$인 상황에 이른다.

\medskip\noindent
증명을 마치기 위해, 다음 조건을 만족하는 쌍마다 서로 다른 원소들로
이루어진 $(p + 1)$-튜플 $i_0, \ldots, i_p \in I$의 개수 $M$에 대해
귀납법을 적용한다. 그 조건은
$i_0 = \alpha(j_0), \ldots, i_p = \alpha(j_p)$이면서
$V_{j_0} \cap \ldots \cap V_{j_p}$가 어떤 $k$에 대해서도
$W_{\alpha(j_0) \ldots \alpha(j_p), k}$에 포함되지 {\bf 않는}
$j_0, \ldots, j_p \in J$의 $(p + 1)$-튜플이 존재하는 것이다. 이를 위해
앞 문단에서 수행한 연산이 $M$을 늘리지 않음을 주장한다. 이는 사상
$\alpha' : J' \to I$가 $V'_{j'} \subset V_{\beta(j')}$를 만족하는 사상
$\beta : J' \to J$를 거쳐 분해된다는 사실에서 형식적으로 따른다.
\end{proof}

\begin{lemma}
\label{topology-lemma-lift-covering-of-a-closed}
$X$를 하우스도르프 국소 준콤팩트 공간이라 하자.
$Z \subset X$를 준콤팩트인(따라서 닫힌) 부분집합이라 하자. 정수
$p \geq 0$, 집합 $I$, 각 $i \in I$에 대한 열린집합 $U_i \subset X$,
그리고 $I$의 각 $(p + 1)$-튜플 $i_0, \ldots, i_p$에 대한 열린집합
$W_{i_0 \ldots i_p} \subset U_{i_0} \cap \ldots \cap U_{i_p}$가
주어지고 다음을 만족한다고 하자.
\begin{enumerate}
\item $Z \subset \bigcup U_i$이다.
\item 모든 $i_0, \ldots, i_p$에 대하여
$W_{i_0 \ldots i_p} \cap Z = U_{i_0} \cap \ldots \cap U_{i_p} \cap Z$이다.
\end{enumerate}
그러면 다음 조건들을 만족하는 $X$의 열린집합들 $V_i$가 존재한다.
$Z \subset \bigcup V_i$이고, 모든 $i$에 대하여
$\overline{V_i} \subset U_i$이며, 모든 $(p + 1)$-튜플
$i_0, \ldots, i_p$에 대하여
$V_{i_0} \cap \ldots \cap V_{i_p} \subset W_{i_0 \ldots i_p}$이다.
\end{lemma}

\begin{proof}
$Z$가 준콤팩트이므로 $I$가 유한한 경우로 환원된다(세부사항은 생략한다).
$X$가 국소 준콤팩트이고 $Z$가 준콤팩트이므로 준콤팩트인 근방
$Z \subset E$를 찾을 수 있다. 즉, $E$는 준콤팩트이고 $X$ 안에서
$Z$의 열린 근방을 포함한다. $X$를 $E$로 바꾼 뒤 결과를 증명하면 원래
결과도 따르므로, $X$가 준콤팩트라고 가정해도 좋다.

\medskip\noindent
$I$가 유한하고 $X$가 준콤팩트인 경우를 $p$에 대한 귀납법으로 증명한다.
기초 단계는 $p = 0$인 경우이다. 이때
$X = (X \setminus Z) \cup \bigcup W_i$이다. Lemma
\ref{topology-lemma-relatively-compact-refinement}에 의해 열린집합
$V_i \subset W_i$ 및 $V \subset X \setminus Z$로 이루어진 덮개
$X = V \cup \bigcup V_i$를 모든 $i$에 대하여
$\overline{V_i} \subset W_i$가 되도록 찾을 수 있다. 그러면 보조정리에서
제시한 문제의 해를 얻는다.

\medskip\noindent
귀납 단계를 보자. 보조정리가 $p - 1$에 대하여 증명되었다고 가정하자.
$W_{j_0 \ldots j_{p - 1}}$를
$\{j_0, \ldots, j_{p - 1}\} = \{i_0, \ldots, i_p\}$인 모든
$W_{i_0 \ldots i_p}$의 교집합으로 놓는다(여기서 등식은 집합의 등식이다).
귀납가정에 의해 이 열린집합들에 대한 해, 이를테면
$V_i \subset U_i$가 존재한다. $W_{j_0 \ldots j_{p - 1}}$의 선택에 의해,
어떤 $0 \leq a < b \leq p$에 대하여 $i_a = i_b$인 모든
$(p + 1)$-튜플 $i_0, \ldots, i_p$에 대해
$V_{i_0} \cap \ldots \cap V_{i_p} \subset W_{i_0 \ldots i_p}$가 성립한다.
따라서 쌍마다 서로 다른 원소들로 이루어진 어떤 $(p + 1)$-튜플
$i_0, \ldots, i_p$에 대하여
$V_{i_0} \cap \ldots \cap V_{i_p} \not \subset W_{i_0 \ldots i_p}$인
경우에만 $V_i$의 선택을 수정하면 된다. 이 경우
$$
T =
\overline{V_{i_0} \cap \ldots \cap V_{i_p} \setminus W_{i_0 \ldots i_p}}
\subset 
\overline{V_{i_0}} \cap \ldots \cap \overline{V_{i_p}} \setminus
W_{i_0 \ldots i_p}
$$
는 $U_{i_0} \cap \ldots \cap U_{i_p}$에 포함되고 $Z$와 만나지 않는
$X$의 닫힌 부분집합이다. 따라서 $V_{i_0}$를
$V_{i_0} \setminus T$로 바꾸어 문제를 ``고칠'' 수 있다. 문제가 되는
각 튜플에 대하여 이를 유한 번 되풀이하면 보조정리가 증명된다.
\end{proof}

\begin{lemma}
\label{topology-lemma-lift-covering-of-quasi-compact-hausdorff-subset}
$X$를 위상공간이라 하자. $Z \subset X$를, $Z$의 임의의 두 점이 $X$에서
서로소인 열린 근방들을 갖는 준콤팩트 부분집합이라 하자. 정수
$p \geq 0$, 집합 $I$, 각 $i \in I$에 대한 열린집합 $U_i \subset X$,
그리고 $I$의 각 $(p + 1)$-튜플 $i_0, \ldots, i_p$에 대한 열린집합
$W_{i_0 \ldots i_p} \subset U_{i_0} \cap \ldots \cap U_{i_p}$가
주어지고 다음을 만족한다고 하자.
\begin{enumerate}
\item $Z \subset \bigcup U_i$이다.
\item 모든 $i_0, \ldots, i_p$에 대하여
$W_{i_0 \ldots i_p} \cap Z = U_{i_0} \cap \ldots \cap U_{i_p} \cap Z$이다.
\end{enumerate}
그러면 다음을 만족하는 $X$의 열린집합들 $V_i$가 존재한다.
\begin{enumerate}
\item $Z \subset \bigcup V_i$이다.
\item 모든 $i$에 대하여 $V_i \subset U_i$이다.
\item 모든 $i$에 대하여 $\overline{V_i} \cap Z \subset U_i$이다.
\item 모든 $(p + 1)$-튜플 $i_0, \ldots, i_p$에 대하여
$V_{i_0} \cap \ldots \cap V_{i_p} \subset W_{i_0 \ldots i_p}$이다.
\end{enumerate}
\end{lemma}

\begin{proof}
$Z$가 준콤팩트이므로 $I$가 유한한 경우로 환원된다(세부사항은 생략한다).
$I$가 유한한 경우를 $p$에 대한 귀납법으로 증명한다.

\medskip\noindent
기초 단계는 $p = 0$인 경우이다. $z \in Z \cap U_i$ 및
$z' \in Z \setminus U_i$에 대하여 서로소인 $X$의 열린집합
$z \in V_{z, z'}$와 $z' \in W_{z, z'}$가 존재한다.
$Z \setminus U_i$가 준콤팩트이므로(Lemma
\ref{topology-lemma-closed-in-quasi-compact}),
$Z \setminus U_i \subset W_{z, z'_1} \cup \ldots \cup W_{z, z'_r}$인
유한 개의 점 $z'_1, \ldots, z'_r$을 택할 수 있다. 그러면
$V_z = V_{z, z'_1} \cap \ldots \cap V_{z, z'_r} \cap U_i$는 $U_i$에
포함되는 $z$의 열린 근방이고
$\overline{V_z} \cap Z \subset U_i$라는 성질을 갖는다. $z$와 $i$가
임의적이었고 $Z$가 준콤팩트이므로, 유한한 목록
$z_1, i_1, \ldots, z_t, i_t$와 열린집합들
$V_{z_j} \subset U_{i_j}$를
$\overline{V_{z_j}} \cap Z \subset U_{i_j}$ 및
$Z \subset \bigcup V_{z_j}$가 되도록 찾을 수 있다. 이제
$V_i = W_i \cap (\bigcup_{j : i = i_j} V_{z_j})$로 놓으면
$p = 0$인 경우의 문제가 풀린다.

\medskip\noindent
귀납 단계를 보자. 보조정리가 $p - 1$에 대하여 증명되었다고 가정하자.
$W_{j_0 \ldots j_{p - 1}}$를
$\{j_0, \ldots, j_{p - 1}\} = \{i_0, \ldots, i_p\}$인 모든
$W_{i_0 \ldots i_p}$의 교집합으로 놓는다(여기서 등식은 집합의 등식이다).
귀납가정에 의해 이 열린집합들에 대한 해, 이를테면
$V_i \subset U_i$가 존재한다. $W_{j_0 \ldots j_{p - 1}}$의 선택에 의해,
어떤 $0 \leq a < b \leq p$에 대하여 $i_a = i_b$인 모든
$(p + 1)$-튜플 $i_0, \ldots, i_p$에 대해
$V_{i_0} \cap \ldots \cap V_{i_p} \subset W_{i_0 \ldots i_p}$가 성립한다.
따라서 쌍마다 서로 다른 원소들로 이루어진 어떤 $(p + 1)$-튜플
$i_0, \ldots, i_p$에 대하여
$V_{i_0} \cap \ldots \cap V_{i_p} \not \subset W_{i_0 \ldots i_p}$인
경우에만 $V_i$의 선택을 수정하면 된다. 이 경우
$$
T =
\overline{V_{i_0} \cap \ldots \cap V_{i_p} \setminus W_{i_0 \ldots i_p}}
\subset
\overline{V_{i_0}} \cap \ldots \cap \overline{V_{i_p}} \setminus
W_{i_0 \ldots i_p}
$$
는 위 열린집합들 $V_i$의 성질 (3)에 의해 $Z$와 만나지 않는 $X$의 닫힌
부분집합이다. 따라서 $V_{i_0}$를 $V_{i_0} \setminus T$로 바꾸어 문제를
``고칠'' 수 있다. 문제가 되는 각 튜플에 대하여 이를 유한 번 되풀이하면
보조정리가 증명된다.
\end{proof}

\section{공간의 극한}
\label{topology-section-limits}

\noindent
위상공간의 범주는 곱을 갖는다. 구체적으로, $I$가 집합이고 각
$i \in I$에 대하여 위상공간 $X_i$가 주어지면
$\prod_{i \in I} X_i$에 {\it 곱위상}을 준다. 이 위상의 기저로는
$U_i \subset X_i$가 열려 있고 거의 모든 $i$에 대하여 $U_i = X_i$인
꼴의 집합 $\prod U_i$들을 사용한다.

\medskip\noindent
위상공간의 범주는 등화자를 갖는다. 구체적으로,
$a, b : X \to Y$가 위상공간들의 사상이면 $a$와 $b$의 등화자는
유도위상을 준 부분집합 $\{x \in X \mid a(x) = b(x)\} \subset X$이다.

\begin{lemma}
\label{topology-lemma-limits}
위상공간의 범주는 극한을 가지며, 집합으로 가는 망각 함자는 이 극한들과
교환한다.
\end{lemma}

\begin{proof}
이는 위 논의와 Categories, Lemma
\ref{categories-lemma-limits-products-equalizers}에서 따른다. 위 설명으로부터
망각 함자가 극한과 교환함도 따른다. 이를 보는 다른 방법은 Categories,
Lemma \ref{categories-lemma-adjoint-exact}를 적용하고, 망각 함자가 집합에
대응하는 이산 위상공간을 대응시키는 함자를 왼쪽 수반함자로 갖는다는
사실을 이용하는 것이다.
\end{proof}

\begin{lemma}
\label{topology-lemma-describe-limits}
$\mathcal{I}$를 쌍대여과 범주라 하자. $i \mapsto X_i$를
$\mathcal{I}$ 위의 위상공간 도표라 하자. $X = \lim X_i$를 사영사상
$f_i : X \to X_i$를 갖는 극한이라 하자.
\begin{enumerate}
\item $X$의 임의의 열린집합은 어떤 부분집합 $J \subset I$와 열린집합들
% P01-E1185: the frozen formula uses the undefined set I rather than Ob(mathcal I).
$U_j \subset X_j$에 대하여 $\bigcup_{j \in J} f_j^{-1}(U_j)$ 꼴이다.
\item $X$의 임의의 준콤팩트 열린집합은 어떤 $i$와 어떤 열린집합
$U_i \subset X_i$에 대하여 $f_i^{-1}(U_i)$ 꼴이다.
\end{enumerate}
\end{lemma}

\begin{proof}
위에서 주어진 극한의 구성에 따르면 $X \subset \prod X_i$이고 그 위상은
유도위상이다. $\prod X_i$의 위상의 한 기저는 $U_i \subset X_i$가 열려
있고 거의 모든 $i$에 대하여 $U_i = X_i$인 열린집합들 $\prod U_i$로
이루어진다. $U_{i_j} \not = X_{i_j}$인 대상들이
$i_1, \ldots, i_n \in \Ob(\mathcal{I})$라고 하자. 그러면
$$
X \cap \prod U_i = f_{i_1}^{-1}(U_{i_1}) \cap \ldots \cap f_{i_n}^{-1}(U_{i_n})
$$
위상공간의 일반적인 극한에서는 이들이 $X$ 위상의 기저를 이룬다. 그러나
보조정리의 가정처럼 $\mathcal{I}$가 쌍대여과이면
$j \in \Ob(\mathcal{I})$와 사상들 $j \to i_l$,
$l = 1, \ldots, n$을 택할 수 있다. 다음과 같이 놓자.
$$
U_j =
(X_j \to X_{i_1})^{-1}(U_{i_1}) \cap \ldots \cap
(X_j \to X_{i_n})^{-1}(U_{i_n})
$$
그러면 $X \cap \prod U_i = f_j^{-1}(U_j)$임이 분명하다. 따라서 $X$의
임의의 열린집합 $W$에 대하여 집합 $A$, 사상
$\alpha : A \to \Ob(\mathcal{I})$, 그리고 열린집합들
$U_a \subset X_{\alpha(a)}$가 존재하여
$W = \bigcup f_{\alpha(a)}^{-1}(U_a)$가 된다. $J = \Im(\alpha)$로 놓고,
$j \in J$에 대하여 $U_j = \bigcup_{\alpha(a) = j} U_a$로 놓으면
$W = \bigcup_{j \in J} f_j^{-1}(U_j)$임을 알 수 있다. 이로써 (1)이
증명된다.

\medskip\noindent
(2)를 보이기 위해 $\bigcup_{j \in J} f_j^{-1}(U_j)$가 준콤팩트라고 하자.
그러면 어떤 $j_1, \ldots, j_m \in J$에 대하여 이는
$f_{j_1}^{-1}(U_{j_1}) \cup \ldots \cup f_{j_m}^{-1}(U_{j_m})$와 같다.
$\mathcal{I}$가 쌍대여과이므로 $i \in \Ob(\mathcal{I})$와 사상들
$i \to j_l$, $l = 1, \ldots, m$을 택할 수 있다. 다음과 같이 놓자.
$$
U_i =
(X_i \to X_{j_1})^{-1}(U_{j_1}) \cup \ldots \cup
(X_i \to X_{j_m})^{-1}(U_{j_m})
$$
그러면 원하는 대로 우리의 열린집합은 $f_i^{-1}(U_i)$와 같다.
\end{proof}

\begin{lemma}
\label{topology-lemma-characterize-limit}
$\mathcal{I}$를 쌍대여과 범주라 하자. $i \mapsto X_i$를
$\mathcal{I}$ 위의 위상공간 도표라 하자. $X$를 다음 조건을 만족하는
위상공간이라 하자.
\begin{enumerate}
\item 집합으로서 $X = \lim X_i$이다(사영사상들을 $f_i$라 쓰자).
\item $i \in \Ob(\mathcal{I})$ 및 열린집합 $U_i \subset X_i$에 대한 집합들
$f_i^{-1}(U_i)$가 $X$의 위상의 기저를 이룬다.
\end{enumerate}
그러면 $X$는 위상공간으로서 $X_i$들의 극한이다.
\end{lemma}

\begin{proof}
이는 Lemma \ref{topology-lemma-describe-limits}의 극한 위상에 대한 설명에서 따른다.
\end{proof}

\begin{theorem}[티호노프]
\label{topology-theorem-tychonov}
준콤팩트 공간들의 곱은 준콤팩트이다.
\end{theorem}

\begin{proof}
$I$를 집합이라 하고 각 $i \in I$에 대하여 $X_i$를 준콤팩트 위상공간이라
하자. $X = \prod X_i$로 놓자. $\mathcal{B}$를
$U_i \times \prod_{i' \in I, i' \not = i} X_{i'}$ 꼴의 $X$의 부분집합들로
이루어진 집합이라 하자. 여기서 $U_i \subset X_i$는 열려 있다. 구성에 의해
이 모음은 $X$의 위상의 부분기저이다. Lemma
\ref{topology-lemma-subbase-theorem}에 의해 $\mathcal{B}$의 원소들 $B_j$로 이루어진
임의의 덮개 $X = \bigcup_{j \in J} B_j$가 유한 세분을 가짐을 보이면
충분하다. $J = \coprod J_i$로 분해하여 $j \in J_i$이면
$B_j = U_j \times \prod_{i' \not = i} X_{i'}$이고
$U_j \subset X_i$가 열려 있도록 할 수 있다. 만약
$X_i = \bigcup_{j \in J_i} U_j$이면 유한 세분이 존재하고,
$X = \bigcup_{j \in J} B_j$가 유한 세분을 갖는다고 결론 내린다.
그렇지 않으면 모든 $i$에 대하여 $\bigcup_{j \in J_i} U_j$에 속하지 않는
점 $x_i \in X_i$를 택할 수 있다. 그러나 그러면 점
$x = (x_i)_{i \in I}$는 $\bigcup_{j \in J} B_j$에 포함되지 않는 $X$의
원소이므로 모순이다.
\end{proof}

\noindent
다음 보조정리는 공간들 $X_i$가 하우스도르프라는 가정을 버리면 성립하지
않는다. Examples, Section
\ref{examples-section-lim-not-quasi-compact}를 보라.

\begin{lemma}
\label{topology-lemma-inverse-limit-quasi-compact}
$\mathcal{I}$를 범주라 하고 $i \mapsto X_i$를 위상공간의 범주에서
$\mathcal{I}$ 위의 도표라 하자. 각 $X_i$가 준콤팩트 하우스도르프이면
$\lim X_i$는 준콤팩트이다.
\end{lemma}

\begin{proof}
$\lim X_i$가 $\prod X_i$의 부분공간임을 상기하자. Theorem
\ref{topology-theorem-tychonov}에 의해 이 곱은 준콤팩트이다. 따라서
$\lim X_i$가 $\prod X_i$의 닫힌 부분공간임을 보이면 충분하다(Lemma
\ref{topology-lemma-closed-in-quasi-compact}). $\varphi : j \to k$가
$\mathcal{I}$의 사상이면, 대응하는 연속사상 $X_j \to X_k$의 그래프를
$\Gamma_\varphi \subset X_j \times X_k$라 쓰자. Lemma
\ref{topology-lemma-graph-closed}에 의해 이 그래프는 닫혀 있다.
$\lim X_i$가 다음 닫힌 부분집합들의 교집합임은 분명하다.
$$
\Gamma_\varphi \times \prod\nolimits_{l \not = j, k} X_l
\subset \prod X_i
$$
따라서 결과가 따른다.
\end{proof}

\noindent
다음 보조정리는 Categories, Lemma
\ref{categories-lemma-nonempty-limit}를 일반화하고, Lemma
\ref{topology-lemma-intersection-closed-in-quasi-compact}를 부분적으로 일반화한다.

\begin{lemma}
\label{topology-lemma-nonempty-limit}
$\mathcal{I}$를 쌍대여과 범주라 하고 $i \mapsto X_i$를 위상공간의 범주에서
$\mathcal{I}$ 위의 도표라 하자. 각 $X_i$가 준콤팩트이고 하우스도르프이며
공집합이 아니면 $\lim X_i$도 공집합이 아니다.
\end{lemma}

\begin{proof}
Lemma \ref{topology-lemma-inverse-limit-quasi-compact}의 증명에서
$X = \lim X_i$는 준콤팩트 공간 $\prod X_i$ 안에서 다음 닫힌 부분집합들의
교집합임을 보았다.
$$
Z_\varphi = \Gamma_\varphi \times \prod\nolimits_{l \not = j, k} X_l
$$
여기서 $\varphi : j \to k$는 $\mathcal{I}$의 사상이고
$\Gamma_\varphi \subset X_j \times X_k$는 대응하는 사상
$X_j \to X_k$의 그래프이다. 따라서 Lemma
\ref{topology-lemma-intersection-closed-in-quasi-compact}에 의해 이 부분집합들의
임의의 유한 교집합이 공집합이 아님을 보이면 충분하다.
$\varphi_t : j_t \to k_t$, $t = 1, \ldots, n$이 $\mathcal{I}$의 사상들로
이루어진 유한한 모음이라고 하자. $\mathcal{I}$가 쌍대여과이므로 대상
$j$와 각 $t$에 대한 사상 $\psi_t : j \to j_t$를 택할 수 있다.
(a) $j_t = j_{t'}$, (b) $j_t = k_{t'}$, (c) $k_t = k_{t'}$ 중 하나를
만족하는 각 쌍 $t, t'$에 대하여 두 사상 $j \to l$을 얻는다. 여기서
(a), (b)의 경우에는 $l = j_t$이고 (c)의 경우에는 $l = k_t$이다.
$\mathcal{I}$가 쌍대여과이고 쌍 $(t, t')$가 유한 개뿐이므로, 이러한 모든
쌍 $(t, t')$에 대하여 이 두 사상을 등화하는 사상 $j' \to j$를 택할 수
있다. 원소 $x \in X_{j'}$를 택하고 각 $t$에 대하여 $x_{j_t}$ 및 각각
$x_{k_t}$를 $X_{j'} \to X_j \to X_{j_t}$ 및 각각
$X_{j'} \to X_j \to X_{j_t} \to X_{k_t}$에 의한 $x$의 상이라 하자.
어떤 $t$에 대해서도 $j_t$나 $k_t$와 같지 않은 각 첨자
$l \in \Ob(\mathcal{I})$에 대하여 임의의 원소 $x_l \in X_l$을 택한다
(선택공리를 사용한다). 그러면 구성에 의해
$(x_i)_{i \in \Ob(\mathcal{I})}$는 교집합
$$
Z_{\varphi_1} \cap \ldots \cap Z_{\varphi_n}
$$
에 속하므로 증명이 끝난다.
\end{proof}

\section{구성가능 집합}
\label{topology-section-constructible}

\begin{definition}
\label{topology-definition-constructible}
$X$를 위상공간이라 하고 $E \subset X$를 $X$의 부분집합이라 하자.
\begin{enumerate}
\item $E$가 $X$에서 열린 동시에 레트로콤팩트인 $U, V \subset X$에 대한
$U \cap V^c$ 꼴의 부분집합들의 유한 합집합이면, $E$가 $X$에서
{\it 구성가능}하다고 한다.\footnote{EGA I 제2판 \cite{EGA1-second}에서는
이를 ``전역 구성가능'' 집합이라 불렀고, 여기서 국소 구성가능 집합이라
부르는 것을 ``구성가능'' 집합이라 불렀다.}
\item 열린 덮개 $X = \bigcup V_i$가 존재하여 각 $E \cap V_i$가 $V_i$에서
구성가능하면, $E$가 $X$에서 {\it 국소 구성가능}하다고 한다.
\end{enumerate}
\end{definition}

\begin{lemma}
\label{topology-lemma-constructible}
구성가능 집합들의 모음은 유한 교집합, 유한 합집합, 여집합에 대하여
닫혀 있다.
\end{lemma}

\begin{proof}
$U_1$, $U_2$가 $X$에서 열린 동시에 레트로콤팩트이면 $U_1 \cup U_2$도
그러함에 유의하라. 실제로 $X$의 두 준콤팩트 부분집합의 합집합은
준콤팩트이다. $U_1 \cap U_2$도 레트로콤팩트이다. 이를 보이기 위해
$U \subset X$가 준콤팩트 열린집합이라고 하자. $U_2$가 $X$에서
레트로콤팩트이므로 $U_2 \cap U$는 준콤팩트이고, 이어서 $U_1$이 $X$에서
레트로콤팩트이므로 $U_1 \cap (U_2 \cap U)$도 준콤팩트라고 결론 내린다.
이로부터 구성가능 집합의 여집합이 구성가능하고, 구성가능 집합들의 유한
합집합이 구성가능하며, 구성가능 집합들의 유한 교집합이 구성가능함을
형식적으로 보일 수 있다.
\end{proof}

\begin{lemma}
\label{topology-lemma-inverse-images-constructibles}
$f : X \to Y$를 위상공간들의 연속사상이라 하자. $Y$의 모든 레트로콤팩트
열린 부분집합의 역상이 $X$에서 레트로콤팩트이면, 구성가능 집합의 역상도
구성가능하다.
\end{lemma}

\begin{proof}
이는 $f^{-1}(U \cap V^c) = f^{-1}(U) \cap f^{-1}(V)^c$라는 사실과
구성가능 집합의 정의에서 따른다.
\end{proof}

\begin{lemma}
\label{topology-lemma-open-immersion-constructible-inverse-image}
$U \subset X$를 열린집합이라 하자. 구성가능 집합 $E \subset X$에 대하여
교집합 $E \cap U$는 $U$에서 구성가능하다.
\end{lemma}

\begin{proof}
$V \subset X$가 $X$에서 레트로콤팩트 열린집합이라고 하자. Lemma
\ref{topology-lemma-inverse-images-constructibles}에 의해 $V \cap U$가 $U$에서
레트로콤팩트임을 보이면 충분하다. 이를 보이기 위해 $W \subset U$가 열린
준콤팩트 집합이라고 하자. 그러면 $W$는 $X$에서도 열리고 준콤팩트이다.
따라서 $V$가 $X$에서 레트로콤팩트이므로
$V \cap W = V \cap U \cap W$는 준콤팩트이다.
\end{proof}

\begin{lemma}
\label{topology-lemma-quasi-compact-open-immersion-constructible-image}
$U \subset X$를 레트로콤팩트 열린집합이라 하고 $E \subset U$라 하자.
$E$가 $U$에서 구성가능하면 $E$는 $X$에서 구성가능하다.
\end{lemma}

\begin{proof}
$V, W \subset U$가 $U$에서 레트로콤팩트 열린집합들이라고 하자. 그러면
$V, W$는 $X$에서도 레트로콤팩트 열린집합들이다(Lemma
\ref{topology-lemma-composition-quasi-compact}). 따라서
$V \cap (U \setminus W) = V \cap (X \setminus W)$는 $X$에서
구성가능하다. $U$의 모든 구성가능 부분집합은
$V \cap (U \setminus W)$ 꼴의 부분집합들의 유한 합집합이므로 결론이
따른다.
\end{proof}

\begin{lemma}
\label{topology-lemma-collate-constructible}
$X$를 위상공간이라 하고 $E \subset X$를 부분집합이라 하자.
$X = V_1 \cup \ldots \cup V_m$을 레트로콤팩트 열린집합들로 이루어진 유한
덮개라 하자. $E$가 $X$에서 구성가능할 필요충분조건은 각
$j = 1, \ldots, m$에 대하여 $E \cap V_j$가 $V_j$에서 구성가능한 것이다.
\end{lemma}

\begin{proof}
$E$가 $X$에서 구성가능하면 Lemma
\ref{topology-lemma-open-immersion-constructible-inverse-image}에 의해 모든 $j$에
대하여 $E \cap V_j$가 $V_j$에서 구성가능하다. 역으로, 각
$j = 1, \ldots, m$에 대하여 $E \cap V_j$가 $V_j$에서 구성가능하다고 하자.
% P01-E1190: the frozen finite union has no indexing condition or range.
그러면 $E = \bigcup E \cap V_j$는 Lemma
\ref{topology-lemma-quasi-compact-open-immersion-constructible-image}에 의해
구성가능한 집합들의 유한 합집합이고, 따라서 구성가능하다.
\end{proof}

\begin{lemma}
\label{topology-lemma-intersect-constructible-with-closed}
$X$를 위상공간이라 하고 $Z \subset X$를 $X \setminus Z$가 준콤팩트인 닫힌
부분집합이라 하자. 구성가능 집합 $E \subset X$에 대하여 교집합
$E \cap Z$는 $Z$에서 구성가능하다.
\end{lemma}

\begin{proof}
$V \subset X$가 $X$에서 레트로콤팩트 열린집합이라고 하자. Lemma
\ref{topology-lemma-inverse-images-constructibles}에 의해 $V \cap Z$가 $Z$에서
레트로콤팩트임을 보이면 충분하다. 이를 보이기 위해 $W \subset Z$가 열린
준콤팩트 집합이라고 하자. 부분집합
$W' = W \cup (X \setminus Z)$는 준콤팩트이고 열려 있으며
$W = Z \cap W'$이다. 따라서 $V$가 $X$에서 레트로콤팩트이므로
$V \cap Z \cap W = V \cap Z \cap W'$는 준콤팩트 열린집합
$V \cap W'$의 닫힌 부분집합이다. 그러므로 Lemma
\ref{topology-lemma-closed-in-quasi-compact}에 의해 $V \cap Z \cap W$는
준콤팩트이다.
\end{proof}

\begin{lemma}
\label{topology-lemma-intersect-constructible-with-retrocompact}
$X$를 위상공간이라 하고 $T \subset X$를 부분집합이라 하자. 다음을
가정하자.
\begin{enumerate}
\item $T$는 $X$에서 레트로콤팩트이다.
\item 준콤팩트 열린집합들이 $X$의 위상의 기저를 이룬다.
\end{enumerate}
그러면 구성가능 집합 $E \subset X$에 대하여 교집합 $E \cap T$는 $T$에서
구성가능하다.
\end{lemma}

\begin{proof}
$V \subset X$가 $X$에서 레트로콤팩트 열린집합이라고 하자. Lemma
\ref{topology-lemma-inverse-images-constructibles}에 의해 $V \cap T$가 $T$에서
레트로콤팩트임을 보이면 충분하다. 이를 보이기 위해 $W \subset T$가 열린
준콤팩트 집합이라고 하자. 가정 (2)에 의해 $W = T \cap W'$인 준콤팩트
열린집합 $W' \subset X$를 찾을 수 있다(세부사항은 생략한다). 따라서
$V \cap T \cap W = V \cap T \cap W'$는 $V$가 $X$에서
레트로콤팩트이므로 준콤팩트 열린집합 $V \cap W'$와 $T$의 교집합이다.
그러므로 $V \cap T \cap W$는 준콤팩트이다.
\end{proof}

\begin{lemma}
\label{topology-lemma-closed-constructible-image}
$Z \subset X$를 여집합이 레트로콤팩트 열린집합인 닫힌 부분집합이라 하자.
$E \subset Z$라 하자. $E$가 $Z$에서 구성가능하면 $E$는 $X$에서
구성가능하다.
\end{lemma}

\begin{proof}
$V \subset Z$가 $Z$에서 레트로콤팩트 열린집합이라고 하자. $X$의 열린
부분집합 $\tilde V = V \cup (X \setminus Z)$를 생각하자.
$W \subset X$를 준콤팩트 열린집합이라 하자. 그러면
$$
W \cap \tilde V =
\left(V \cap W\right) \cup \left((X \setminus Z) \cap W\right).
$$
첫 부분은 준콤팩트이다. 실제로 $V \cap W = V \cap (Z \cap W)$이고,
$(Z \cap W)$는 $Z$에서 준콤팩트 열린집합이며(Lemma
\ref{topology-lemma-closed-in-quasi-compact}), $V$는 $Z$에서 레트로콤팩트이다.
둘째 부분은 $(X \setminus Z)$가 $X$에서 레트로콤팩트이므로
준콤팩트이다. 이로써 $\tilde V$가 $X$에서 레트로콤팩트임을 알 수 있다.
따라서 $V_1, V_2 \subset Z$가 레트로콤팩트 열린집합들이면
$$
V_1 \cap (Z \setminus V_2) = \tilde V_1 \cap (X \setminus \tilde V_2)
$$
는 $X$에서 구성가능하다. $Z$의 모든 구성가능 부분집합은
$V_1 \cap (Z \setminus V_2)$ 꼴의 부분집합들의 유한 합집합이므로 결론이
따른다.
\end{proof}

\begin{lemma}
\label{topology-lemma-constructible-is-retrocompact}
$X$를 위상공간이라 하자. $X$의 모든 구성가능 부분집합은
레트로콤팩트이다.
\end{lemma}

\begin{proof}
$U_i, V_i$가 $X$에서 레트로콤팩트 열린집합이고
$E = \bigcup_{i = 1, \ldots, n} U_i \cap V_i^c$라고 하자.
$W \subset X$를 준콤팩트 열린집합이라 하자. 그러면
$E \cap W = \bigcup_{i = 1, \ldots, n} U_i \cap V_i^c \cap W$이다.
따라서 $U, V$가 레트로콤팩트 열린집합이고 $W$가 준콤팩트 열린집합일 때
$U \cap V^c \cap W$가 준콤팩트임을 보이면 충분하다. 이는
$U \cap V^c \cap W$가 준콤팩트 집합 $U \cap W$의 닫힌 부분집합이기
때문에 성립하며, Lemma \ref{topology-lemma-closed-in-quasi-compact}가 적용된다.
\end{proof}

\noindent
질문: $X$가 준콤팩트 위상공간이라고만 가정해도 다음 보조정리가 성립하는가?
Lemma \ref{topology-lemma-intersect-constructible-with-closed}와 비교하라.

\begin{lemma}
\label{topology-lemma-intersect-constructible-with-constructible}
$X$를 위상공간이라 하고 $X$가 준콤팩트 열린집합들로 이루어진 기저를
갖는다고 가정하자. $X$에서 구성가능한 $E, E'$에 대하여 교집합
$E \cap E'$은 $E$에서 구성가능하다.
\end{lemma}

\begin{proof}
Lemmas \ref{topology-lemma-intersect-constructible-with-retrocompact}와
\ref{topology-lemma-constructible-is-retrocompact}를 결합하라.
\end{proof}

\begin{lemma}
\label{topology-lemma-constructible-in-constructible}
$X$를 위상공간이라 하고 $X$가 준콤팩트 열린집합들로 이루어진 기저를
갖는다고 가정하자. $E$가 $X$에서 구성가능하고 $F \subset E$가 $E$에서
구성가능하다고 하자. 그러면 $F$는 $X$에서 구성가능하다.
\end{lemma}

\begin{proof}
$X$의 임의의 레트로콤팩트 부분집합 $T$는 준콤팩트 열린집합들로 이루어진
유도위상의 기저를 가짐을 관찰하라. 특히 이는 임의의 구성가능 부분집합에
대하여 성립한다(Lemma \ref{topology-lemma-constructible-is-retrocompact}).
$U_i, V_i \subset X$가 레트로콤팩트 열린집합이고
$E_i = U_i \cap V_i^c$일 때 $E = E_1 \cup \ldots \cup E_n$으로 쓰자.
Lemma \ref{topology-lemma-intersect-constructible-with-constructible}에 의해
$E_i = E \cap E_i$는 $E$에서 구성가능함에 유의하라. 따라서 Lemma
\ref{topology-lemma-intersect-constructible-with-constructible}에 의해
$F \cap E_i$는 $E_i$에서 구성가능하다. 그러므로
$U, V \subset X$가 레트로콤팩트 열린집합이고 $E = U \cap V^c$인 경우의
보조정리를 증명하면 충분하다. 이 경우 포함사상 $E \subset X$는 다음과
같이 합성된다.
$$
E = U \cap V^c \to U \to X
$$
이제 첫 포함사상에는 Lemma \ref{topology-lemma-closed-constructible-image}를,
둘째 포함사상에는 Lemma
\ref{topology-lemma-quasi-compact-open-immersion-constructible-image}를 적용할 수
있다.
\end{proof}

\begin{lemma}
\label{topology-lemma-locally-closed-constructible-image}
$X$를 준콤팩트 열린집합들로 이루어진 기저를 가지며 임의의 두 준콤팩트
열린집합의 교집합이 준콤팩트인 준콤팩트 위상공간이라 하자.
$T \subset X$를 국소닫힌 부분집합이라 하고, $T$가 준콤팩트이며 $T^c$가
$X$에서 레트로콤팩트라고 하자. 그러면 $T$는 $X$에서 구성가능하다.
\end{lemma}

\begin{proof}
$T$가 준콤팩트이고 $\overline{T}$에서 열려 있음에 유의하라. 준콤팩트
열린집합들로 이루어진 기저를 사용하면 $U$가 $X$의 준콤팩트 열린집합인
$T = U \cap \overline{T}$로 쓸 수 있다. 그러면 $T^c$가 $X$에서
레트로콤팩트이므로 $V = U \setminus T = U \cap T^c$는 $U$에서
레트로콤팩트이다. 따라서 $V$는 준콤팩트이다. 임의의 두 준콤팩트
열린집합의 교집합이 준콤팩트이므로 $X$의 모든 준콤팩트 열린집합은
레트로콤팩트이다. 따라서 $U$ 및 $V = U \setminus T$가 $X$의
레트로콤팩트 열린집합인 $T = U \cap V^c$를 얻는다. 특히 $T$는 $X$에서
구성가능하다.
\end{proof}

\begin{lemma}
\label{topology-lemma-collate-constructible-from-constructible}
$X$를 준콤팩트 열린집합들로 이루어진 위상의 기저를 갖는 위상공간이라
하고 $E \subset X$를 부분집합이라 하자.
$X = E_1 \cup \ldots \cup E_m$을 구성가능 부분집합들로 이루어진 유한
덮개라 하자. $E$가 $X$에서 구성가능할 필요충분조건은 각
$j = 1, \ldots, m$에 대하여 $E \cap E_j$가 $E_j$에서 구성가능한 것이다.
\end{lemma}

\begin{proof}
Lemmas \ref{topology-lemma-intersect-constructible-with-constructible}와
\ref{topology-lemma-constructible-in-constructible}를 결합하라.
\end{proof}

\begin{lemma}
\label{topology-lemma-generic-point-in-constructible}
$X$를 위상공간이라 하자. $Z \subset X$가 기약이라고 하자.
$E \subset X$를 국소닫힌 부분집합들의 유한 합집합이라 하자(예를 들어
$E$가 구성가능한 경우이다). 다음은 서로 동치이다.
\begin{enumerate}
\item 교집합 $E \cap Z$는 $Z$의 열린 조밀 부분집합을 포함한다.
\item 교집합 $E \cap Z$는 $Z$에서 조밀하다.
\end{enumerate}
$Z$가 일반점 $\xi$를 가지면, 이는 다음과도 동치이다.
\begin{enumerate}
\item[(3)] $\xi \in E$이다.
\end{enumerate}
\end{lemma}

\begin{proof}
(1) $\Rightarrow$ (2)는 분명하다. (2)를 가정하자. $E \cap Z$는 $Z$의
국소닫힌 부분집합들 $Z_i$의 유한 합집합임에 유의하라. $Z$가 기약이므로
$Z_i$ 중 하나는 $Z$에서 조밀해야 한다. 그러면 이 $Z_i$는 자신의 폐포에서
열려 있으므로 $Z$에서 조밀한 열린집합이다. 따라서 (1)이 성립한다.

\medskip\noindent
$\xi \in Z$가 일반점이라고 하자. 동치인 조건 (1), (2)가 성립하면
$\xi \in E$이다. 역으로 $\xi \in E$이면 $\xi \in E \cap Z$이고, 따라서
$E \cap Z$는 $Z$에서 조밀하다.
\end{proof}

\section{구성가능 집합과 뇌터 공간}
\label{topology-section-constructible-Noetherian}

\begin{lemma}
\label{topology-lemma-constructible-Noetherian-space}
$X$를 뇌터 위상공간이라 하자. $X$의 구성가능 집합들은 정확히 $X$의
국소닫힌 부분집합들의 유한 합집합들이다.
\end{lemma}

\begin{proof}
이는 Lemma \ref{topology-lemma-Noetherian-quasi-compact}에서 바로 따른다.
\end{proof}

\begin{lemma}
\label{topology-lemma-constructible-map-Noetherian}
$f : X \to Y$를 뇌터 위상공간들의 연속사상이라 하자.
$E \subset Y$가 $Y$에서 구성가능하면 $f^{-1}(E)$는 $X$에서
구성가능하다.
\end{lemma}

\begin{proof}
이는 Lemma \ref{topology-lemma-constructible-Noetherian-space}와 연속사상의 정의에서
바로 따른다.
\end{proof}

\begin{lemma}
\label{topology-lemma-characterize-constructible-Noetherian}
$X$를 뇌터 위상공간이라 하고 $E \subset X$를 부분집합이라 하자. 다음은
서로 동치이다.
\begin{enumerate}
\item $E$는 $X$에서 구성가능하다.
\item 모든 기약 닫힌 부분집합 $Z \subset X$에 대하여 교집합 $E \cap Z$는
$Z$의 공집합이 아닌 열린집합을 포함하거나 $Z$에서 조밀하지 않다.
\end{enumerate}
\end{lemma}

\begin{proof}
$E$가 구성가능하고 $Z \subset X$가 기약 닫힌집합이라고 하자. 그러면
Lemma \ref{topology-lemma-constructible-map-Noetherian}에 의해 $E \cap Z$는 $Z$에서
구성가능하다. 따라서 $E \cap Z$는 $Z$의 공집합이 아닌 국소닫힌
부분집합들 $T_i$의 유한 합집합이다. $T_i$ 가운데 어느 것도 $Z$에서
열려 있지 않으면 $E \cap Z$가 $Z$에서 조밀하지 않음이 분명하다. 이로써
(1)이 (2)를 함의함을 알 수 있다.

\medskip\noindent
역으로 (2)가 성립한다고 가정하자. $E \cap Y$가 $Y$에서 구성가능하지 않은
$X$의 닫힌 부분집합 $Y$들로 이루어진 집합 $\mathcal{S}$를 생각하자.
% P01-E1191: frozen smallest is stronger than the minimal element supplied by Noetherianity.
$\mathcal{S} \not = \emptyset$이면 $X$가 뇌터이므로 이는 최소 원소 $Y$를
갖는다. $Y = Y_1 \cup \ldots \cup Y_r$를 $Y$의 기약성분 분해라 하자.
Lemma \ref{topology-lemma-Noetherian}를 보라. $r > 1$이면 각 $Y_i \cap E$는 $Y_i$에서
구성가능하고, 따라서 $Y_i$의 국소닫힌 부분집합들의 유한 합집합이다.
그러므로 $E \cap Y$ 역시 $Y$의 국소닫힌 부분집합들의 유한 합집합이고,
Lemma \ref{topology-lemma-constructible-Noetherian-space}에 의해 $E \cap Y$가 $Y$에서
구성가능하다고 결론 내린다. 이는 모순이므로 $r = 1$이다. $r = 1$이면
$Y$는 기약이고, 가정 (2)에 의해 $E \cap Y$는 (a) $Y$의 열린집합 $V$를
포함하거나 (b) $Y$에서 조밀하지 않다. 경우 (a)에서는 $Y$의 극소성에
의해 $E \cap (Y \setminus V)$가 $Y \setminus V$의 국소닫힌 부분집합들의
유한 합집합임을 알 수 있다. 따라서 $E \cap Y$는 $Y$의 국소닫힌
부분집합들의 유한 합집합이고 Lemma
\ref{topology-lemma-constructible-Noetherian-space}에 의해 구성가능하다. 이는
모순이므로 경우 (b)이어야 한다. 경우 (b)에서는 어떤 진닫힌 부분집합
$Y' \subset Y$에 대하여 $E \cap Y = E \cap Y'$임을 알 수 있다. $Y$의
극소성에 의해 $E \cap Y'$은 $Y'$의 국소닫힌 부분집합들의 유한
합집합이며, 따라서 $E \cap Y' = E \cap Y$는 $Y$의 국소닫힌 부분집합들의
유한 합집합이고 Lemma \ref{topology-lemma-constructible-Noetherian-space}에 의해
구성가능함을 알 수 있다. 이 모순으로 보조정리의 증명이 끝난다.
\end{proof}

\begin{lemma}
\label{topology-lemma-constructible-neighbourhood-Noetherian}
$X$를 뇌터 위상공간이라 하고 $x \in X$라 하자. $E \subset X$가 $X$에서
구성가능하다고 하자. 다음은 서로 동치이다.
\begin{enumerate}
\item $E$는 $x$의 근방이다.
\item $x$를 포함하는 $X$의 모든 기약 닫힌 부분집합 $Y$에 대하여
교집합 $E \cap Y$는 $Y$에서 조밀하다.
\end{enumerate}
\end{lemma}

\begin{proof}
(1)이 (2)를 함의함은 분명하다. (2)를 가정하자. $x$를 포함하고
$E \cap Y$가 $Y$에서 $x$의 근방이 아닌 $X$의 닫힌 부분집합 $Y$들로
이루어진 집합 $\mathcal{S}$를 생각하자. $\mathcal{S} \not = \emptyset$이면
$X$가 뇌터이므로 이는 극소 원소 $Y$를 갖는다.
$Y = Y_1 \cup Y_2$이고 $Y_1$, $Y_2$가 더 작은 두 공집합이 아닌 닫힌
부분집합이라고 하자. $i = 1, 2$에 대하여 $x \in Y_i$이면
$Y_i \cap E$는 $Y_i$에서 $x$의 근방이고, 따라서 $Y \cap E$가 $Y$에서
$x$의 근방이라고 결론 내리는데 이는 모순이다. 가령 $x \in Y_1$이지만
$x \not\in Y_2$이면 $Y_1 \cap E$는 $Y_1$에서 $x$의 근방이고, 따라서
$Y$에서도 그러하여 역시 모순이다. 그러므로 $Y$는 기약 닫힌집합이다.
가정 (2)에 의해 $E \cap Y$는 $Y$에서 조밀하다. 따라서 Lemma
\ref{topology-lemma-characterize-constructible-Noetherian}에 의해 $E \cap Y$는
$Y$의 열린집합 $V$를 포함한다. $x \in V$이면 $E \cap Y$는 $Y$에서
$x$의 근방이므로 모순이다. $x \not \in V$이면
$Y' = Y \setminus V$는 $x$를 포함하는 $Y$의 진닫힌 부분집합이다. $Y$의
극소성에 의해 $E \cap Y'$은 $Y'$에서 $x$의 열린 근방
$V' \subset Y'$를 포함한다. 그러나 그러면 $V' \cup V$는 $E$에 포함되는
$Y$에서 $x$의 열린 근방이므로 모순이다. 이 모순으로 보조정리의 증명이
끝난다.
\end{proof}

\begin{lemma}
\label{topology-lemma-characterize-open-Noetherian}
$X$를 뇌터 위상공간이라 하고 $E \subset X$를 부분집합이라 하자. 다음은
서로 동치이다.
\begin{enumerate}
\item $E$는 $X$에서 열려 있다.
\item $X$의 모든 기약 닫힌 부분집합 $Y$에 대하여 교집합 $E \cap Y$는
공집합이거나 $Y$의 공집합이 아닌 열린집합을 포함한다.
\end{enumerate}
\end{lemma}

\begin{proof}
이는 Lemmas \ref{topology-lemma-characterize-constructible-Noetherian}와
\ref{topology-lemma-constructible-neighbourhood-Noetherian}에서 형식적으로 따른다.
\end{proof}

\section{proper 사상의 특징짓기}
\label{topology-section-proper}

\noindent
보통 위상수학에서 proper 사상의 개념을 논의하는 절을 포함한다. 위상공간의
연속사상이 보편적으로 닫혀 있고 분리되어 있으면 이를 {\it proper}라고
정의한다. 이는 대수기하학의 proper 사상의 정의와 잘 부합하지만 부르바키의
정의와는 다르다. 여기서 정의한 위상공간의 proper 사상에 대해서는 별도의
가정 없이 proper 기저변환 정리(Cohomology, Theorem
\ref{cohomology-theorem-proper-base-change})가 성립한다. 또한
$\mathbf{C}$ 위의 유한형 스킴들의 사상 $f : X \to Y$가 주어졌을 때 다음이
성립한다. $f$가 스킴의 사상으로서 proper일 필요충분조건은 고전적 위상을
준 $\mathbf{C}$-점 위의 연속사상 $f : X(\mathbf{C}) \to Y(\mathbf{C})$가
proper인 것이다. 이는 \cite[Exp. XII, Prop.\ 3.2(v)]{SGA1}에 설명되어
있으며, 거기에는 위상수학의 proper성을 분리성을 덧붙인 부르바키의 개념으로
취한다는 각주도 있다.

\medskip\noindent
분리됨, 보편적으로 닫힘, 이 둘을 함께 만족함(즉 proper성)이라는 세 가지
서로 다른 개념에 이름을 붙여 두는 것이 유용하다. 국소 콤팩트 하우스도르프
공간들의 연속사상 $f : X \to Y$에 대해서는 ``proper''라는 말이 오래전부터
``$f^{-1}($콤팩트$) =$ 콤팩트''라는 개념을 가리키는 데 쓰였고, 이처럼 좋은
공간들에서는 이 개념이 보편적으로 닫힘과 동치이다. 실제로 일반적인
경우에도 사상이 닫혀 있다는 가정을 포함하면, 명료하게 ``준콤팩트''라는
말로 서술한 역상 조건이 보편적으로 닫힘과 동치임을 볼 것이다. 또한
\cite[Exercises 22-26 in Chapter II]{LangReal}를 보라. 다만 Lang은 부르바키와
마찬가지로 ``proper''를 ``보편적으로 닫힌''의 동의어로 사용함에 주의하라.

\begin{lemma}[튜브 보조정리]
\label{topology-lemma-tube}
$X$와 $Y$를 위상공간이라 하자. $A \subset X$와 $B \subset Y$를 준콤팩트
부분집합들이라 하자. $A \times B \subset W \subset X \times Y$이고 $W$가
$X \times Y$에서 열려 있다고 하자. 그러면 $U \times V \subset W$인
열린집합들 $A \subset U \subset X$와 $B \subset V \subset Y$가 존재한다.
\end{lemma}

\begin{proof}
모든 $a \in A$와 $b \in B$에 대하여 $X$의 열린집합 $U_{(a, b)}$와
$Y$의 열린집합 $V_{(a, b)}$가 존재하여
$(a, b) \in U_{(a, b)} \times V_{(a, b)} \subset W$이다. $b$를 고정하면
$A \subset U_{(a_1, b)} \cup \ldots \cup U_{(a_n, b)}$인 유한 개의
$a_1, \ldots, a_n$이 존재한다. 따라서
$$
A \times \{b\} \subset
(U_{(a_1, b)} \cup \ldots \cup U_{(a_n, b)}) \times
(V_{(a_1, b)} \cap \ldots \cap V_{(a_n, b)}) \subset W.
$$
그러므로 모든 $b \in B$에 대하여
$A \times \{b\} \subset U_b \times V_b \subset W$인 열린집합들
$U_b \subset X$와 $V_b \subset Y$가 존재한다. 위와 마찬가지로
$B \subset V_{b_1} \cup \ldots \cup V_{b_m}$인 유한 개의
$b_1, \ldots, b_m$이 존재한다. 이제
$A \times B \subset
(U_{b_1} \cap \ldots \cap U_{b_m}) \times
(V_{b_1} \cup \ldots \cup V_{b_m})$이므로 원하는 결과를 얻는다.
\end{proof}

\noindent
다음 정의의 용어법은 독자에게 익숙한 것과 조금 다를 수 있다.

\begin{definition}
\label{topology-definition-proper-map}
$f : X\to Y$를 위상공간 사이의 연속사상이라 하자.
\begin{enumerate}
\item 모든 닫힌 부분집합의 상이 닫혀 있으면 사상 $f$를 {\it 닫힌} 사상이라
한다.
\item 임의의 위상공간 $Z$에 대하여 사상 $Z \times X\to Z \times Y$가
닫혀 있으면 사상 $f$를 {\it 부르바키-proper}라고 한다.\footnote{이는
\cite{Bourbaki}에서 쓰는 용어법이다. 문헌에서는 때때로 이 성질을
``보편적으로 닫힌''이라고 부르기도 한다.}
\item 모든 준콤팩트 부분집합 $V \subset Y$의 역상 $f^{-1}(V)$가
준콤팩트이면 사상 $f$를 {\it 준-proper}라고 한다.
\item 임의의 연속사상 $g: Z \to Y$에 대하여 사상
$f': Z \times_Y X \to Z$가 닫혀 있으면 $f$를 {\it 보편적으로 닫힌}
사상이라 한다.
\item $f$가 분리되어 있고 보편적으로 닫혀 있으면 $f$를 {\it proper}라고
한다.
\end{enumerate}
\end{definition}

\noindent
다음 보조정리는 뒤에서 유용하다.

\begin{lemma}
\label{topology-lemma-characterize-quasi-compact}
\begin{reference}
\cite[I, p. 75, Lemme 1]{Bourbaki}와
\cite[I, p. 76, Corrolaire 1]{Bourbaki}의 결합이다.
\end{reference}
위상공간 $X$가 준콤팩트일 필요충분조건은 임의의 위상공간 $Z$에 대하여
사영사상 $Z \times X \to Z$가 닫혀 있는 것이다.
\end{lemma}

\begin{proof}
(아래의 비고도 보라.)
$X$가 준콤팩트가 아니면 유한 개의 $U_i$로는 $X$를 덮을 수 없는 열린 덮개
$X = \bigcup_{i \in I} U_i$가 존재한다. $I$의 멱집합 $\mathcal{P}(I)$에서
$I$와 $I$의 공집합이 아닌 모든 유한 부분집합으로 이루어진 부분집합을
$Z$라 하자. 다음 집합들을 위상의 기저로 삼아 $Z$에 위상을 정의한다.
\begin{enumerate}
\item $Z\setminus\{I\}$의 모든 부분집합.
\item $I$의 각 유한 부분집합 $K$에 대한 집합
% P01-E0018: the frozen set-builder formula has an unmatched closing parenthesis.
$U_K := \{J\subset I \mid J \in Z, \ K\subset J \})$.
\end{enumerate}
이들이 위상의 기저임을 확인하는 것은 독자에게 맡긴다. 다음 공식으로
정의되는 $Z \times X$의 부분집합을 생각하자.
$$
M = \{(J, x) \mid J \in Z, \ x \in \bigcap\nolimits_{i \in J} U_i^c)\}
$$
$(J, x) \not \in M$이면 어떤 $i \in J$에 대하여 $x \in U_i$이다. 따라서
$U_{\{i\}} \times U_i \subset Z \times X$는 $(J, x)$를 포함하고 $M$과
만나지 않는 열린 부분집합이다. 그러므로 $M$은 닫혀 있다. $M$의 $Z$로의
% P01-E1195: the frozen set difference below uses arithmetic minus.
사영은 닫혀 있지 않은 $Z-\{I\}$이다. 따라서 $Z \times X \to Z$는 닫힌
사상이 아니다.

\medskip\noindent
$X$가 준콤팩트라고 하자. $Z$를 위상공간이라 하고
$M \subset  Z \times X$를 닫힌집합이라 하자. $z \in Z$를
$\text{pr}_1(M)$에 속하지 않는 점이라 하자. Tube Lemma
\ref{topology-lemma-tube}에 의해 $U \times X$가 $M$의 여집합에 포함되는 열린집합
$U \subset Z$가 존재한다. 따라서 $\text{pr}_1(M)$은 닫혀 있다.
\end{proof}

\begin{remark}
\label{topology-remark-lemma-literature}
Lemma \ref{topology-lemma-characterize-quasi-compact}는
\cite[I, p. 75, Lemme 1]{Bourbaki}와
\cite[I, p. 76, Corollaire 1]{Bourbaki}의 결합이다.
\end{remark}

\begin{theorem}
\label{topology-theorem-characterize-proper}
\begin{reference}
\cite[I, p. 75, Theorem 1]{Bourbaki}에는
(2) $\Leftrightarrow$ (4)가 있다.
\cite[I, p. 77, Proposition 6]{Bourbaki}에는
(2) $\Rightarrow$ (1)이 있다.
\end{reference}
$f: X\to Y$를 위상공간 사이의 연속사상이라 하자. 다음 조건들은 서로
동치이다.
\begin{enumerate}
\item 사상 $f$는 준-proper이고 닫혀 있다.
\item 사상 $f$는 부르바키-proper이다.
\item 사상 $f$는 보편적으로 닫혀 있다.
\item 사상 $f$는 닫혀 있고 모든 $y\in Y$에 대하여 $f^{-1}(y)$가
준콤팩트이다.
\end{enumerate}
\end{theorem}

\begin{proof}
(아래의 비고도 보라.)
사상 $f$가 (1)을 만족하면 모든 한 점 집합이 준콤팩트이므로 자동으로 (4)를
만족한다.

\medskip\noindent
사상 $f$가 (4)를 만족한다고 하자. 이 사상이 보편적으로 닫혀 있음, 즉
(3)이 성립함을 증명하겠다. $g : Z \to Y$를 위상공간들의 연속사상이라 하고
다음 도표를 생각하자.
$$
\xymatrix{
Z \times_Y X \ar[r]_{g'} \ar[d]_{f'} & X \ar[d]^f \\
Z \ar[r]^g & Y
}
$$
증명 중에 $Z \times_Y X \to Z \times X$가 그 상 위의 위상동형이라는
사실을 사용할 것이다. 즉 $Z \times_Y X$를 유도위상을 준
$Z \times X$의 대응하는 부분집합과 동일시할 수 있다.
$f' : Z \times_Y X \to Z$의 상은
$\Im(f') = \{z : g(z) \in f(X)\}$이다. $f(X)$가 닫혀 있으므로
$\Im(f')$가 $Z$의 닫힌 부분공간임을 알 수 있다. 닫힌 부분집합
$P \subset Z \times_Y X$를 생각하자. $z \in Z$이고
$z \not \in f'(P)$라 하자. $z \not \in \Im(f')$이면
$Z \setminus \Im(f')$는 $f'(P)$를 피하는 열린 근방이다. $z$가
$\Im(f')$에 속하면
$(f')^{-1}\{z\} = \{z\} \times f^{-1}\{g(z)\}$이고,
$f^{-1}\{g(z)\}$는 가정에 의해 준콤팩트이다. $P$가
$Z \times_Y X$의 닫힌 부분집합이므로 $P = P' \cap Z \times_Y X$인
$Z \times X$의 닫힌집합 $P'$이 존재한다.
$(f')^{-1}\{z\}$는 $P^c = P'^c \cup (Z \times_Y X)^c$의 부분집합이고
$(f')^{-1}\{z\}$가 $(Z \times_Y X)^c$와 서로소이므로,
$(f')^{-1}\{z\}$는 $P'^c$에
포함된다. Tube Lemma \ref{topology-lemma-tube}를 다음 포함관계에 적용할 수 있다.
$(f')^{-1}\{z\} = \{z\} \times f^{-1}\{g(z)\}
\subset (P')^c \subset Z \times X$.
그러면 $(f')^{-1}\{z\}$를 포함하는 $V \times U$를 얻는다. 여기서 $U$와
$V$는 각각 $X$와 $Z$의 열린집합이고 $V \times U$는 $P'$과의 교집합이
공집합이다. 이제
% P01-E1195: the frozen set difference below also uses arithmetic minus.
$V \cap g^{-1}(Y-f(U^c))$는 $f$가 닫혀 있으므로 $Z$에서 열려 있고,
$z$를 포함하며 $P$의 상과의 교집합이 공집합이다. 따라서 $f'(P)$는 닫혀
있다. 다시 말해 사상 $f$는 보편적으로 닫혀 있다.

\medskip\noindent
(3) $\Rightarrow$ (2)는 자명하다. 실제로 임의의 위상공간 $Z$가 주어지면
사영사상 $g : Z \times Y \to Y$를 생각하자. 그러면 $f'$이 사상
$Z \times X \to Z \times Y$임, 다시 말해
$(Z \times Y) \times_Y X = Z \times X$임을 쉽게 알 수 있다. 이 동일시는
순전히 범주론적인 성질이며 위상공간 자체와는 아무 관계가 없다.

\medskip\noindent
$f$가 (2)를 만족한다고 하자. (1)을 만족함을 증명하겠다. $f$는 닫혀 있다고
가정된 사상 $\{pt\} \times X \to \{pt\} \times Y$와 $f$를 동일시할 수 있으므로
닫혀 있음에 유의하라. 임의의 준콤팩트 부분집합 $K \subset Y$를 택하자.
$Z$를 임의의 위상공간이라 하자. $Z \times X \to Z \times Y$가 닫혀
있으므로 사상 $Z \times f^{-1}(K) \to Z \times K$가 닫혀 있음을 알 수
있다. 실제로 $T$가 $Z \times f^{-1}(K)$에서 닫혀 있으면 어떤 닫힌집합
$T' \subset Z \times X$에 대하여
$T = Z \times f^{-1}(K) \cap T'$로 쓸 수 있다. $K$가 준콤팩트이므로
Lemma \ref{topology-lemma-characterize-quasi-compact}에 의해 $K \times Z\to Z$가
닫혀 있다. 따라서 합성 $Z \times f^{-1}(K)\to Z \times K \to Z$가 닫혀
있고, 다시 Lemma \ref{topology-lemma-characterize-quasi-compact}에 의해
$f^{-1}(K)$는 준콤팩트이어야 한다.
\end{proof}

\begin{remark}
\label{topology-remark-proof-literature}
관련 문헌을 몇 가지 제시한다. \cite[I, p. 75, Theorem 1]{Bourbaki}에는
(2) $\Leftrightarrow$ (4)가 있다. \cite[I, p. 77, Proposition 6]{Bourbaki}에는
(2) $\Rightarrow$ (1)이 있다. 물론 자명하게 (1) $\Rightarrow$ (4)도
성립한다. 따라서 (1), (2), (4)는 서로 동치이다. (3)과 (4)의 동치는
\cite[Chapter II, Exercise 25]{LangReal}이다.
\end{remark}

\begin{lemma}
\label{topology-lemma-closed-map}
\begin{slogan}
콤팩트 공간에서 하우스도르프 공간으로 가는 사상은 보편적으로 닫혀 있다.
\end{slogan}
$f : X \to Y$를 위상공간들의 연속사상이라 하자. $X$가 준콤팩트이고
$Y$가 하우스도르프이면 $f$는 보편적으로 닫혀 있다.
\end{lemma}

\begin{proof}
$Y$의 모든 점은 닫혀 있으므로 Lemma
\ref{topology-lemma-closed-in-quasi-compact}에 의해 모든 $y \in Y$에 대하여 $X$의
닫힌 부분집합 $f^{-1}(y)$가 준콤팩트임을 알 수 있다. 따라서 Theorem
\ref{topology-theorem-characterize-proper}에 의해 $f$가 닫혀 있음을 보이면 충분하다.
$E \subset X$가 닫혀 있으면 이는 준콤팩트이고(Lemma
\ref{topology-lemma-closed-in-quasi-compact}), 따라서 $f(E) \subset Y$도
준콤팩트이며(Lemma \ref{topology-lemma-image-quasi-compact}), 따라서 $f(E)$는
$Y$에서 닫혀 있다(Lemma \ref{topology-lemma-quasi-compact-in-Hausdorff}).
\end{proof}

\begin{lemma}
\label{topology-lemma-bijective-map}
$f : X \to Y$를 위상공간들의 연속사상이라 하자. $f$가 전단사이고 $X$가
준콤팩트이며 $Y$가 하우스도르프이면 $f$는 위상동형이다.
\end{lemma}

\begin{proof}
$f$가 닫혀 있음을 보이면 $f^{-1}$이 연속임이 따르므로 충분하다.
$T \subset X$가 닫혀 있으면 Lemma \ref{topology-lemma-closed-in-quasi-compact}에 의해
$T$는 준콤팩트이고, Lemma \ref{topology-lemma-image-quasi-compact}에 의해 $f(T)$는
준콤팩트이며, Lemma \ref{topology-lemma-quasi-compact-in-Hausdorff}에 의해 $f(T)$는
닫혀 있다.
\end{proof}

\section{Jacobson 공간}
\label{topology-section-space-jacobson}

\begin{definition}
\label{topology-definition-space-jacobson}
$X$를 위상공간이라 하자. $X_0$를 $X$의 닫힌점들의 집합이라 하자.
모든 닫힌 부분집합 $Z \subset X$가 $Z \cap X_0$의 폐포이면
$X$를 {\it Jacobson}이라 한다.
\end{definition}

\noindent
위상공간 $X$가 Jacobson일 필요충분조건은 $X$의 모든 비공집합 국소닫힌
부분집합이 $X$에서 닫힌 점을 가지는 것임에 유의하라.

\medskip\noindent
$X$를 Jacobson 공간이라 하고, $X_0$를 유도위상을 준 $X$의 닫힌점들의
집합이라 하자. 정의에 의해 사상 $X_0 \to X$가 $X_0$의 닫힌 부분집합과
$X$의 닫힌 부분집합 사이의 전단사를 유도함은 분명하다. 따라서 $X$의
많은 성질을 $X_0$가 물려받는다. 예를 들어 $X$와 $X_0$의 크룰 차원은 같다.

\begin{lemma}
\label{topology-lemma-jacobson-check-irreducible-closed}
$X$를 위상공간이라 하고, $X_0$를 $X$의 닫힌점들의 집합이라 하자.
모든 점 $x\in X$에 대하여 교집합 $X_0 \cap \overline{\{x\}}$가
$\overline{\{x\}}$에서 조밀하다고 하자. 그러면 $X$는 Jacobson이다.
\end{lemma}

\begin{proof}
$Z$를 $X$의 닫힌 부분집합, $U$를 $X$의 열린 부분집합이라 하고
$U\cap Z$가 비어 있지 않다고 하자. 그러면 $x\in U\cap Z$에 대하여
$\overline{\{x\}}\cap U$는 $Z\cap U$의 비공집합 부분집합이고,
가정에 의해 이는 필요한 대로 $X$에서 닫힌 점을 포함한다.
\end{proof}

\begin{lemma}
\label{topology-lemma-non-jacobson-Noetherian-characterize}
$X$를 준콤팩트 열린집합들로 이루어진 기저를 가지는 콜모고로프 위상공간이라
하자. $X$가 Jacobson이 아니면, $\{x\}$가 국소닫힌집합인 닫히지 않은 점
$x \in X$가 존재한다.
\end{lemma}

\begin{proof}
$X$가 Jacobson이 아니므로 $Z \cap U$가 비어 있지 않고 $X$에서 닫힌 점을
포함하지 않도록 하는 $X$의 닫힌집합 $Z$와 열린집합 $U$가 존재한다.
$X$가 준콤팩트 열린집합들로 이루어진 기저를 가지므로 $U$를 $Z\cap U$의
한 점의 준콤팩트 열린근방으로 바꿀 수 있고, 따라서 $U$가 준콤팩트
열린집합이라고 가정해도 된다. Lemma \ref{topology-lemma-quasi-compact-closed-point}에
의해 $Z \cap U$에서 닫힌 점 $x \in Z \cap U$가 존재한다. 따라서
$\{x\}$는 국소닫혀 있지만 $X$에서는 닫혀 있지 않다.
\end{proof}

\begin{lemma}
\label{topology-lemma-jacobson-local}
$X$를 위상공간이라 하고 $X = \bigcup U_i$를 열린 덮개라 하자.
$X$가 Jacobson일 필요충분조건은 각 $U_i$가 Jacobson인 것이다.
또한 이 경우 $X_0 = \bigcup U_{i, 0}$이다.
\end{lemma}

\begin{proof}
$X$를 위상공간이라 하자. $X_0$를 $X$의 닫힌점들의 집합이라 하고,
$U_{i, 0}$를 $U_i$의 닫힌점들의 집합이라 하자. 그러면
$X_0 \cap U_i \subset U_{i, 0}$이지만 일반적으로 등호가 성립하지는 않는다.

\medskip\noindent
먼저 각 $U_i$가 Jacobson이라고 하자. 이 경우
$X_0 \cap U_i = U_{i, 0}$이라고 주장한다. 실제로 $x \in U_{i, 0}$,
즉 $x$가 $U_i$에서 닫혀 있다고 하자. $\overline{\{x\}}$를 $X$에서의
폐포라 하고 $\overline{\{x\}} \cap U_j$를 생각하자. $x \not \in U_j$이면
$\overline{\{x\}} \cap U_j = \emptyset$이다. $x \in U_j$이면
$U_i \cap U_j \subset U_j$는 $x$를 포함하는 $U_j$의 열린 부분집합이다.
$T' = U_j \setminus U_i \cap U_j$이고 $T = \{x\} \amalg T'$라 하자.
$x$가 $U_i$에서 닫혀 있으므로 $x$는 $U_i \cap U_j$에서도 닫혀 있고,
따라서 $T$, $T'$은 $U_j$의 닫힌 부분집합이며 $T$는 $x$를 포함한다.
$U_j$가 Jacobson이므로 $U_j$의 닫힌점들은 $T$에서 조밀하다.
$T = \{x\} \amalg T'$이므로 이는 $x$가 $U_j$에서 닫혀 있을 때만 가능하다.
따라서 $\overline{\{x\}} \cap U_j = \{x\}$이다. 그러므로 원하는 대로
$\overline{\{x\}} = \{ x \}$이다.

\medskip\noindent
여전히 각 $U_i$가 Jacobson이라고 가정하고 $Z \subset X$를 닫힌
부분집합이라 하자. 이제 $X_0 \cap Z  \cap U_i
= U_{i, 0} \cap Z$가 $Z \cap U_i$에서 조밀함을 알았으므로
$X_0 \cap Z$가 $Z$에서 조밀함이 곧바로 따른다.

\medskip\noindent
반대로 $X$가 Jacobson이라고 하자. $Z \subset U_i$가 닫혀 있다고 하자.
$X_0 \cap \overline{Z}$는 $\overline{Z}$에서 조밀하다.
$\overline{Z} \setminus Z$가 닫혀 있으므로 $X_0 \cap Z$도 $Z$에서 조밀하다.
$X_0 \cap U_i \subset U_{i, 0}$이므로 $U_{i, 0} \cap Z$가 $Z$에서
조밀하다. 따라서 원하는 대로 $U_i$는 Jacobson이다.
\end{proof}

\begin{lemma}
\label{topology-lemma-jacobson-inherited}
$X$가 Jacobson이라고 하자. 다음 유형의 부분집합 $T \subset X$는
Jacobson이다.
\begin{enumerate}
\item 열린 부분공간.
\item 닫힌 부분공간.
\item 국소닫힌 부분공간.
\item 국소닫힌 부분공간들의 합집합.
\item 구성가능 집합.
\item $X$ 위에서 국소적으로 국소닫힌 부분집합들의 합집합인 임의의
부분집합 $T \subset X$.
\end{enumerate}
이 각각의 경우 $T$의 닫힌점들은 $X$에서 닫혀 있다.
\end{lemma}

\begin{proof}
$X_0$를 $X$의 닫힌점들의 집합이라 하자. 임의의 부분집합 $T \subset X$에
대하여 $(*)$는 다음 성질을 뜻한다고 하자.
\begin{itemize}
\item[(*)] $T$의 모든 비공집합 국소닫힌 부분집합은 $X$에서 닫힌 점을 가진다.
\end{itemize}
항상 $X_0 \cap T \subset T_0$임에 유의하라. 따라서 성질 $(*)$는 $T$가
Jacobson임을 함의한다. 또한 이는 $T$의 모든 닫힌점이 $X$에서 닫혀 있음을
분명히 함의한다.

\medskip\noindent
$T=\bigcup_i T_i$이고 $T_i$가 $X$에서 국소닫혀 있다고 하자.
$A\subset T$를 $T$의 비공집합 국소닫힌 부분집합이라 하자. 그러면
$A\cap T_i$가 비어 있지 않은 $T_i$가 존재하며, 이는 $T_i$에서 국소닫혀
있으므로 $X$에서도 국소닫혀 있다. $X$가 Jacobson이므로 $A$는 $X$에서
닫힌 점을 가진다.
\end{proof}

\begin{lemma}
\label{topology-lemma-finite-jacobson}
유한 Jacobson 공간은 이산공간이다.
유한개의 닫힌점들을 가지는 Jacobson 공간은 이산공간이다.
\end{lemma}

\begin{proof}
$X$가 유한하면 닫힌점들의 집합 $X_0 \subset X$는 유한하다.
$X_0$가 유한하고 $X$가 Jacobson이라고 하자. 그러면 Jacobson 성질에 의해
$\overline{X_0} = X$이다. 이제
$X_0 =\{x_1, \ldots, x_n\} = \bigcup_{i = 1, \ldots, n} \{x_i\}$는
닫힌집합들의 유한 합집합이므로 닫혀 있고, 따라서
$X = \overline{X_0} = X_0$이다. 모든 점은 닫혀 있으며 유한성에 의해
모든 점은 열려 있다.
\end{proof}

\begin{lemma}
\label{topology-lemma-jacobson-equivalent-locally-closed}
\begin{slogan}
Jacobson 공간에서는 닫힌점들이 위상에 관한 모든 것을 포착한다.
\end{slogan}
$X$를 Jacobson 위상공간이라 하자. $X_0$를 $X$의 닫힌점들의 집합이라 하자.
$E \mapsto E \cap X_0$로 주어지는, 전단사이고 포함관계를 보존하는 대응
$$
\{\text{finite unions loc. closed subsets of } X\}
\leftrightarrow
\{\text{finite unions loc. closed subsets of } X_0\}
$$
이 존재한다. 이 대응은 국소닫힌 부분집합, 열린 부분집합, 닫힌 부분집합의
각 모임을 보존한다.
\end{lemma}

\begin{proof}
대응 $E \mapsto E \cap X_0$가 단사임만 보이면 된다. 실제로 $E\neq E'$이면
일반성을 잃지 않고 $E\setminus E'$가 비어 있지 않다고 할 수 있으며,
이는 국소닫힌집합들의 유한 합집합이다(자세한 내용은 생략한다).
$X$가 Jacobson이므로
$(E \setminus E') \cap X_0 = E \cap X_0 \setminus E' \cap X_0$는 비어 있지 않다.
\end{proof}

\begin{lemma}
\label{topology-lemma-jacobson-equivalent-constructible}
$X$를 Jacobson 위상공간이라 하자. $X_0$를 $X$의 닫힌점들의 집합이라 하자.
$E \mapsto E \cap X_0$로 주어지는, 전단사이고 포함관계를 보존하는 대응
$$
\{\text{constructible subsets of } X\}
\leftrightarrow
\{\text{constructible subsets of } X_0\}
$$
이 존재한다. 이 대응은 레트로콤팩트 열린 부분집합들의 모임과 그들의
여집합들의 모임을 보존한다.
\end{lemma}

\begin{proof}
위의 Lemma \ref{topology-lemma-jacobson-equivalent-locally-closed}에 의해,
$U$가 $X$에서 열려 있을 때 $U\cap X_0$가 $X_0$에서 레트로콤팩트일
필요충분조건이 $U$가 $X$에서 레트로콤팩트인 것임을 보이면 된다.
이는 $U$가 $X$에서 열려 있을 때 $U\cap X_0$가 준콤팩트일 필요충분조건이
$U$가 준콤팩트인 것임을 보이면 따른다. Lemma
\ref{topology-lemma-jacobson-inherited}에 의해 $X$를 $U$로 바꾸고 $U = X$라고
가정해도 된다. 마지막으로 $X$의 열린집합들의 임의의 모임 $\mathcal{U}$가
$X$를 덮을 필요충분조건은 $X_0$를 덮는 것임에 유의하라. 실제로
$X\setminus \bigcup \mathcal{U}$가 비어 있지 않다면, 이 닫힌집합에 $X$의
Jacobson 성질을 적용하여 그 안에서 $X_0$의 한 점을 찾을 수 있다.
\end{proof}

\section{특수화}
\label{topology-section-specialization}

\begin{definition}
\label{topology-definition-specialization}
$X$를 위상공간이라 하자.
\begin{enumerate}
\item $x, x' \in X$일 때 $x \in \overline{\{x'\}}$이면, $x$는 $x'$의
{\it 특수화}라고 하거나 $x'$은 $x$의 {\it 일반화}라고 한다.
표기는 $x' \leadsto x$이다.
\item 부분집합 $T \subset X$에 대하여, 모든 $x' \in T$와 모든 특수화
$x' \leadsto x$에 대해 $x \in T$이면 그 부분집합은 {\it 특수화에 대해 안정}이라고 한다.
\item 부분집합 $T \subset X$에 대하여, 모든 $x \in T$와 모든 일반화
$x' \leadsto x$에 대해 $x' \in T$이면 그 부분집합은 {\it 일반화에 대해 안정}이라고 한다.
\end{enumerate}
\end{definition}

\begin{lemma}
\label{topology-lemma-open-closed-specialization}
$X$를 위상공간이라 하자.
\begin{enumerate}
\item $X$의 모든 닫힌 부분집합은 특수화에 대해 안정이다.
\item $X$의 모든 열린 부분집합은 일반화에 대해 안정이다.
\item 부분집합 $T \subset X$가 특수화에 대해 안정일 필요충분조건은
여집합 $T^c$가 일반화에 대해 안정인 것이다.
\end{enumerate}
\end{lemma}

\begin{proof}
$F$를 $X$의 닫힌 부분집합이라 하자. $y\in F$이면 $\{y\} \subset F$이고,
$F$가 닫혀 있으므로 $\overline{\{y\}} \subset \overline{F} = F$이다.
따라서 $y$의 모든 특수화 $x$에 대하여 $x\in F$이다.

\medskip\noindent
$x, y\in X$이고 $x\in \overline{\{y\}}$라고 하며, $T$를 $X$의
부분집합이라 하자. $T$가 특수화에 대해 안정이라는 것은 $y\in T$이면
$x\in T$라는 뜻이고, 반대로 $T$가 일반화에 대해 안정이라는 것은
$x\in T$이면 $y\in T$라는 뜻이다. 따라서 대우를 사용하면 (3)이 증명된다.

\medskip\noindent
두 번째 성질은 여집합을 생각하면 (1)과 (3)에서 따른다.
\end{proof}

\begin{lemma}
\label{topology-lemma-stable-specialization}
$T \subset X$를 위상공간 $X$의 부분집합이라 하자. 다음은 서로 동치이다.
\begin{enumerate}
\item $T$는 특수화에 대해 안정이다.
\item $T$는 $X$의 닫힌 부분집합들의 (유향) 합집합이다.
\end{enumerate}
\end{lemma}

\begin{proof}
$T$가 특수화에 대해 안정이라고 하자. 그러면 모든 $y\in T$에 대하여
$\overline{\{y\}} \subset T$이다. 따라서
$T = \bigcup_{y\in T} \overline{\{y\}}$이며, 이는 $X$의 닫힌 부분집합들의
합집합이다. 반대로 $T = \bigcup_{i\in I}F_i$이고 $F_i$들이 $X$의 닫힌
부분집합들이라고 하자. $y\in T$이면 $y\in F_i$인 $i\in I$가 존재한다.
$F_i$가 닫혀 있으므로 $\overline{\{y\}} \subset F_i \subset T$이고,
이는 $T$가 특수화에 대해 안정임을 증명한다.
\end{proof}

\begin{definition}
\label{topology-definition-lift-specializations}
$f : X \to Y$를 위상공간들의 연속사상이라 하자.
\begin{enumerate}
\item $Y$에서 $y' \leadsto y$가 주어지고 $f(x') = y'$인 임의의
$x'\in X$가 주어질 때, $f(x) = y$를 만족하는 $X$에서의 $x'$의 특수화
$x' \leadsto x$가 존재하면 {\it 특수화가 $f$를 따라 올려진다}고 하거나
$f$를 {\it 특수화 사상}이라고 한다.
\item $Y$에서 $y' \leadsto y$가 주어지고 $f(x) = y$인 임의의 $x\in X$가
주어질 때, $f(x') = y'$을 만족하는 $X$에서의 $x$의 일반화
$x' \leadsto x$가 존재하면 {\it 일반화가 $f$를 따라 올려진다}고 하거나
$f$를 {\it 일반화 사상}이라고 한다.
\end{enumerate}
\end{definition}

\begin{lemma}
\label{topology-lemma-lift-specialization-composition}
$f : X \to Y$와 $g : Y \to Z$가 위상공간들의 연속사상이라고 하자.
특수화가 $f$와 $g$를 따라 모두 올려지면 특수화는 $g \circ f$를 따라
올려진다. ``일반화가 따라 올려진다''에 대해서도 마찬가지이다.
\end{lemma}

\begin{proof}
$z'\leadsto z$를 $Z$에서의 특수화라 하고, $g\circ f (x') = z'$인
$x' \in X$가 주어졌다고 하자. 특수화가 $g$를 따라 올려지므로
$g(y) = z$인 $Y$에서의 $f(x')$의 특수화 $f(x') \leadsto y$가 존재한다.
마찬가지로 특수화가 $f$를 따라 올려지므로 $f(x) = y$인 $X$에서의
$x'$의 특수화 $x' \leadsto x$가 존재한다. 이는 $g\circ f(x) = z$인
$X$에서의 $x'$의 특수화 $x' \leadsto x$를 준다. 다시 말해 특수화는
$g\circ f$를 따라 올려진다.
\end{proof}

\begin{lemma}
\label{topology-lemma-lift-specializations-images}
$f : X \to Y$를 위상공간들의 연속사상이라 하자.
\begin{enumerate}
\item 특수화가 $f$를 따라 올려지고 $T \subset X$가 특수화에 대해
안정이면 $f(T) \subset Y$도 특수화에 대해 안정이다.
\item 일반화가 $f$를 따라 올려지고 $T \subset X$가 일반화에 대해
안정이면 $f(T) \subset Y$도 일반화에 대해 안정이다.
\end{enumerate}
\end{lemma}

\begin{proof}
$y'\in f(T)$인 $Y$에서의 특수화 $y' \leadsto y$가 주어지고,
$f(x') = y'$인 $x'\in T$를 택하자. 특수화가 $f$를 따라 올려지므로
$f(x) = y$인 $X$에서의 $x'$의 특수화 $x'\leadsto x$가 존재한다.
그런데 $T$가 특수화에 대해 안정이므로 $x\in T$이고 따라서
$y \in f(T)$이다. 그러므로 $f(T)$는 특수화에 대해 안정이다.

\medskip\noindent
(2)의 증명은 일반화가 $f$를 따라 올려진다는 것을 사용하면 동일하다.
\end{proof}

\begin{lemma}
\label{topology-lemma-closed-open-map-specialization}
$f : X \to Y$를 위상공간들의 연속사상이라 하자.
\begin{enumerate}
\item $f$가 닫힌사상이면 특수화가 $f$를 따라 올려진다.
\item $f$가 열린사상이고 $X$가 뇌터 위상공간이며 $X$의 각 기약 닫힌
부분집합이 일반점을 가지고 $Y$가 콜모고로프이면 일반화가 $f$를 따라
올려진다.
\end{enumerate}
\end{lemma}

\begin{proof}
$f$가 닫힌사상이라고 하자. $Y$에서 $y' \leadsto y$가 주어지고
$f(x') = y'$인 임의의 $x'\in X$가 주어졌다고 하자. $X$의 닫힌 부분집합
$T = \overline{\{x'\}}$를 생각하자. 그러면 $f(T) \subset Y$는 닫힌
부분집합이고 $y' \in f(T)$이다. 따라서 $y \in f(T)$이기도 하다.
그러므로 $x \in T$인 점에 대하여 $y = f(x)$이다. 즉 $x' \leadsto x$이다.

\medskip\noindent
$f$가 열린사상이고 $X$가 뇌터이며 $X$의 모든 기약 닫힌 부분집합이
일반점을 가지고 $Y$가 콜모고로프라고 하자. $Y$에서 $y' \leadsto y$가
주어지고 $f(x) = y$인 임의의 $x \in X$가 주어졌다고 하자.
$T = f^{-1}(\{y'\}) \subset X$를 생각하자. $x$의 열린근방
$x \in U \subset X$를 택하자. 그러면 $f(U) \subset Y$는 열려 있고
$y \in f(U)$이다. 따라서 $y' \in f(U)$이기도 하다. 다시 말해
$T \cap U \not = \emptyset$이다. 이는 $x \in \overline{T}$임을 증명한다.
$X$가 뇌터이므로 $T$는 뇌터이다(Lemma \ref{topology-lemma-Noetherian}).
따라서 $T = T_1 \cup \ldots \cup T_n$인 기약성분 분해가 존재한다.
이에 따라 $\overline{T} = \overline{T_1} \cup \ldots \cup \overline{T_n}$이다.
위에서 본 바에 의해 어떤 $i$에 대하여 $x \in \overline{T_i}$이다.
가정에 의해 일반점 $x' \in \overline{T_i}$가 존재하며 $x' \leadsto x$임을
알 수 있다. $x' \in \overline{T}$이므로 $f(x') \in \overline{\{y'\}}$임을
알 수 있다. 다음에 유의하라.
$f(\overline{T_i}) = f(\overline{\{x'\}}) \subset \overline{\{f(x')\}}$.
$f(x') \not = y'$이면, $Y$가 콜모고로프이므로 $f(x')$은 기약 닫힌 부분집합
$\overline{\{y'\}}$의 일반점이 아니고 포함관계
$\overline{\{f(x')\}} \subset \overline{\{y'\}}$는 엄밀하다.
즉 $y' \not \in f(\overline{T_i})$이다. 이는 $f(T_i) = \{y'\}$라는 사실과
모순이다. 따라서 $f(x') = y'$이고 증명이 끝난다.
\end{proof}

\begin{lemma}
\label{topology-lemma-quotient-kolmogorov}
$s, t : R \to U$와 $\pi : U \to X$가 다음을 만족하는 위상공간들의
연속사상이라고 하자.
\begin{enumerate}
\item $\pi$는 열린사상이다.
\item $U$는 소버이다.
\item $s, t$의 섬유들은 유한하다.
\item 일반화가 $s, t$를 따라 올려진다.
\item $(t, s)(R) \subset U \times U$는 $U$ 위의 동치관계이고, $X$는
(집합으로서) 이 동치관계에 의한 $U$의 몫이다.
\end{enumerate}
그러면 $X$는 콜모고로프이다.
\end{lemma}

\begin{proof}
성질 (3)과 (5)에 의해 점 $x$는 이 동치관계의 유한 동치류
$\{u_1, \ldots, u_n\} \subset U$에 대응한다. $x' \in X$를 동치류
$\{u'_1, \ldots, u'_m\} \subset U$에 대응하는 두 번째 점이라고 하자.
어떤 $i, j$에 대하여 $u_i \leadsto u'_j$라고 하자. 그러면
$s(r') = u'_j$인 임의의 $r' \in R$에 대하여 (4)에 의해
$s(r) = u_i$인 $r \leadsto r'$을 찾을 수 있다. 따라서
$t(r) \leadsto t(r')$이다.
$\{u'_1, \ldots, u'_m\} = t(s^{-1}(\{u'_j\}))$이므로
$\{u'_1, \ldots, u'_m\}$의 모든 원소가 $\{u_1, \ldots, u_n\}$의 한 원소의
특수화임을 알 수 있다. 따라서
$\overline{\{u_1\}} \cup \ldots \cup \overline{\{u_n\}}$는 동치류들의
합집합이고, 따라서 어떤 부분집합 $Z \subset X$에 대하여
$\pi^{-1}(Z)$ 꼴이다. (1)에 의해 $Z$가 $X$에서 닫혀 있음을 알 수 있고,
각 $i$에 대해 $\pi(\overline{\{u_i\}}) \subset \overline{\{x\}}$이므로 실제로
$Z = \overline{\{x\}}$이다. 다시 말해, $x \leadsto x'$일 필요충분조건은
$U$에서 $x$의 어떤 올림이 $U$에서 $x'$의 어떤 올림으로 특수화되는 것이고,
이는 다시 $U$에서 $x'$의 모든 올림이 $U$에서 $x$의 어떤 올림의
특수화인 것과 동치이다.

\medskip\noindent
$x \leadsto x'$과 $x' \leadsto x$가 모두 성립한다고 하자. 위와 같이 $x$가
$\{u_1, \ldots, u_n\}$에 대응하고 $x'$이 $\{u'_1, \ldots, u'_m\}$에
대응한다고 하자. 그러면 앞 문단의 결과에 의해 다음 수열을 찾을 수 있다.
$$
\ldots \leadsto u'_{j_3} \leadsto u_{i_3} \leadsto u'_{j_2} \leadsto
u_{i_2} \leadsto u'_{j_1} \leadsto u_{i_1}
$$
이 수열은 반복되어야 하므로, (2)에 의해
$\{u_1, \ldots, u_n\} = \{u'_1, \ldots, u'_m\}$, 즉 $x = x'$이라고
결론 내린다. 따라서 $X$는 콜모고로프이다.
\end{proof}

\begin{lemma}
\label{topology-lemma-dimension-specializations-lift}
$f : X \to Y$를 위상공간들의 사상이라 하자. $Y$가 소버 위상공간이고
$f$가 전사라고 하자. 특수화나 일반화 중 하나가 $f$를 따라 올려지면
$\dim(X) \geq \dim(Y)$이다.
\end{lemma}

\begin{proof}
특수화가 $f$를 따라 올려진다고 하자.
% P01-E0019: the frozen chain below omits an inclusion after the ellipsis.
$Z_0 \subset Z_1 \subset \ldots Z_e \subset Y$를
% P01-E0020/P01-E1206: frozen prose says X although the displayed chain lies in Y.
$X$의 기약 닫힌 부분집합들의 사슬이라 하자. $\xi_e \in X$를 $Z_e$의
일반점으로 사상되는 점이라 하자. 가정에 의해 $X$에서의 특수화
$\xi_e \leadsto \xi_{e - 1}$이 존재하며, $\xi_{e - 1}$은 $Z_{e - 1}$의
일반점으로 사상된다.
이와 같이 계속하면 $\xi_i$가 $Z_i$의 일반점으로 사상되는 특수화들의 수열
$$
\xi_e \leadsto \xi_{e - 1} \leadsto \ldots \leadsto \xi_0
$$
을 얻는다. 이는 기약 닫힌 부분집합들의 수열
$$
\overline{\{\xi_0\}} \subset
\overline{\{\xi_1\}} \subset \ldots
\overline{\{\xi_e\}}
$$
이 $X$에서 길이 $e$인 사슬임을 분명히 함의한다.
일반화가 $f$를 따라 올려지는 경우도 비슷하다.
\end{proof}

\begin{lemma}
\label{topology-lemma-characterize-closed-Noetherian}
$X$를 뇌터 소버 위상공간이라 하자. $E \subset X$를 $X$의 부분집합이라 하자.
\begin{enumerate}
\item $E$가 구성가능하고 특수화에 대해 안정이면 $E$는 닫혀 있다.
\item $E$가 구성가능하고 일반화에 대해 안정이면 $E$는 열려 있다.
\end{enumerate}
\end{lemma}

\begin{proof}
$E$가 구성가능하고 일반화에 대해 안정이라고 하자. $Y \subset X$를 일반점
$\xi \in Y$를 가지는 기약 닫힌 부분집합이라 하자. $E \cap Y$가 비어 있지
않으면 일반화에 대한 안정성에 의해 이는 $\xi$를 포함하고, 따라서 $Y$에서
조밀하며, 따라서 $Y$의 비공집합 열린집합을 포함한다. Lemma
\ref{topology-lemma-characterize-constructible-Noetherian}를 보라. 그러므로 Lemma
\ref{topology-lemma-characterize-open-Noetherian}에 의해 $E$는 열려 있다. 이로써
(2)가 증명된다. (1)을 증명하려면 (2)를 $X$에서의 $E$의 여집합에 적용하라.
\end{proof}

\section{차원 함수}
\label{topology-section-dimension-function}

\noindent
생각하는 공간이 소버가 아니라면 차원 함수를 생각하는 것은 거의 의미가 없다
(Definition \ref{topology-definition-generic-point}). 따라서 아래 정의는 $X$에 대응하는
소버 위상공간을 생각함으로써 개선할 수 있다. 스킴의 바탕 위상공간은
소버이므로 여기서는 이 개선을 다루지 않는다.

\begin{definition}
\label{topology-definition-dimension-function}
$X$를 위상공간이라 하자.
\begin{enumerate}
\item $x, y \in X$이고 $x \not = y$라고 하자. $x \leadsto y$, 즉 $y$가
$x$의 특수화라고 하자. $x \leadsto z$이고 $z \leadsto y$인
$z \in X \setminus \{x, y\}$가 존재하지 않으면, $y$를 $x$의
{\it 직접 특수화}라고 한다.
\item 다음을 만족하는 사상 $\delta : X \to \mathbf{Z}$를
{\it 차원 함수}\footnote{이는 아마 비표준적인 표기일 것이다. 이 개념은 보통
(국소) 뇌터 스킴에 대해서만 도입되며, 이 경우 조건 (a)는 (b)에서 따른다.}라고
한다.
\begin{enumerate}
\item $x \leadsto y$이고 $x \not = y$일 때마다
$\delta(x) > \delta(y)$이다.
\item $X$에서의 모든 직접 특수화 $x \leadsto y$에 대하여
$\delta(x) = \delta(y) + 1$이다.
\end{enumerate}
\end{enumerate}
\end{definition}

\noindent
$\delta$가 차원 함수이면, 임의의 $t \in \mathbf{Z}$에 대하여
$\delta + t$도 차원 함수임이 분명하다. 다음은 재미있는 보조정리이다.

\begin{lemma}
\label{topology-lemma-dimension-function-catenary}
$X$를 위상공간이라 하자. $X$가 소버이고 차원 함수를 가지면 $X$는
카테너리이다. 또한 임의의 $x \leadsto y$에 대하여
$$
\delta(x) - \delta(y) =
\text{codim}\left(\overline{\{y\}}, \ \overline{\{x\}}\right).
$$
\end{lemma}

\begin{proof}
% P01-E1207: the frozen directions/signs below conflict with the definition and statement.
$Y \subset Y' \subset X$가 기약 닫힌 부분집합들이라고 하자.
$\xi \in Y$, $\xi' \in Y'$를 그 일반점들이라 하자. 그러면 정의로부터
$\text{codim}(Y, Y') \leq \delta(\xi) - \delta(\xi') < \infty$임을 곧바로
알 수 있다. 사실 첫 번째 부등식은 등식이다. 실제로
$$
Y = Y_0 \subset Y_1 \subset \ldots \subset Y_e = Y'
$$
가 기약 닫힌 부분집합들의 임의의 극대 사슬이라고 하자.
$\xi_i \in Y_i$를 일반점이라 하자. 그러면
$\xi_i \leadsto \xi_{i + 1}$이 직접 특수화임을 알 수 있다.
따라서 원하는 대로 $e = \delta(\xi) - \delta(\xi')$임을 알 수 있다.
이는 보조정리의 마지막 명제도 증명한다.
\end{proof}

\begin{lemma}
\label{topology-lemma-dimension-function-unique}
$X$를 위상공간이라 하자. $\delta$, $\delta'$을 $X$ 위의 두 차원 함수라 하자.
$X$가 국소 뇌터이고 소버이면 $\delta - \delta'$은 $X$ 위에서 국소상수이다.
\end{lemma}

\begin{proof}
$x \in X$를 한 점이라 하자. $\delta - \delta'$이 $x$의 한 근방에서
상수임을 보이겠다. $X$를 $X$에서 뇌터인 $x$의 열린근방으로 바꾸어도 된다.
따라서 $X$가 뇌터이고 소버라고 가정해도 된다. $Z_1, \ldots, Z_r$를 $x$를
지나는 $X$의 기약성분들이라 하자. ($X$가 뇌터이므로 이들은 유한개이다.
Lemma \ref{topology-lemma-Noetherian}을 보라.) $\xi_i \in Z_i$를 일반점이라 하자.
% P01-E1208: the frozen parenthetical says this finite union need not be closed.
$Z_1 \cup \ldots \cup Z_r$는 $X$에서 $x$의 근방임에 유의하라(반드시
닫힌 것은 아니다). $\delta - \delta'$이 $Z_1 \cup \ldots \cup Z_r$에서
상수라고 주장한다. 실제로 $y \in Z_i$이면
$$
\delta(x) - \delta(y) = \delta(x) - \delta(\xi_i) + \delta(\xi_i) - \delta(y)
= - \text{codim}(\overline{\{x\}}, Z_i)
+ \text{codim}(\overline{\{y\}}, Z_i)
$$
이다. 이는 Lemma \ref{topology-lemma-dimension-function-catenary}에서 따른다.
$\delta'$에 대해서도 마찬가지이다. 따라서 결과가 따른다.
\end{proof}

\begin{lemma}
\label{topology-lemma-locally-dimension-function}
$X$가 국소 뇌터이고 소버이며 카테너리라고 하자. 그러면 모든 점은
차원 함수를 가지는 열린근방 $U \subset X$를 가진다.
\end{lemma}

\begin{proof}
카테너리 공간의 열린 부분공간이 카테너리라는 것(Lemma
\ref{topology-lemma-catenary})과 뇌터 위상공간이 유한개의 기약성분을 가진다는 것
(Lemma \ref{topology-lemma-Noetherian})을 반복해서 사용하겠다. Lemma
\ref{topology-lemma-dimension-function-unique}의 증명에서 이러한 함수를 구성하는
방법을 보았다. 먼저 $X$를 $x$의 뇌터 열린근방으로 바꾼다. 다음으로
$Z_1, \ldots, Z_r \subset X$를 $X$의 기약성분들이라 하자.
% P01-E1209: the frozen decomposition below leaves its union unindexed.
$$
Z_i \cap Z_j = \bigcup Z_{ijk}
$$
를 기약성분 분해라 하자. $X$를
$$
X \setminus \left(
\bigcup\nolimits_{x \not \in Z_i} Z_i
\cup
\bigcup\nolimits_{x \not \in Z_{ijk}} Z_{ijk}
\right)
$$
로 바꾸면, 모든 $i$에 대해 $x \in Z_i$이고 모든 $i, j, k$에 대해
$x \in Z_{ijk}$이라고 가정해도 된다. $y \in X$에 대하여 $y \in Z_i$인
임의의 $i$를 택하고
$$
\delta(y) = - \text{codim}(\overline{\{y\}}, Z_i)
+ \text{codim}(\overline{\{x\}}, Z_i).
$$
로 둔다. 이것이 차원 함수라고 주장한다. 먼저 이것이 잘 정의됨, 즉 $i$의
선택과 무관함을 보이자. 실제로 어떤 $i, j, k$에 대하여
$y \in Z_{ijk}$라고 하자. 그러면 (Lemma
\ref{topology-lemma-catenary-in-codimension}을 사용하여) 다음을 얻는다.
\begin{align*}
\delta(y) & =
- \text{codim}(\overline{\{y\}}, Z_i)
+ \text{codim}(\overline{\{x\}}, Z_i) \\
& =
- \text{codim}(\overline{\{y\}}, Z_{ijk})
- \text{codim}(Z_{ijk}, Z_i)
+ \text{codim}(\overline{\{x\}}, Z_{ijk})
+ \text{codim}(Z_{ijk}, Z_i) \\
& =
- \text{codim}(\overline{\{y\}}, Z_{ijk})
+ \text{codim}(\overline{\{x\}}, Z_{ijk})
\end{align*}
이고, 이는 $i$와 $j$에 대해 대칭이다. 이것이 차원 함수라는 증명은 생략한다.
\end{proof}

\begin{remark}
\label{topology-remark-obstruction-to-dimension-function}
Lemmas \ref{topology-lemma-dimension-function-unique}와
\ref{topology-lemma-locally-dimension-function}을 결합하면, 카테너리이고 국소 뇌터이며
소버인 위상공간에서 차원 함수를 가지는 것에 대한 장애물이
$H^1(X, \mathbf{Z})$의 한 원소임을 알 수 있다.
\end{remark}

\section{어디에서도 조밀하지 않은 집합}
\label{topology-section-nowhere-dense}

\begin{definition}
\label{topology-definition-nowhere-dense}
$X$를 위상공간이라 하자.
\begin{enumerate}
\item 부분집합 $T \subset X$가 주어졌을 때, $T$의 {\it 내부}는 $T$에
포함되는 $X$의 가장 큰 열린 부분집합이다.
\item 부분집합 $T \subset X$의 폐포의 내부가 비어 있으면 $T$를
{\it 어디에서도 조밀하지 않다}고 한다.
\end{enumerate}
\end{definition}

\begin{lemma}
\label{topology-lemma-nowhere-dense}
$X$를 위상공간이라 하자. 유한개의 어디에서도 조밀하지 않은 집합들의
합집합은 어디에서도 조밀하지 않은 집합이다.
\end{lemma}

\begin{proof}
일반적인 결과는 귀납법으로 따르므로, 어디에서도 조밀하지 않은 두 집합에
대해 보이면 충분하다. $A,B \subset X$를 어디에서도 조밀하지 않은
부분집합들이라 하자. 다음이 성립한다.
$\overline{A \cup B} = \overline{A} \cup \overline{B}$.
따라서 $U \subset \overline{A \cup B}$가 $X$의 열린 부분집합이면,
% P01-E0021/P01-E1210: the frozen set differences omit grouping parentheses.
$U \setminus U \cap \overline{B}$는 $U$와 $X$의 열린 부분집합이고
$\overline{A}$에 포함되므로 비어 있다. 마찬가지로
$U \setminus U \cap \overline{A}$도 비어 있다. 따라서 원하는 대로
$U = \emptyset$이다.
\end{proof}

\begin{lemma}
\label{topology-lemma-image-nowhere-dense-open}
$X$를 위상공간이라 하자. $U \subset X$를 열린 부분집합이라 하고,
$T \subset U$를 부분집합이라 하자. $T$가 $U$에서 어디에서도 조밀하지
않으면 $T$는 $X$에서도 어디에서도 조밀하지 않다.
\end{lemma}

\begin{proof}
$T$가 $U$에서 어디에서도 조밀하지 않다고 하자. $x \in X$가 $X$에서의
$T$의 폐포 $\overline{T}$의 내부점이라고 하자. $V \subset X$가 $X$에서
열려 있고 $x \in V \subset \overline{T}$라고 하자. $\overline{T} \cap U$는
$U$에서의 $T$의 폐포임에 유의하라. 따라서 $\overline{T} \cap U$의 내부가
비어 있다는 사실은 $V \cap U = \emptyset$을 함의한다. 그러므로 $x$는
$U$의 폐포에 속할 수 없고, 하물며 $T$의 폐포에도 속할 수 없다. 이는 모순이다.
\end{proof}

\begin{lemma}
\label{topology-lemma-nowhere-dense-local}
$X$를 위상공간이라 하자. $X = \bigcup U_i$를 열린 덮개라 하고,
$T \subset X$를 부분집합이라 하자. 모든 $i$에 대하여 $T \cap U_i$가
$U_i$에서 어디에서도 조밀하지 않으면 $T$는 $X$에서 어디에서도
조밀하지 않다.
\end{lemma}

\begin{proof}
$\overline{T}_i$를 $U_i$에서의 $T \cap U_i$의 폐포라 하자.
$\overline{T} \cap U_i = \overline{T}_i$이다. 내부를 취하는 연산은
열린집합과의 교집합과 교환하므로
$$
(\text{interior of }\overline{T}) \cap U_i =
\text{interior of }(\overline{T} \cap U_i) =
\text{interior in }U_i\text{ of }\overline{T}_i
$$
이다. 가정에 의해 이들 중 마지막 것은 비어 있다. 따라서 $T$는 $X$에서
어디에서도 조밀하지 않다.
\end{proof}

\begin{lemma}
\label{topology-lemma-closed-image-nowhere-dense}
$f : X \to Y$를 위상공간들의 연속사상이라 하고 $T \subset X$를
부분집합이라 하자. $f$가 $X$에서 $Y$의 닫힌 부분집합 위로의 위상동형이고
$T$가 $X$에서 어디에서도 조밀하지 않으면 $f(T)$도 $Y$에서 어디에서도
조밀하지 않다.
\end{lemma}

\begin{proof}
$f(X)$가 $Y$에서 닫혀 있고 $f$가 $X$에서 $f(X)$ 위로의 위상동형이므로,
$Y$에서의 $f(T)$의 폐포는 $f(\overline{T})$와 같다. 따라서
$V \subset Y$가 열려 있고 $f(T)$의 폐포에 포함되면,
$U = f^{-1}(V)$는 열려 있고 $\overline{T}$에 포함된다. 그러므로
$U = \emptyset$이고, 이는 다시 원하는 대로 $V = \emptyset$임을 보인다.
\end{proof}

\begin{lemma}
\label{topology-lemma-open-inverse-image-closed-nowhere-dense}
$f : X \to Y$를 위상공간들의 연속사상이라 하고 $T \subset Y$를
부분집합이라 하자. $f$가 열린사상이고 $T$가 $Y$의 닫힌, 어디에서도
조밀하지 않은 부분집합이면 $f^{-1}(T)$도 $X$의 닫힌, 어디에서도
조밀하지 않은 부분집합이다. $f$가 전사인 열린사상이면, $T$가 닫혀 있고
어디에서도 조밀하지 않을 필요충분조건은 $f^{-1}(T)$가 닫혀 있고
어디에서도 조밀하지 않은 것이다.
\end{lemma}

\begin{proof}
생략한다. (힌트: 첫 번째 경우에는 $f^{-1}(T)$의 내부가 $T$의 내부로
사상되고, 두 번째 경우에는 $f^{-1}(T)$의 내부가 $T$의 내부 위로 사상된다.)
\end{proof}

\begin{lemma}
\label{topology-lemma-dense-in-constructible}
$X$를 위상공간이라 하자. $E \subset X$를 국소닫힌 부분집합들의 유한
합집합이라 하자(예를 들어 $E$가 구성가능인 경우). 그러면 $E$는 자신의
폐포의 조밀 열린 부분집합을 포함한다.
\end{lemma}

\begin{proof}
$X$를 $E$의 폐포로 바꾼 뒤에는 $E$가 $X$에서 조밀하다고 가정해도 되며,
$E$가 조밀 열린집합을 포함함을 보여야 한다. 그러면 $E$의 여집합 $T$는
국소닫힌 부분집합들의 유한 합집합이고, 이를
$T = T_1 \cup \ldots \cup T_n$이라 하자. 또한 $T$는 $X$에서 어디에서도
조밀하지 않다. $T_i$는 그 폐포 $Z_i$에서 열린 조밀 부분집합이므로
$Z_i$도 $X$에서 어디에서도 조밀하지 않다. 그러면 $T$의 폐포
$Z = Z_1 \cup \ldots \cup Z_n$은 Lemma \ref{topology-lemma-nowhere-dense}에 의해
$X$에서 어디에서도 조밀하지 않다. 따라서 $E$는 조밀 열린집합
$X \setminus Z$를 포함한다.
\end{proof}







\section{프로유한 공간}
\label{topology-section-profinite}

\noindent
정의는 다음과 같다.

\begin{definition}
\label{topology-definition-profinite}
위상공간이 유한 이산공간들로 이루어진 도표의 극한과 위상동형이면
이를 {\it 프로유한} 공간이라 한다.
\end{definition}

\noindent
이것은 프로유한 공간을 특징짓는 가장 편리한 방법은 아니다.

\begin{lemma}
\label{topology-lemma-profinite}
$X$를 위상공간이라 하자. 다음은 동치이다.
\begin{enumerate}
\item $X$는 프로유한 공간이다.
\item $X$는 하우스도르프이고 준콤팩트하며 완전불연결이다.
\end{enumerate}
이들이 성립하면 $X$는 유한 이산공간들의 쌍대여과 극한이다.
\end{lemma}

\begin{proof}
(1)을 가정하자. $X = \lim X_i$가 되도록 유한 이산공간들의 도표
$i \mapsto X_i$를 택한다. 각 $X_i$가 하우스도르프이고 준콤팩트하므로,
Lemma \ref{topology-lemma-inverse-limit-quasi-compact}에 의해 $X$가 준콤팩트임을
알 수 있다. 서로 다른 점 $x, x' \in X$가 있으면 $x$와 $x'$는 어떤
$X_i$에서 서로 다른 점으로 사상된다. 따라서 $x$와 $x'$는 서로소인
열린근방들을 가지므로 $X$는 하우스도르프이다. 정확히 같은 방법으로
$X$가 완전불연결임을 알 수 있다.

\medskip\noindent
(2)를 가정하자. 모든 $i \in I$에 대하여 $U_i$가 공집합이 아닌 열린
(따라서 닫힌) 집합인 유한 서로소 합집합 분해
$X = \coprod_{i \in I} U_i$들의 집합을 $\mathcal{I}$라 하자.
각 $I \in \mathcal{I}$에 대하여 $U_i$의 점을 $i$로 보내는 연속사상
$X \to I$가 있다. $I, I' \in \mathcal{I}$에 대하여 $I'$에 대응하는
덮개가 $I$에 대응하는 덮개를 세분할 필요충분조건으로 $I \leq I'$라
정의하여 부분순서를 준다. 이 경우 표준사상 $I' \to I$를 얻는다.
즉, $\mathcal{I}$ 위의 유한 이산공간들의 역계를 얻는다. 사상들
$X \to I$는 서로 양립하므로 연속사상
$$
X \longrightarrow \lim_{I \in \mathcal{I}} I
$$
를 얻는다. 이 사상이 위상동형이라고 주장하며, 그러면 증명이 끝난다.
(부분순서집합 $\mathcal{I}$가 유향이므로 마지막 주장도 따라온다. 실제로
$X$의 두 서로소 합집합 분해가 주어지면 둘 다 세분하는 세 번째 분해를
찾을 수 있다.) 이 사상은 단사이다. 실제로 $X$가 완전불연결이므로
$\{x\}$는 $x$를 포함하는 $X$의 모든 열리고 닫힌 부분집합들의 교집합이다
(Lemma \ref{topology-lemma-connected-component-intersection-compact-Hausdorff}).
또한 Lemma \ref{topology-lemma-intersection-closed-in-quasi-compact}에 의해 이 사상은
전사이다. Lemma \ref{topology-lemma-bijective-map}에 의해 이 사상은 위상동형이다.
\end{proof}

\begin{lemma}
\label{topology-lemma-directed-inverse-limit-profinite}
프로유한 공간들의 극한은 프로유한이다.
\end{lemma}

\begin{proof}
$\mathcal{I}$를 지표범주로 하는 프로유한 공간들의 도표를
$i \mapsto X_i$라 하자. Lemma \ref{topology-lemma-profinite}의 프로유한 공간에
대한 특징짓기를 사용하자. 특히 각 $X_i$는 하우스도르프이고 준콤팩트하며
완전불연결이다. Lemma \ref{topology-lemma-limits}에 의해 극한 $X = \lim X_i$가
존재한다. Lemma \ref{topology-lemma-inverse-limit-quasi-compact}에 의해 극한 $X$는
준콤팩트이다. 서로 다른 점 $x, x' \in X$를 택하자. 그러면 어떤 $i$가
존재하여, 사영사상 $X \to X_i$ 아래에서 $x$와 $x'$의 $X_i$에서의 상 $x_i$와
$x'_i$가 서로 다르다. 그러면 $x_i$와 $x'_i$는 $X_i$에서 서로소인
열린근방들을 가진다. 이 열린집합들의 역상을 취하면 $X$가 하우스도르프임을
알 수 있다. 마찬가지로 $x_i$와 $x'_i$는 $X_i$의 서로 다른 연결성분에
속하므로, $x$와 $x'$도 반드시 $X$의 서로 다른 연결성분에 속한다.
따라서 $X$는 완전불연결이다. 이것으로 증명이 끝난다.
\end{proof}

\begin{lemma}
\label{topology-lemma-profinite-refine-open-covering}
$X$를 프로유한 공간이라 하자. $X$의 모든 열린 덮개는 각 $U_i$가
열리고 닫힌 집합인 유한 덮개 $X = \coprod U_i$에 의한 세분을 가진다.
\end{lemma}

\begin{proof}
이 공간을 유향집합 $I$ 위의 유한 이산공간들의 역계의 극한
$X = \lim X_i$로 나타내자(Lemma \ref{topology-lemma-profinite}). 사영사상을
$f_i : X \to X_i$라 하자. 각 점 $x = (x_i) \in X$에 대하여
$f_i^{-1}(\{x_i\})$들의 모임은 열린근방들의 기본계이다. 따라서 $X$가
준콤팩트이므로, 다음과 같은 열린 덮개가 주어졌다고 가정해도 된다.
$$
X = f_{i_1}^{-1}(\{x_{i_1}\}) \cup \ldots \cup f_{i_n}^{-1}(\{x_{i_n}\})
$$
$I$가 유향집합이므로 가능한 바와 같이, $j = 1, \ldots, n$에 대하여
$i \geq i_j$인 $i \in I$를 택한다. 그러면 덮개
$$
X = \coprod\nolimits_{t \in X_i} f_i^{-1}(\{t\})
$$
가 주어진 덮개를 세분하고 원하는 꼴임을 알 수 있다.
\end{proof}

\begin{lemma}
\label{topology-lemma-pi0-profinite}
$X$를 위상공간이라 하자. $X$가 준콤팩트이고 $X$의 모든 연결성분이
그 연결성분을 포함하는 열리고 닫힌 부분집합들의 교집합이면,
$\pi_0(X)$는 프로유한 공간이다.
\end{lemma}

\begin{proof}
이를 증명하기 위해 Lemma \ref{topology-lemma-profinite}을 사용하겠다.
$\pi_0(X)$는 준콤팩트 공간의 상이므로 준콤팩트이다
(Lemma \ref{topology-lemma-image-quasi-compact}). 구성에 의해 이는 완전불연결이다
(Lemma \ref{topology-lemma-space-connected-components}). $C, D \subset X$를
$X$의 서로 다른 연결성분이라 하자. $C$를 포함하는 $X$의 열리고 닫힌
부분집합들의 교집합으로 $C = \bigcap U_\alpha$라 쓰자.
$U_\alpha$들의 임의의 유한 교집합도 같은 꼴이다.
$\bigcap U_\alpha \cap D = \emptyset$이므로 어떤 $\alpha$에 대하여
$U_\alpha \cap D = \emptyset$임을 알 수 있다(Lemmas
\ref{topology-lemma-connected-components}, \ref{topology-lemma-closed-in-quasi-compact},
\ref{topology-lemma-intersection-closed-in-quasi-compact}을 사용한다).
$U_\alpha$는 열리고 닫혀 있으므로 자신이 포함하는 연결성분들의 합집합이다.
즉, $U_\alpha$는 어떤 열리고 닫힌 부분집합
$V_\alpha \subset \pi_0(X)$의 역상이다.
이는 $C$와 $D$에 대응하는 점들이 서로소인 열린 부분집합들에 포함됨을
보인다. 따라서 $\pi_0(X)$는 하우스도르프이다.
\end{proof}




\section{스펙트럴 공간}
\label{topology-section-spectral}

\noindent
이 절의 내용은 \cite{Hochster}와 \cite{Hochster-thesis}에서 가져왔다.
Hochster는 그의 학위논문에서, 특히 스펙트럴 공간들이 어떤 환의 스펙트럼으로
나타나는 위상공간들과 정확히 일치함을 증명한다.

\begin{definition}
\label{topology-definition-spectral-space}
위상공간 $X$가 소버이고 준콤팩트이며, 두 준콤팩트 열린집합의 교집합이
준콤팩트이고, 준콤팩트 열린집합들의 모임이 위상의 기저를 이루면 이를
{\it 스펙트럴}이라고 한다. 스펙트럴 공간들의 연속사상
$f : X \to Y$에 대하여, 준콤팩트 열린집합의 역상이 준콤팩트이면 이 사상을
{\it 스펙트럴}이라고 한다.
\end{definition}

\noindent
다시 말해, 스펙트럴 공간들의 연속사상이 스펙트럴일 필요충분조건은 그 사상이
준콤팩트인 것이다(Definition \ref{topology-definition-quasi-compact}).

\medskip\noindent
$X$를 스펙트럴 공간이라 하자. $X$ 위의 {\it 구성가능 위상}은, $U$가
$X$의 준콤팩트 열린집합일 때 집합들 $U$와 $U^c$를 열린집합의 부분기저로
갖는 위상이다. $X$가 스펙트럴이므로 열린집합 $U \subset X$가
레트로콤팩트일 필요충분조건은 $U$가 준콤팩트인 것임에 유의하라. 따라서
구성가능 위상은 $X$의 모든 구성가능 부분집합이 열리고 닫히게 하는 가장
엉성한 위상으로도 특징지을 수 있다(구성가능 집합의 정의와 성질은 Section
\ref{topology-section-constructible}을 보라). 이에 따라 $X$의 부분집합이 구성가능
위상에서 열려 있을, 각각 닫혀 있을 필요충분조건은 그것이 구성가능
부분집합들의 합집합, 각각 교집합인 것이다. 준콤팩트 열린집합들의 모임은
$X$의 위상의 기저이므로, 구성가능 위상은 $X$에 주어진 위상보다 강하다.

\begin{lemma}
\label{topology-lemma-constructible-hausdorff-quasi-compact}
$X$를 스펙트럴 공간이라 하자. 구성가능 위상은 하우스도르프이고
완전불연결이며 준콤팩트이다.
\end{lemma}

\begin{proof}
$x, y \in X$이고 $x \not = y$라고 하자. $X$가 소버이므로 두 점 $x, y$ 중
정확히 하나를 포함하는 열린 부분집합 $U$가 있다. $x \in U$라고 하자.
$U$를 $U$에 포함되는 $x$의 준콤팩트 열린근방으로 바꾸어도 된다. 그러면
$U$와 $U^c$는 구성가능 위상에서 열리고 닫혀 있다. 따라서 구성가능 위상에서
$x \in U$와 $y \in U^c$가 서로소인 열린집합들이므로 $X$는 구성가능
위상에서 하우스도르프이다. $U$의 존재는 $x$와 $y$가 구성가능 위상에서
서로 다른 연결성분에 속함도 함의하므로, $X$는 구성가능 위상에서
완전불연결이다.

\medskip\noindent
$\mathcal{B}$를 $B \subset X$가 준콤팩트 열린집합이거나, 준콤팩트인
여집합을 갖는 닫힌집합인 그러한 부분집합 $B$들의 모임이라 하자.
$B \in \mathcal{B}$이면 $B^c \in \mathcal{B}$이다. Lemma
\ref{topology-lemma-subbase-theorem}에 의해, $B_i \in \mathcal{B}$인 모든 덮개
$X = \bigcup_{i \in I} B_i$가 유한 세분을 가짐을 보이면 충분하다.
여집합을 취하면, 모든 $n$과 모든 $i_1, \ldots, i_n \in I$에 대하여
$B_{i_1} \cap \ldots \cap B_{i_n} \not = \emptyset$을 만족하는
$\mathcal{B}$의 원소들의 임의의 족 $\{B_i\}_{i \in I}$가 공통 교점을
가짐을 보이면 된다. $i \not = i'$이면 $B_i \not = B_{i'}$라고 가정해도 된다.

\medskip\noindent
모순을 얻기 위해, $\{B_i\}_{i \in I}$가 유한교차성을 갖지만 교집합은
공집합인 $\mathcal{B}$의 원소들의 족이라고 하자. 초른 보조정리를 적용하면
우리의 족이 극대라고 가정해도 된다(세부사항은 생략한다).
$I' \subset I$를 $B_i$가 닫힌집합인 지표들의 집합이라 하고
$Z = \bigcap_{i \in I'} B_i$로 놓자. 이는 $X$의 닫힌 부분집합이고 Lemma
\ref{topology-lemma-intersection-closed-in-quasi-compact}에 의해 공집합이 아니다.
$Z$가 가약이면, 어느 것도 $Z$와 같지 않은 두 닫힌 부분집합의 합집합으로
$Z = Z' \cup Z''$라고 쓸 수 있다. 이는 특히 $Z'$와 만나지만 $Z''$와는
만나지 않는 준콤팩트 열린집합 $U' \subset X$를 찾을 수 있음을 뜻한다.
마찬가지로 $Z''$와 만나지만 $Z'$와는 만나지 않는 준콤팩트 열린집합
$U'' \subset X$를 찾을 수 있다. $B' = X \setminus U'$ 및
$B'' = X \setminus U''$로 놓자. $Z'' \subset B'$이고
$Z' \subset B''$임에 유의하라. 유한개의 지표 $i_1, \ldots, i_n \in I$가
존재하여 $B' \cap B_{i_1} \cap \ldots \cap B_{i_n} = \emptyset$이고,
또 유한개의 지표 $j_1, \ldots, j_m \in I$가 존재하여
$B'' \cap B_{j_1} \cap \ldots \cap B_{j_m} = \emptyset$이면 다음을 얻는다.
$Z \cap B_{i_1} \cap \ldots \cap B_{i_n} \cap B_{j_1} \cap \ldots \cap B_{j_m}
= \emptyset$.
그러나 집합
$B_{i_1} \cap \ldots \cap B_{i_n} \cap B_{j_1} \cap \ldots \cap B_{j_m}$
은 준콤팩트이므로, 유한개의 지표 $i'_1, \ldots, i'_l \in I'$에 대하여
$B_{i_1} \cap \ldots \cap B_{i_n} \cap B_{j_1} \cap \ldots \cap
B_{j_m} \cap B_{i'_1} \cap \ldots \cap B_{i'_l} = \emptyset$이 되어 모순이다.
따라서 주어진 족에 $B'$ 또는 $B''$ 중 하나를 더할 수 있으며, 이는 극대성에
모순이다. 그러므로 $Z$는 기약이다. 그러나 이 역시 모순을 낳는다. 위와 같은
논증에 의해 모든 $i \in I \setminus I'$에 대하여 공집합이 아닌 열린집합
$Z \cap B_i$가 $Z$의 유일한 일반점을 포함하기 때문이다. 이 모순으로
보조정리가 증명되었다.
\end{proof}

\begin{lemma}
\label{topology-lemma-fibres-spectral-map-quasi-compact}
$f : X \to Y$를 스펙트럴 공간들의 스펙트럴 사상이라 하자. 그러면
\begin{enumerate}
\item $f$는 구성가능 위상에서 연속이다.
\item $f$의 섬유들은 준콤팩트이다.
\item 그 상은 구성가능 위상에서 닫혀 있다.
\end{enumerate}
\end{lemma}

\begin{proof}
$X'$와 $Y'$를 각각 구성가능 위상을 부여한 $X$와 $Y$라 하자. Lemma
\ref{topology-lemma-constructible-hausdorff-quasi-compact}에 의해 이들은 준콤팩트
하우스도르프 공간들이다. (1)은 $X' \to Y'$가 연속이라는 말이며 정의에서
곧바로 따른다. (3)은 $f(X')$가 하우스도르프 공간 $Y'$의 준콤팩트
부분집합이라는 사실에서 따른다. Lemma
\ref{topology-lemma-quasi-compact-in-Hausdorff}를 보라. 위상공간들의 연속사상으로
이루어진 가환도표
$$
\xymatrix{
X' \ar[r] \ar[d] & X \ar[d] \\
Y' \ar[r] & Y
}
$$
를 얻는다. $Y'$가 하우스도르프이므로 섬유 $X'_y$들은 $X'$에서 닫혀 있다.
$X'$가 준콤팩트이므로 $X'_y$는 준콤팩트이다(Lemma
\ref{topology-lemma-closed-in-quasi-compact}). $X'_y \to X_y$가 전사인
연속사상이므로 $X_y$가 준콤팩트라고 결론 내린다(Lemma
\ref{topology-lemma-image-quasi-compact}).
\end{proof}

\begin{lemma}
\label{topology-lemma-spectral-if-continuous-wrt-constructible-top}
$X$와 $Y$를 스펙트럴 공간이라 하자. $f : X \to Y$를 연속사상이라 하자.
그러면 $f$가 스펙트럴일 필요충분조건은 $f$가 구성가능 위상에서 연속인 것이다.
\end{lemma}

\begin{proof}
필요성은 Lemma \ref{topology-lemma-fibres-spectral-map-quasi-compact}이다.
$f$가 구성가능 위상에서 연속이라고 가정하자. $V \subset Y$를 준콤팩트
열린집합이라 하자. 그러면 $V$는 구성가능 위상에서 열리고 닫혀 있다.
따라서 $f^{-1}(V)$는 구성가능 위상에서 열리고 닫혀 있다. Lemma
\ref{topology-lemma-constructible-hausdorff-quasi-compact}에 의해 $X$가 구성가능
위상에서 준콤팩트이므로 $f^{-1}(V)$는 구성가능 위상에서 준콤팩트이다.
항등사상 $f^{-1}(V) \to f^{-1}(V)$가 구성가능 위상에서 보통 위상으로 가는
전사인 연속사상이므로, Lemma \ref{topology-lemma-image-quasi-compact}에 의해
$f^{-1}(V)$가 $X$의 위상에서 준콤팩트라고 결론 내린다. 이것으로 증명이 끝난다.
\end{proof}

\begin{lemma}
\label{topology-lemma-spectral-sub}
$X$를 스펙트럴 공간이라 하자. $E \subset X$가 구성가능 위상에서 닫혀
있다고 하자(예를 들어 구성가능이거나 닫혀 있다고 하자). 그러면 유도위상을
부여한 $E$는 스펙트럴 공간이다.
\end{lemma}

\begin{proof}
$Z \subset E$를 기약 닫힌 부분집합이라 하자. $X$에서 $Z$의 폐포
$\overline{Z}$의 일반점을 $\eta$라 하자. $E$가 소버임을 보이기 위해
$\eta \in E$임을 보이자. 그렇지 않다면 $E$가 구성가능 위상에서 닫혀
있으므로, $\eta \in F$이고 $F \cap E = \emptyset$인 구성가능 부분집합
$F \subset X$가 존재한다. Lemma \ref{topology-lemma-generic-point-in-constructible}에
의해 이는 $F \cap \overline{Z}$가 $\overline{Z}$의 공집합이 아닌 열린
부분집합을 포함함을 뜻한다. 그러나 $\overline{Z}$가 $Z$의 폐포이고
$Z \cap F = \emptyset$이므로 이는 불가능하다.

\medskip\noindent
$E$가 구성가능 위상에서 닫혀 있으므로 구성가능 위상에서 준콤팩트이다
(Lemmas \ref{topology-lemma-closed-in-quasi-compact}와
\ref{topology-lemma-constructible-hausdorff-quasi-compact}). 따라서 하물며 $X$에서
오는 위상에서도 준콤팩트이다. $U \subset X$가 준콤팩트 열린집합이면
$E \cap U$는 구성가능 위상에서 닫혀 있고, 따라서 위에서 보았듯이
준콤팩트이다. 이에 따라 $E$의 준콤팩트 열린 부분집합들은 $X$의 준콤팩트
열린집합 $U$에 대한 교집합들 $E \cap U$이고, 이들은 위상의 기저를 이룬다.
마지막으로 준콤팩트 열린집합들 $U, U' \subset X$가 주어졌을 때,
$(E \cap U) \cap (E \cap U') = E \cap (U \cap U')$이고 $X$가
스펙트럴이므로 $U \cap U'$는 준콤팩트이다. 이것으로 증명이 끝난다.
\end{proof}

\begin{lemma}
\label{topology-lemma-constructible-stable-specialization-closed}
$X$를 스펙트럴 공간이라 하자. $E \subset X$를 구성가능 위상에서 닫힌
부분집합이라 하자(예를 들어 구성가능 부분집합이라 하자).
\begin{enumerate}
\item $x \in \overline{E}$이면, $x$는 $E$의 한 점의 특수화이다.
\item $E$가 특수화에 대해 안정이면 $E$는 닫혀 있다.
\item $E' \subset X$가 구성가능 위상에서 열려 있고(예를 들어 구성가능이고)
일반화에 대해 안정이면 $E'$은 열려 있다.
\end{enumerate}
\end{lemma}

\begin{proof}
(1)의 증명. $x \in \overline{E}$라고 하자. $\{U_i\}$를 $x$의 준콤팩트
열린근방들의 집합이라 하자. $U_i$들의 유한 교집합도 같은 꼴이다.
모든 $i$에 대하여 교집합 $U_i \cap E$는 공집합이 아니다. 부분집합들
$U_i \cap E$가 구성가능 위상에서 닫혀 있으므로, Lemma
\ref{topology-lemma-constructible-hausdorff-quasi-compact}와 Lemma
\ref{topology-lemma-intersection-closed-in-quasi-compact}에 의해
$\bigcap (U_i \cap E)$가 공집합이 아님을 알 수 있다. $\{U_i\}$가 $x$의
열린근방들의 기본계이므로 $\bigcap U_i$는 $x$의 일반화들의 집합이다.
따라서 $x$는 $E$의 한 점의 특수화이다.

\medskip\noindent
(2)는 (1)에서 곧바로 따른다.

\medskip\noindent
(3)의 증명. $E'$이 (3)과 같다고 가정하자. $E'$의 여집합은 구성가능
위상에서 닫혀 있고(Lemma \ref{topology-lemma-constructible}), 특수화에 대해
닫혀 있다(Lemma \ref{topology-lemma-open-closed-specialization}). 따라서 (2)에
의해 그 여집합은 닫혀 있다. 즉, $E'$은 열려 있다.
\end{proof}

\begin{lemma}
\label{topology-lemma-two-points}
% P01-E1213: the frozen wording requires a third common generalization.
$X$를 스펙트럴 공간이라 하자. $x, y \in X$라 하자. 그러면 $x$와 $y$ 둘
모두로 특수화되는 세 번째 점이 존재하거나, $x$와 $y$를 포함하는 서로소인
열린근방들이 존재한다.
\end{lemma}

\begin{proof}
$\{U_i\}$를 $x$의 준콤팩트 열린근방들의 집합이라 하자. $U_i$들의 유한
교집합도 같은 꼴이다. $\{V_j\}$를 $y$의 준콤팩트 열린근방들의 집합이라
하자. $V_j$들의 유한 교집합도 같은 꼴이다. 어떤 $i, j$에 대하여
$U_i \cap V_j$가 공집합이면 증명이 끝난다. 그렇지 않으면 모든 $i$와
$j$에 대하여 교집합 $U_i \cap V_j$가 공집합이 아니다. 집합들
$U_i \cap V_j$는 $X$ 위의 구성가능 위상에서 닫혀 있다. Lemma
\ref{topology-lemma-constructible-hausdorff-quasi-compact}에 의해
$\bigcap (U_i \cap V_j)$가 공집합이 아님을 알 수 있다(Lemma
\ref{topology-lemma-intersection-closed-in-quasi-compact}). $X$가 소버 공간이고
$\{U_i\}$가 $x$의 열린근방들의 기본계이므로, $\bigcap U_i$는 $x$의
일반화들의 집합이다. 마찬가지로 $\bigcap V_j$는 $y$의 일반화들의
집합이다. 따라서 $\bigcap (U_i \cap V_j)$의 임의의 원소는 $x$와 $y$
둘 모두로 특수화된다.
\end{proof}

\begin{lemma}
\label{topology-lemma-characterize-profinite-spectral}
$X$를 스펙트럴 공간이라 하자. 다음은 동치이다.
\begin{enumerate}
\item $X$는 프로유한이다.
\item $X$는 하우스도르프이다.
\item $X$는 완전불연결이다.
\item 모든 준콤팩트 열린집합은 닫혀 있다.
\item 점들 사이에 자명하지 않은 특수화가 없다.
\item $X$의 모든 점은 닫혀 있다.
\item $X$의 모든 점은 $X$의 한 기약성분의 일반점이다.
\item 구성가능 위상은 $X$에 주어진 위상과 같다.
\item 여기에 더 추가한다.
\end{enumerate}
\end{lemma}

\begin{proof}
Lemma \ref{topology-lemma-profinite}은 (1) $\Rightarrow$ (3)을 보인다. 기약성분은
닫혀 있으므로 $X$가 완전불연결이면 모든 점은 닫혀 있다. 따라서 (3)은
(6)을 함의한다. (6)과 (5)의 동치는 곧바로 알 수 있고, $X$가 소버이므로
(6) $\Leftrightarrow$ (7)이 성립한다. (5)를 가정하자. 그러면 $X$의
모든 구성가능 부분집합은 닫혀 있고(Lemma
\ref{topology-lemma-constructible-stable-specialization-closed}), 특히 모든
준콤팩트 열린집합은 닫혀 있다. 따라서 (5)는 (4)를 함의한다. $X$가
소버이므로, 임의의 두 점에 대하여 정확히 하나의 점을 포함하는 준콤팩트
열린집합이 있다. 따라서 (4)는 (2)를 함의한다. (4)와 (8)은 구성가능
위상의 정의에 의해 동치이다. 이제 (2)가 (1)을 함의함을 보이는 일만
남았다. $X$가 하우스도르프라고 하자. 모든 준콤팩트 열린집합은 또한
닫혀 있다(Lemma \ref{topology-lemma-quasi-compact-in-Hausdorff}). 이는 $X$가
완전불연결임을 함의한다. 따라서 Lemma \ref{topology-lemma-profinite}에 의해
프로유한이다.
\end{proof}

\begin{lemma}
\label{topology-lemma-spectral-pi0}
$X$가 스펙트럴 공간이면 $\pi_0(X)$는 프로유한 공간이다.
\end{lemma}

\begin{proof}
Lemmas \ref{topology-lemma-connected-component-intersection}와
\ref{topology-lemma-pi0-profinite}를 결합하라.
\end{proof}

\begin{lemma}
\label{topology-lemma-product-spectral-spaces}
두 스펙트럴 공간의 곱은 스펙트럴이다.
\end{lemma}

\begin{proof}
$X$, $Y$를 스펙트럴 공간이라 하자. 사영사상들을 각각
$p : X \times Y \to X$와 $q : X \times Y \to Y$라 하자.
$Z \subset X \times Y$를 기약 닫힌 부분집합이라 하자. 그러면
$p(Z) \subset X$는 기약이고 $q(Z) \subset Y$도 기약이다.
% P01-E0023/P01-E1215: the frozen proof uses p(X), q(Y), not p(Z), q(Z).
$p(X)$의 폐포의 일반점을 $x \in X$라 하고 $q(Y)$의 폐포의 일반점을
$y \in Y$라 하자. $(x, y) \not \in Z$이면, $x \in U \subset X$와
$y \in V \subset Y$를 만족하고
% P01-E1216: the frozen product-neighbourhood intersection lacks parentheses.
$Z \cap U \times V = \emptyset$인 열린집합들이 존재한다. 따라서 $Z$는
$(X \setminus U) \times Y \cup X \times (Y \setminus V)$에 포함된다.
$Z$가 기약이므로 $Z \subset (X \setminus U) \times Y$이거나
$Z \subset X \times (Y \setminus V)$임을 알 수 있다. 첫 번째 경우에는
$p(Z) \subset (X \setminus U)$이고 두 번째 경우에는
$q(Z) \subset (Y \setminus V)$이다. $x$가 $p(Z)$의 폐포에 있고 $y$가
$q(Z)$의 폐포에 있으므로 두 경우 모두 모순이다. 따라서
$(x, y) \in Z$라고 결론 내리며, 이는 $(x, y)$가 $Z$의 일반점이라는
뜻이다.

\medskip\noindent
$X \times Y$의 위상의 한 기저는 $U \subset X$와 $V \subset Y$가
준콤팩트 열린집합인 그러한 $U \times V$ 꼴의 열린집합들로 이루어진다
(여기서는 $X$와 $Y$가 스펙트럴임을 사용한다). 그러면 $U \times V$는
준콤팩트 공간들의 곱이므로 준콤팩트이다. 또한 $X \times Y$의 임의의
준콤팩트 열린집합은 그러한 준콤팩트 직사각형들 $U \times V$의 유한
합집합이다. 두 그러한 집합의 교집합도 다시 준콤팩트임이 따라온다
($X$와 $Y$가 스펙트럴이기 때문이다). 이것으로 증명이 끝난다.
\end{proof}

\begin{lemma}
\label{topology-lemma-spectral-bijective}
$f : X \to Y$를 위상공간들의 연속사상이라 하자. 다음이 성립한다고 하자.
\begin{enumerate}
\item $X$와 $Y$는 스펙트럴이다.
\item $f$는 스펙트럴이고 전단사이다.
\item 일반화(각각 특수화)를 $f$를 따라 끌어올릴 수 있다.
\end{enumerate}
그러면 $f$는 위상동형이다.
\end{lemma}

\begin{proof}
$f$가 스펙트럴이므로 구성가능 위상을 부여한 $X$와 $Y$ 사이의 연속사상을
정의한다. Lemmas \ref{topology-lemma-constructible-hausdorff-quasi-compact}와
\ref{topology-lemma-bijective-map}에 의해 $X \to Y$가 구성가능 위상에서
위상동형임이 따라온다. $U \subset X$를 준콤팩트 열린집합이라 하자.
그러면 $f(U)$는 $Y$에서 구성가능이다. $y \in Y$가 $f(U)$의 한 점으로
특수화된다고 하자. 마지막 가정에 의해 $f^{-1}(y)$가 $U$의 한 점으로
특수화됨을 알 수 있다. 따라서 $f^{-1}(y) \in U$이다. 그러므로
$y \in f(U)$이다. 이에 따라 $f(U)$는 열려 있다. Lemma
\ref{topology-lemma-constructible-stable-specialization-closed}를 보라. 따라서
$f$는 위상동형이다. 특수화를 $f$를 따라 끌어올릴 수 있는 경우 보조정리를
증명하려면, 대신 $X \setminus Z$가 $X$의 준콤팩트 열린집합일 때
$f(Z)$가 닫혀 있음을 보인다.
\end{proof}

\begin{lemma}
\label{topology-lemma-directed-inverse-limit-finite-sober-spectral-spaces}
유한 소버 위상공간들로 이루어진 유향 역계의 역극한은 스펙트럴
위상공간이다.
\end{lemma}

\begin{proof}
$I$를 유향집합이라 하자. $X_i$를 $I$ 위의 유한 소버 공간들의 역계라 하자.
Lemma \ref{topology-lemma-limits}에 의해 존재하는 $X = \lim X_i$라 하자. 집합으로도
$X = \lim X_i$이다. 사영사상을 $p_i : X \to X_i$라 하자. $I$가
유향이므로 Lemma \ref{topology-lemma-describe-limits}를 적용할 수 있다. 위상의 한
기저는 $U_i \subset X_i$가 열린집합일 때의 열린집합들
$p_i^{-1}(U_i)$로 주어진다. $p_i^{-1}(U_i)$의 열린 덮개는 특히 프로유한
위상에서의 열린 덮개이므로, $p_i^{-1}(U_i)$가 준콤팩트이라고 결론 내린다.
$U_i \subset X_i$와 $U_j \subset X_j$가 주어지면, 어떤 $k \geq i, j$와
열린집합 $U_k \subset X_k$에 대하여
$p_i^{-1}(U_i) \cap p_j^{-1}(U_j) = p_k^{-1}(U_k)$이다. 마지막으로
$Z \subset X$가 기약이고 닫혀 있으면 $p_i(Z) \subset X_i$는 기약이고,
따라서 유일한 일반점 $\xi_i$를 가진다($X_i$가 유한 소버 위상공간이기
때문이다). 그러면 $\xi = \lim \xi_i$는 $Z$의 일반점이다($Z$가 닫혀
있으므로 이는 $Z$의 한 점이다). 이것으로 증명이 끝난다.
\end{proof}

\begin{lemma}
\label{topology-lemma-spectral-closed-in-product-two-point-space}
$W$를 두 점으로 이루어지고 그중 하나는 닫혀 있지만 다른 하나는 닫혀 있지
않은 위상공간이라 하자. 위상공간이 스펙트럴일 필요충분조건은 구성가능
위상에서 닫힌 $W$의 복사본들의 곱의 부분공간과 위상동형인 것이다.
\end{lemma}

\begin{proof}
$W = \{0, 1\}$로 쓰되 $0$은 $1$의 특수화이지만 그 역은 아니라고 하자.
$I$를 집합이라 하자. 공간 $\prod_{i \in I} W$는 Lemma
\ref{topology-lemma-directed-inverse-limit-finite-sober-spectral-spaces}에 의해
스펙트럴이다. 따라서 구성가능 위상에서 닫힌 $\prod_{i \in I} W$의
부분공간은 Lemma \ref{topology-lemma-spectral-sub}에 의해 스펙트럴 공간임을 알 수 있다.

\medskip\noindent
역으로 $X$를 스펙트럴 공간이라 하자. $U \subset X$를 준콤팩트
열린집합이라 하자. $U$의 모든 점을 $1$로 보내고 $X \setminus U$의 모든
점을 $0$으로 보내는 연속사상
$$
f_U : X \longrightarrow W
$$
을 생각하자. 이 사상들의 곱을 취하여 연속사상
$$
f = \prod f_U : X \longrightarrow \prod\nolimits_U W
$$
를 얻는다.
% P01-E0024/P01-E1218: the frozen proof uses Y before defining Y as f(X).
구성에 의해 사상 $f : X \to Y$는 스펙트럴이다. Lemma
\ref{topology-lemma-fibres-spectral-map-quasi-compact}에 의해 $f$의 상은 구성가능
위상에서 닫혀 있다. $x', x \in X$가 서로 다르면, $X$가 소버이므로
$x'$이 $x$의 특수화가 아니거나 그 역이 성립한다. 어느 경우든
($X$의 위상의 한 기저를 준콤팩트 열린집합들이 이루므로)
$f_U(x') \not = f_U(x)$인 준콤팩트 열린집합 $U \subset X$가 존재한다.
따라서 $f$는 단사이다. $Y = f(X)$에 유도위상을 부여하자.
$Y$에서의 특수화를 $y' \leadsto y$라 쓰고, $f(x') = y'$ 및
$f(x) = y$라고 하자. 위와 같이 논증하면 $x' \leadsto x$이다. 그렇지 않으면
$x \in U$이고 $x' \not \in U$인 $U$가 존재하고, 이는
$f_U(x') \not \leadsto f_U(x)$를 함의할 것이기 때문이다. Lemma
\ref{topology-lemma-spectral-bijective}에 의해 $f : X \to Y$가 위상동형이라고
결론 내린다.
\end{proof}

\begin{lemma}
\label{topology-lemma-spectral-inverse-limit-finite-sober-spaces}
위상공간이 스펙트럴일 필요충분조건은 유한 소버 위상공간들의 유향 역극한인
것이다.
\end{lemma}

\begin{proof}
한 방향은 Lemma
\ref{topology-lemma-directed-inverse-limit-finite-sober-spectral-spaces}에서 주어진다.
역으로 $X$가 스펙트럴이라고 가정하자. 그러면 Lemma
\ref{topology-lemma-spectral-closed-in-product-two-point-space}에서처럼
$W = \{0, 1\}$일 때 $X \subset \prod_{i \in I} W$가 구성가능 위상에서
닫힌 부분집합이라고 가정해도 된다. 다음과 같이 쓸 수 있다.
$$
\prod\nolimits_{i \in I} W =
\lim_{J \subset I\text{ finite }} \prod\nolimits_{j \in J} W
$$
이는 쌍대여과 극한이다. 각 $J$에 대하여 $X_J \subset \prod_{j \in J} W$를
$X$의 상이라 하자. $X$가 구성가능 위상을 부여한 곱에서 닫혀 있으므로,
집합으로서 $X = \lim X_J$임을 알 수 있다(세부사항은 생략한다).
극한에 관한 형식적 논증(생략한다)은 위상공간으로서도
$X = \lim X_J$임을 보인다.
\end{proof}

\begin{lemma}
\label{topology-lemma-Noetherian-goes-to-spectral}
$X$를 위상공간이라 하고, $c : X \to X'$를 $X$에서 소버 위상공간으로 가는
보편사상이라 하자. Lemma \ref{topology-lemma-make-sober}를 보라.
\begin{enumerate}
\item $X$가 준콤팩트이면 $X'$도 그러하다.
\item $X$가 준콤팩트이고 준콤팩트 열린집합들로 이루어진 기저를 가지며,
두 준콤팩트 열린집합의 교집합이 준콤팩트이면 $X'$은 스펙트럴이다.
\item $X$가 뇌터이면 $X'$은 뇌터 스펙트럴 공간이다.
\end{enumerate}
\end{lemma}

\begin{proof}
$U \subset X$를 열린집합이라 하고 $U' \subset X'$을 대응하는 열린집합,
즉 $c^{-1}(U') = U$인 열린집합이라 하자. $U$가 준콤팩트일 필요충분조건은
$U'$이 준콤팩트인 것이다. 실제로 $c$에 의한 역상은 합집합과 교환하는
$X$의 열린집합들과 $X'$의 열린집합들 사이의 전단사이기 때문이다. 이는
특히 (1)을 증명한다.

\medskip\noindent
(2)의 증명. 위의 사실로부터 $X'$이 준콤팩트 열린집합들로 이루어진 기저를
가짐이 따라온다. $c^{-1}$는 열린집합들의 쌍의 교집합과도 교환하므로,
$X'$의 두 준콤팩트 열린집합의 교집합은 준콤팩트임을 알 수 있다. 마지막으로
$X'$은 (1)에 의해 준콤팩트이고 구성에 의해 소버이다. 따라서 $X'$은
스펙트럴이다.

\medskip\noindent
(3)의 증명. 이 성질은 $X$에서 성립하는 열린 부분집합들에 대한 오름사슬
조건으로 정의되므로 $X'$이 뇌터라는 것은 곧바로 알 수 있다. (2)에서
$X'$이 스펙트럴임을 이미 보았다.
\end{proof}




\section{스펙트럴 공간들의 극한}
\label{topology-section-limits-spectral}

\noindent
Lemma \ref{topology-lemma-spectral-inverse-limit-finite-sober-spaces}은 모든 스펙트럴
공간이 유한 소버 공간들의 쌍대여과 극한임을 말해 준다. 모든 유한 소버
공간은 스펙트럴 공간이고, 유한 소버 공간들의 모든 연속사상은 스펙트럴
공간들의 스펙트럴 사상이다. 이 절에서는 스펙트럴 사상들로 이루어진
스펙트럴 위상공간들의 계의 극한에 관한 몇 가지 보조정리를 증명한다.

\begin{lemma}
\label{topology-lemma-inverse-limit-spectral-spaces-quasi-compact}
$\mathcal{I}$를 범주라 하자. $i \mapsto X_i$를 스펙트럴 공간들의 도표라
하고, $\mathcal{I}$의 $a : j \to i$에 대응하는 사상
$f_a : X_j \to X_i$가 스펙트럴이라고 하자.
\begin{enumerate}
\item 구성가능 위상에서 닫힌 부분집합들 $Z_i \subset X_i$가 주어지고,
$\mathcal{I}$의 모든 $a : j \to i$에 대하여 $f_a(Z_j) \subset Z_i$이면,
$\lim Z_i$는 준콤팩트이다.
\item 공간 $X = \lim X_i$는 준콤팩트이다.
\end{enumerate}
\end{lemma}

\begin{proof}
극한 $Z = \lim Z_i$는 Lemma \ref{topology-lemma-limits}에 의해 존재한다.
$X'_i$를 구성가능 위상을 부여한 공간 $X_i$라 하고, $Z'_i$를 $X'_i$의
대응하는 부분공간이라 하자. $\mathcal{I}$의 사상 $a : j \to i$를 택하자.
$f_a$가 스펙트럴이므로 연속사상 $f_a : X'_j \to X'_i$를 정의한다. 따라서
$f_a|_{Z_j} : Z'_j \to Z'_i$는 준콤팩트 하우스도르프 공간들의 연속사상이다
(Lemmas \ref{topology-lemma-constructible-hausdorff-quasi-compact}와
\ref{topology-lemma-closed-in-quasi-compact}에 의한다). 그러므로 Lemma
\ref{topology-lemma-inverse-limit-quasi-compact}에 의해 $Z' = \lim Z_i$는
준콤팩트이다. 사상들 $Z'_i \to Z_i$가 연속이므로 $Z' \to Z$는 연속이고
바탕집합들 위에서 전단사이다. 따라서 전사 연속사상 $Z' \to Z$의 상인
$Z$는 준콤팩트이다(Lemma
\ref{topology-lemma-image-quasi-compact}).
\end{proof}

\begin{lemma}
\label{topology-lemma-inverse-limit-spectral-spaces-nonempty}
$\mathcal{I}$를 쌍대여과 범주라 하자. $i \mapsto X_i$를 스펙트럴 공간들의
도표라 하고, $\mathcal{I}$의 $a : j \to i$에 대응하는 사상
$f_a : X_j \to X_i$가 스펙트럴이라고 하자.
\begin{enumerate}
\item 구성가능 위상에서 닫힌 공집합이 아닌 부분집합들 $Z_i \subset X_i$가
주어지고, $\mathcal{I}$의 모든 $a : j \to i$에 대하여
$f_a(Z_j) \subset Z_i$이면 $\lim Z_i$는 공집합이 아니다.
\item 각 $X_i$가 공집합이 아니면 $X = \lim X_i$도 공집합이 아니다.
\end{enumerate}
\end{lemma}

\begin{proof}
$X'_i$를 구성가능 위상을 부여한 공간 $X_i$라 하고, $Z'_i$를 $X'_i$의
대응하는 부분공간이라 하자. $\mathcal{I}$의 사상 $a : j \to i$를 택하자.
$f_a$가 스펙트럴이므로 연속사상 $f_a : X'_j \to X'_i$를 정의한다. 따라서
$f_a|_{Z_j} : Z'_j \to Z'_i$는 준콤팩트 하우스도르프 공간들의 연속사상이다
(Lemmas \ref{topology-lemma-constructible-hausdorff-quasi-compact}와
\ref{topology-lemma-closed-in-quasi-compact}에 의한다). Lemma
\ref{topology-lemma-nonempty-limit}에 의해 공간 $\lim Z'_i$는 공집합이 아니다.
집합으로서 $\lim Z'_i = \lim Z_i$이므로 결론을 얻는다.
\end{proof}

\begin{lemma}
\label{topology-lemma-inverse-limit-spectral-spaces-equal}
$\mathcal{I}$를 쌍대여과 범주라 하자. $i \mapsto X_i$를 스펙트럴 공간들의
도표라 하고, $\mathcal{I}$의 $a : j \to i$에 대응하는 사상
$f_a : X_j \to X_i$가 스펙트럴이라고 하자. 사영사상들이
$p_i : X \to X_i$인 $X = \lim X_i$라 하자. $i \in \Ob(\mathcal{I})$라
하고, $E$는 구성가능 위상에서 닫혀 있고 $F$는 구성가능 위상에서 열려 있는
부분집합들 $E, F \subset X_i$라 하자. 그러면
$p_i^{-1}(E) \subset p_i^{-1}(F)$일 필요충분조건은 $\mathcal{I}$에 사상
$a : j \to i$가 존재하여 $f_a^{-1}(E) \subset f_a^{-1}(F)$가 되는 것이다.
\end{lemma}

\begin{proof}
다음을 관찰하라.
% P01-E1221: the frozen RHS omits grouping around each set difference.
$$
p_i^{-1}(E) \setminus p_i^{-1}(F) =
\lim_{a : j \to i} f_a^{-1}(E) \setminus f_a^{-1}(F)
$$
$f_a$가 스펙트럴 사상이므로 구성가능 위상에서 연속이고, 따라서 집합
$f_a^{-1}(E) \setminus f_a^{-1}(F)$는 구성가능 위상에서 닫혀 있다.
그러므로 Lemma \ref{topology-lemma-inverse-limit-spectral-spaces-nonempty}를 적용하면,
좌변이 공집합이 아닐 필요충분조건은 우변의 공간들 각각이 공집합이 아닌
것임을 알 수 있다.
\end{proof}

\begin{lemma}
\label{topology-lemma-inverse-limit-spectral-spaces-constructible}
$\mathcal{I}$를 쌍대여과 범주라 하자. $i \mapsto X_i$를 스펙트럴 공간들의
도표라 하고, $\mathcal{I}$의 $a : j \to i$에 대응하는 사상
$f_a : X_j \to X_i$가 스펙트럴이라고 하자. 사영사상들이
$p_i : X \to X_i$인 $X = \lim X_i$라 하자. $E \subset X$를 구성가능
부분집합이라 하자. 그러면 어떤 $i \in \Ob(\mathcal{I})$와 구성가능
부분집합 $E_i \subset X_i$가 존재하여 $p_i^{-1}(E_i) = E$가 된다. $E$가
열려 있을 때, 각각 닫혀 있을 때에는 $E_i$도 열려 있게, 각각 닫혀 있게
택할 수 있다.
\end{lemma}

\begin{proof}
$E$가 $X$의 준콤팩트 열린집합이라고 가정하자. Lemma
\ref{topology-lemma-describe-limits}에 의해 어떤 $i$와 어떤 열린집합
$U_i \subset X_i$에 대하여 $E = p_i^{-1}(U_i)$라고 쓸 수 있다.
$U_i = \bigcup U_{i, \alpha}$를 준콤팩트 열린집합들의 합집합으로 쓰자.
$E$가 준콤팩트이므로 $E = p_i^{-1}(U_{i, \alpha_1} \cup \ldots \cup U_{i, \alpha_n})$가
되도록 $\alpha_1, \ldots, \alpha_n$을 찾을 수 있다. 따라서
$E_i = U_{i, \alpha_1} \cup \ldots \cup U_{i, \alpha_n}$이면 된다.

\medskip\noindent
$E$가 구성가능인 닫힌 부분집합이라고 가정하자. 그러면 $E^c$는 준콤팩트
열린집합이다. 따라서 앞 문단의 결과에 의해 어떤 $i$와 준콤팩트 열린집합
$F_i \subset X_i$에 대하여 $E^c = p_i^{-1}(F_i)$이다. 그러면 원하는 대로
$E = p_i^{-1}(F_i^c)$이다.

\medskip\noindent
$E$가 일반적인 경우에는 $U_l$이 구성가능 열린집합이고 $Z_l$이 구성가능
닫힌집합인 그러한 식
$E = \bigcup_{l = 1, \ldots, n} U_l \cap Z_l$로 쓸 수 있다. 앞 문단들의
결과에 의해 $U_l = p_{i_l}^{-1}(U_{l, i_l})$ 및
$Z_l = p_{j_l}^{-1}(Z_{l, j_l})$로 쓸 수 있다. 여기서
$U_{l, i_l} \subset X_{i_l}$는 구성가능 열린집합이고
$Z_{l, j_l} \subset X_{j_l}$는 구성가능 닫힌집합이다. $\mathcal{I}$가
쌍대여과이므로 $\mathcal{I}$의 대상 $k$와 사상들 $a_l : k \to i_l$ 및
$b_l : k \to j_l$을 택할 수 있다. 그러면
$E_k = \bigcup_{l = 1, \ldots, n}
f_{a_l}^{-1}(U_{l, i_l}) \cap f_{b_l}^{-1}(Z_{l, j_l})$로 놓아, $X$에서의
역상이 $E$인 $X_k$의 구성가능 부분집합을 얻는다.
\end{proof}

\begin{lemma}
\label{topology-lemma-directed-inverse-limit-spectral-spaces}
$\mathcal{I}$를 쌍대여과 지표범주라 하자. $i \mapsto X_i$를 스펙트럴
공간들의 도표라 하고, $\mathcal{I}$의 $a : j \to i$에 대응하는 사상
$f_a : X_j \to X_i$가 스펙트럴이라고 하자. 그러면 역극한
$X = \lim X_i$는 스펙트럴 위상공간이고 사영사상들
$p_i : X \to X_i$는 스펙트럴이다.
\end{lemma}

\begin{proof}
극한 $X = \lim X_i$는 존재하고(Lemma \ref{topology-lemma-limits}), Lemma
\ref{topology-lemma-inverse-limit-spectral-spaces-quasi-compact}에 의해 준콤팩트이다.

\medskip\noindent
사영사상을 $p_i : X \to X_i$라 하자. $\mathcal{I}$가 쌍대여과이므로
Lemma \ref{topology-lemma-describe-limits}를 적용할 수 있다. 따라서 $X$의 위상의 한
기저는 $U_i \subset X_i$가 열린집합일 때 열린집합들 $p_i^{-1}(U_i)$로
주어진다. $X_i$의 위상의 한 기저를 준콤팩트 열린집합들이 이루므로, $X$의
위상의 한 기저를 $U_i \subset X_i$가 준콤팩트 열린집합일 때의
$p_i^{-1}(U_i)$들이 이룬다고 결론 내린다. 형식적 논증은 위상공간으로서
다음이 성립함을 보인다.
$$
p_i^{-1}(U_i) = \lim_{a : j \to i} f_a^{-1}(U_i)
$$
각 $f_a$가 스펙트럴이므로 집합들 $f_a^{-1}(U_i)$는 $X_j$의 구성가능
위상에서 닫혀 있고, 따라서 Lemma
\ref{topology-lemma-inverse-limit-spectral-spaces-quasi-compact}에 의해
$p_i^{-1}(U_i)$는 준콤팩트이다. 그러므로 $X$의 위상은 준콤팩트
열린집합들로 이루어진 기저를 가진다.

\medskip\noindent
$X$의 임의의 준콤팩트 열린집합 $U$는 어떤 $i$와 어떤 준콤팩트 열린집합
$U_i \subset X_i$에 대하여 $U = p_i^{-1}(U_i)$ 꼴이다(Lemma
\ref{topology-lemma-inverse-limit-spectral-spaces-constructible}을 보라).
$U_i \subset X_i$와 $U_j \subset X_j$가 준콤팩트 열린집합들이면, 어떤
$k$와 준콤팩트 열린집합 $U_k \subset X_k$에 대하여
$p_i^{-1}(U_i) \cap p_j^{-1}(U_j) = p_k^{-1}(U_k)$이다. 실제로 $k$와
사상들 $k \to i$ 및 $k \to j$를 택하고, $U_k$를 $U_i$와 $U_j$의
$X_k$로의 역상들의 교집합으로 놓으면 된다. 이로써 $X$의 두 준콤팩트
열린집합의 교집합이 준콤팩트 열린집합임을 알 수 있다.

\medskip\noindent
마지막으로 $Z \subset X$가 기약이고 닫혀 있다고 하자. 그러면
$p_i(Z) \subset X_i$는 기약이고, 따라서 $Z_i = \overline{p_i(Z)}$는
유일한 일반점 $\xi_i$를 가진다($X_i$가 스펙트럴 공간이기 때문이다).
$\mathcal{I}$의 $a : j \to i$에 대하여 $\overline{f_a(Z_j)} = Z_i$이므로
$f_a(\xi_j) = \xi_i$이다. 따라서 $\xi = \lim \xi_i$는 $X$의 한 점이다.
주장: $\xi \in Z$이다. 그렇지 않으면 $\xi$를 포함하고 $Z$와 만나지 않는
준콤팩트 열린집합을 찾을 수 있다. 이는 어떤 $i$와 준콤팩트 열린집합
$U_i \subset X_i$에 대하여 $p_i^{-1}(U_i)$ 꼴일 것이다. 그러면
$\xi_i \in U_i$이지만 $p_i(Z) \cap U_i = \emptyset$이며, 이는
$\xi_i \in \overline{p_i(Z)}$에 모순이다. 따라서 $\xi \in Z$이고
$\overline{\{\xi\}} \subset Z$이다. 역으로 모든 $z \in Z$는 $\xi$의
폐포에 속한다. 실제로 $z$의 준콤팩트 열린근방 $U$가 주어지면, 어떤 $i$와
준콤팩트 열린집합 $U_i \subset X_i$에 대하여 $U = p_i^{-1}(U_i)$라고
쓴다. $p_i(z) \in U_i$이고, 따라서 $\xi_i \in U_i$이며, 따라서
$\xi \in U$임을 알 수 있다. 그러므로 $\xi$는 $Z$의 일반점이다.
$\xi$가 $Z$의 유일한 일반점이라는 증명은 생략한다(힌트: 두 번째 일반점은
$X_i$에서 $\xi_i$로 사상되어야 하므로 $\xi$와 같아야 함을 보이라. 이는
$X_i$의 스펙트럴성에 의해 기약집합 $Z_i$가 유일한 일반점을 갖기 때문이다).
이것으로 증명이 끝난다.
\end{proof}

\begin{lemma}
\label{topology-lemma-descend-opens}
$\mathcal{I}$를 쌍대여과 지표범주라 하자. $i \mapsto X_i$를 스펙트럴
공간들의 도표라 하고, $\mathcal{I}$의 $a : j \to i$에 대응하는 사상
$f_a : X_j \to X_i$가 스펙트럴이라고 하자. $X = \lim X_i$로 놓고
사영사상을 $p_i : X \to X_i$라 하자.
\begin{enumerate}
\item 임의의 준콤팩트 열린집합 $U \subset X$가 주어지면, 어떤
$i \in \Ob(\mathcal{I})$와 준콤팩트 열린집합 $U_i \subset X_i$가
존재하여 $p_i^{-1}(U_i) = U$가 된다.
\item 준콤팩트 열린집합들 $U_i \subset X_i$와 $U_j \subset X_j$가
주어지고 $p_i^{-1}(U_i) \subset p_j^{-1}(U_j)$라 하자. 그러면 어떤
$k \in \Ob(\mathcal{I})$와 사상들 $a : k \to i$ 및 $b : k \to j$가
존재하여 $f_a^{-1}(U_i) \subset f_b^{-1}(U_j)$가 된다.
\item $U_i, U_{1, i}, \ldots, U_{n, i} \subset X_i$가 준콤팩트
열린집합들이고
$p_i^{-1}(U_i) = p_i^{-1}(U_{1, i}) \cup \ldots \cup p_i^{-1}(U_{n, i})$
이면, $\mathcal{I}$의 어떤 사상 $a : j \to i$에 대하여
$f_a^{-1}(U_i) = f_a^{-1}(U_{1, i}) \cup \ldots \cup f_a^{-1}(U_{n, i})$
이다.
\item 교집합에 대해서도 (3)과 같은 명제가 성립한다.
\end{enumerate}
\end{lemma}

\begin{proof}
(1)은 Lemma \ref{topology-lemma-inverse-limit-spectral-spaces-constructible}의 특수한
경우이다. 준콤팩트 열린집합은 구성가능 위상에서 열리고 닫혀 있으므로,
(2)는 Lemma \ref{topology-lemma-inverse-limit-spectral-spaces-equal}의 특수한
경우이다. (3)과 (4)는 (1), (2), 그리고 부분집합의 역상을 취하는 연산이
합집합 및 교집합을 취하는 연산과 교환한다는 사실로부터 형식적으로 따라온다.
\end{proof}

\begin{lemma}
\label{topology-lemma-make-spectral-space}
$W$를 스펙트럴 공간 $X$의 부분집합이라 하자. 다음은 동치이다.
\begin{enumerate}
\item $W$는 구성가능 집합들의 교집합이고 일반화에 대해 닫혀 있다.
\item $W$는 준콤팩트이고 일반화에 대해 닫혀 있다.
\item 준콤팩트 부분집합 $E \subset X$가 존재하여, $W$는 $E$로
특수화되는 점들의 집합이다.
\item $W$는 준콤팩트 열린 부분집합들의 교집합이다.
\item
\label{topology-item-intersection-quasi-compact-open}
공집합이 아닌 집합 $I$와 준콤팩트 열린집합들 $U_i \subset X$,
$i \in I$가 존재하여 $W = \bigcap U_i$이고, 모든 $i, j \in I$에 대하여
$U_k \subset U_i \cap U_j$인 $k \in I$가 존재한다.
\end{enumerate}
이 경우 (a) $W$는 스펙트럴 공간이고, (b) 위상공간으로서
$W = \lim U_i$이며, (c) $W$를 포함하는 임의의 열린집합 $U$에 대하여
$U_i \subset U$인 $i$가 존재한다.
\end{lemma}

\begin{proof}
$W \subset X$가 (1)을 만족한다고 하자. 그러면 $W$는 구성가능 위상에서
닫혀 있고, 따라서 구성가능 위상에서 준콤팩트이며(Lemmas
\ref{topology-lemma-constructible-hausdorff-quasi-compact}와
\ref{topology-lemma-closed-in-quasi-compact}), 따라서 $X$의 위상에서도
준콤팩트이다($X$의 열린집합들이 구성가능 위상에서 열려 있기 때문이다).
그러므로 (2)가 성립한다.

\medskip\noindent
$E = W$로 놓으면 (2)가 (3)을 함의한다는 것은 분명하다.

\medskip\noindent
$X$를 스펙트럴 공간이라 하고 $E \subset W$가 (3)과 같다고 하자. $W$의
모든 점이 $E$의 한 점으로 특수화되므로 $E$를 포함하는 $W$의 열린집합은
$W$와 같다. 따라서 $E$가 준콤팩트이므로 $W$도 준콤팩트이다.
$x \in X$이고 $x \not \in W$이면 $Z = \overline{\{x\}}$는 $W$와
서로소이다. $W$가 준콤팩트이므로 $W \subset U$이고
$U \cap Z = \emptyset$인 준콤팩트 열린집합 $U$를 찾을 수 있다. 따라서
(4)가 성립한다고 결론 내린다.

\medskip\noindent
$W = \bigcap_{j \in J} U_j$이면, $I$를 $J$의 유한 부분집합들의 집합으로
놓고 $i = \{j_1, \ldots, j_r\}$에 대하여
$U_i = U_{j_1} \cap \ldots \cap U_{j_r}$로 놓음으로써 (4)가 (5)를
함의함을 알 수 있다. (5)가 (1)을 함의하는 것은 곧바로 알 수 있다.

\medskip\noindent
$I$와 $U_i$가 (5)와 같다고 하자. $W = \bigcap U_i$이므로 극한의
보편성질에 의해 $W = \lim U_i$이다. 그러면 Lemma
\ref{topology-lemma-directed-inverse-limit-spectral-spaces}에 의해 $W$는 스펙트럴
공간이다. $U \subset X$를 $W$의 열린근방이라 하자. 그러면
$E_i = U_i \cap (X \setminus U)$는 스펙트럴 공간 $Z = X \setminus U$의
구성가능 부분집합들의 족이고 그 교집합은 공집합이다.
% P01-E0025: the frozen proof calls the constructible topology spectral.
$Z$ 위의 스펙트럴 위상이 준콤팩트임을 이용하면(Lemma
\ref{topology-lemma-constructible-hausdorff-quasi-compact}), Lemma
\ref{topology-lemma-intersection-closed-in-quasi-compact}에 의해 어떤 $i$에 대하여
$E_i = \emptyset$이라고 결론 내린다.
\end{proof}

\begin{lemma}
\label{topology-lemma-make-spectral-space-minus}
$X$를 스펙트럴 공간이라 하자. $E \subset X$를 구성가능 부분집합이라 하자.
$W \subset X$를 $E$의 한 점으로 특수화되는 $X$의 점들의 집합이라 하자.
그러면 $W \setminus E$는 스펙트럴 공간이다. Lemma
\ref{topology-lemma-make-spectral-space} (\ref{topology-item-intersection-quasi-compact-open})에서와
같은 $U_i$들에 대하여 $W = \bigcap U_i$이면
$W \setminus E = \lim (U_i \setminus E)$이다.
\end{lemma}

\begin{proof}
$E$가 구성가능이므로 준콤팩트이고, 따라서 Lemma
\ref{topology-lemma-make-spectral-space}를 $W$에 적용할 수 있다. $E$가
구성가능이면 모든 $i \in I$에 대하여 $E$는 $U_i$에서 구성가능이다.
따라서 Lemma \ref{topology-lemma-spectral-sub}에 의해 $U_i \setminus E$는
스펙트럴이다. $W \setminus E = \bigcap (U_i \setminus E)$이므로,
% P01-E0026: the frozen limit expression below drops the grouping around U_i minus E.
극한의 보편성질에 의해 $W \setminus E = \lim U_i \setminus E$이다.
그러면 Lemma \ref{topology-lemma-directed-inverse-limit-spectral-spaces}에 의해
$W \setminus E$는 스펙트럴 공간이다.
\end{proof}

\section{스톤--체흐 콤팩트화}
\label{topology-section-stone-cech}

\noindent
위상공간 $X$의 스톤--체흐 콤팩트화는 $X$에서 하우스도르프
준콤팩트 공간 $\beta(X)$로 가는 사상 $X \to \beta(X)$로서, 이러한
사상들에 대하여 보편적인 것이다. 다음의 간단한 보조정리를 사용하는
표준적인 논증으로 그 존재를 증명한다.

\begin{lemma}
\label{topology-lemma-dense-image}
$f : X \to Y$를 위상공간들의 연속사상이라 하자. $f(X)$가 $Y$에서
조밀하고 $Y$가 하우스도르프라고 가정하자. 그러면 $Y$의 기수는, $P$를
멱집합 연산이라 할 때, $P(P(X))$의 기수 이하이다.
\end{lemma}

\begin{proof}
$S = f(X) \subset Y$로 놓자. $\mathcal{D}$를 $Y$의 모든 닫힌 영역,
즉 자기 내부의 폐포와 같은 부분집합 $D \subset Y$들의 집합이라 하자.
$Y$의 열린 부분집합의 폐포는 닫힌 영역임에 유의하라. $y \in Y$에
대하여 다음 집합을 생각하자.
$$
I_y = \{T \subset S \mid \text{ there exists }
D \in \mathcal{D}\text{ with }T = S \cap D\text{ and }y \in D\}.
$$
$S$가 $Y$에서 조밀하므로 모든 닫힌 영역 $D$에 대하여 $S \cap D$는
$D$에서 조밀하다. 따라서 $D, D' \in \mathcal{D}$에 대하여
$D \cap S = D' \cap S$이면 $D = D'$이다. 그러므로 $I_y = I_{y'}$이면
$y = y'$이다. 실제로 하우스도르프 조건에 의해 $y$를 포함하지만
$y'$은 포함하지 않는 닫힌 영역을 찾을 수 있다. 이로부터 결과가 따른다.
\end{proof}

\noindent
$X$를 위상공간이라 하자. Lemma \ref{topology-lemma-dense-image}에 의해, 상이
조밀하고 $Y$가 하우스도르프이면서 준콤팩트인 연속사상
$f : X \to Y$들의 동형류로 이루어진 집합 $I$가 존재한다. 각
$i \in I$에 대하여 대표원 $f_i : X \to Y_i$를 선택하자. 사상
$$
\prod f_i : X \longrightarrow \prod\nolimits_{i \in I} Y_i
$$
를 생각하고 그 상의 폐포를 $\beta(X)$로 나타내자. 각 $Y_i$가
하우스도르프이므로 $\beta(X)$도 하우스도르프이다. 각 $Y_i$가
준콤팩트이므로 $\beta(X)$도 준콤팩트이다(Theorem
\ref{topology-theorem-tychonov}와 Lemma \ref{topology-lemma-closed-in-quasi-compact}를
사용하라).

\medskip\noindent
표준사상 $X \to \beta(X)$가 하우스도르프 준콤팩트 공간으로 가는
사상들에 대한 보편성질을 만족함을 보이자. 즉, $f : X \to Y$를 그러한
사상이라 하자. $Z \subset Y$를 $f(X)$의 폐포라 하자. 그러면
$X \to Z$는 사상 $f_i : X \to Y_i$ 가운데 하나, 이를테면
$f_{i_0} : X \to Y_{i_0}$와 동형이다. 따라서 원하는 대로 $f$는
$X \to \beta(X) \to \prod Y_i \to Y_{i_0} \cong Z \to Y$로 분해된다.

\begin{lemma}
\label{topology-lemma-one-point-compactification}
$X$를 하우스도르프 국소 준콤팩트 공간이라 하자. $X$를 준콤팩트
하우스도르프 공간 $X^*$의 열린 부분공간과 동일시하며
$X^* \setminus X$가 한원소 집합이 되게 하는 사상 $X \to X^*$가 존재한다
(한 점 콤팩트화). 특히 사상 $X \to \beta(X)$는 $X$를 $\beta(X)$의
열린 부분공간과 동일시한다.
\end{lemma}

\begin{proof}
$X^* = X \amalg \{\infty\}$로 놓자. $X^*$의 부분집합 $V$가 열린다는
것을 다음과 같이 선언한다. $V \subset X$이고 V가 $X$에서 열려 있거나,
혹은 $\infty \in V$이고 $U = V \cap X$가 $X$의 열린집합이며
$X \setminus U$가 준콤팩트이다. 이것이 위상을 정의한다는 확인은
생략한다.
% P01-E0027: the frozen proof names X, rather than X-star, as the ambient open subspace.
$X \to X^*$가 $X$를 $X$의 열린 부분공간과 동일시함은 분명하다.

\medskip\noindent
$X$가 국소 준콤팩트이므로 모든 점 $x \in X$에는 준콤팩트 근방
$x \in E \subset X$가 존재한다. 그러면 $E$는 닫혀 있고(Lemma
\ref{topology-lemma-quasi-compact-in-Hausdorff}의 (1)),
$V = (X \setminus E) \amalg \{\infty\}$는 $E$의 내부와 서로소인
$\infty$의 열린근방이다. 따라서 $X^*$는 하우스도르프이다.

\medskip\noindent
$X^* = \bigcup V_i$를 열린덮개라 하자. 그러면 어떤 $i$, 이를테면
$i_0$에 대하여 $\infty \in V_{i_0}$이다. 구성에 의해
$Z = X^* \setminus V_{i_0}$는 준콤팩트이다. 따라서 덮개
$Z \subset \bigcup_{i \not = i_0} Z \cap V_i$에는 유한 세분이 존재하고,
이는 주어진 $X^*$의 덮개에도 유한 세분이 존재함을 뜻한다. 그러므로
$X^*$는 준콤팩트이다.

\medskip\noindent
스톤--체흐 콤팩트화의 보편성질에 의해 사상 $X \to X^*$는
$X \to \beta(X) \to X^*$로 분해된다. $\varphi : \beta(X) \to X^*$를
이 분해사상이라 하자. 그러면 $X \to \varphi^{-1}(X)$는
$\varphi^{-1}(X) \to X$의 절단이므로 그 상은 닫혀 있다(Lemma
\ref{topology-lemma-section-closed}). $X \to \beta(X)$의 상이 조밀하므로
$X = \varphi^{-1}(X)$라고 결론 내린다.
\end{proof}









\section{극단적인 비연결 공간}
\label{topology-section-extremally-disconnected}

\noindent
이 절의 내용은 \cite{Gleason}에서 가져왔다(\cite{Rainwater}에서와 같이
약간 수정하였다). Gleason의 논문에서는 준콤팩트 하우스도르프 공간들의
범주에서 ``사영 대상''이 정확히 극단적인 비연결 공간임을 보인다.

\begin{definition}
\label{topology-definition-extremally-disconnected}
위상공간 $X$의 모든 열린 부분집합의 폐포가 열려 있으면 $X$를
{\it 극단적으로 비연결}이라고 한다.
\end{definition}

\noindent
$X$가 하우스도르프이고 극단적으로 비연결이면 $X$는 완전불연결이다
(일반적으로는 참이 아니다). $X$가 준콤팩트이고 하우스도르프이며
극단적으로 비연결이면 Lemma \ref{topology-lemma-profinite}에 의해 $X$는
프로유한이지만, 그 역은 일반적으로 성립하지 않는다. 예를 들어 $p$진 정수
$\mathbf{Z}_p = \lim \mathbf{Z}/p^n\mathbf{Z}$는 프로유한 공간이지만
극단적으로 비연결이지 않다. 실제로 $U \subset \mathbf{Z}_p$를 부치가
짝수인 영이 아닌 원소들의 집합이라 하면 $U$는 열려 있지만, 그 폐포
$U \cup \{0\}$은 열려 있지 않다.

\begin{lemma}
\label{topology-lemma-image-open-technical}
$f : X \to Y$를 위상공간들의 연속사상이라 하자. $f$가 전사이고 모든
진닫힌 부분집합 $E \subset X$에 대하여 $f(E) \not = Y$라고 가정하자.
그러면 열린집합 $U \subset X$에 대하여 부분집합 $f(U)$는
$Y \setminus f(X \setminus U)$의 폐포에 포함된다.
\end{lemma}

\begin{proof}
$y \in f(U)$를 택하고 $V \subset Y$를 $y$의 임의의 열린근방이라 하자.
$V$가 $Y \setminus f(X \setminus U)$와 만남을 보이겠다.
$W = U \cap f^{-1}(V)$는 $X$의 공집합이 아닌 열린 부분집합이므로
$f(X \setminus W) \not = Y$임에 유의하라. $y' \in Y$이고
$y' \not \in f(X \setminus W)$인 y'을 택하자. 그러면
$y' \in V$이고 $y' \in Y \setminus f(X \setminus U)$임을 초등적으로
확인할 수 있다.
\end{proof}

\begin{lemma}
\label{topology-lemma-intersection-empty}
$X$를 극단적으로 비연결인 공간이라 하자. $U, V \subset X$가 서로소인
열린 부분집합들이면 $\overline{U}$와 $\overline{V}$도 서로소이다.
\end{lemma}

\begin{proof}
가정에 의해 $\overline{U}$는 열려 있으므로 $V \cap \overline{U}$는
열려 있고 $U$와 서로소이다. 그런데 $\overline{U}$는 $U$를 포함하는
$X$의 모든 닫힌 부분집합의 교집합이므로 이는 공집합이다. 이는 열린집합
$\overline{V} \cap \overline{U}$가 $V$를 피함을 뜻하므로, 같은 논증에
의해 이 집합도 공집합이다.
\end{proof}

\begin{lemma}
\label{topology-lemma-isomorphism}
$f : X \to Y$를 하우스도르프 준콤팩트 위상공간들의 연속사상이라 하자.
$Y$가 극단적으로 비연결이고 $f$가 전사이며, $X$의 모든 진닫힌 부분집합
$Z$에 대하여 $f(Z) \not = Y$이면 $f$는 위상동형사상이다.
\end{lemma}

\begin{proof}
Lemma \ref{topology-lemma-bijective-map}에 의해 $f$가 단사임을 보이면 충분하다.
$x, x' \in X$가 서로 다른 점들이고 $y = f(x) = f(x')$이라고 가정하자.
$x, x'$의 서로소인 열린근방들 $U, U' \subset X$를 택하자. $f$가
닫힌사상임에 유의하라(Lemma \ref{topology-lemma-closed-map}). 따라서
$T = f(X \setminus U)$와 $T' = f(X \setminus U')$는 $Y$에서 닫혀 있다.
$X$가 $X \setminus U$와 $X \setminus U'$의 합집합이므로
$Y = T \cup T'$이다. Lemma \ref{topology-lemma-image-open-technical}에 의해
$y$는 $Y \setminus T$의 폐포와 $Y \setminus T'$의 폐포에 모두
포함된다. 한편 Lemma \ref{topology-lemma-intersection-empty}에 의해 이 교집합은
공집합이다. 이와 같이 원하는 모순을 얻는다.
\end{proof}

\begin{lemma}
\label{topology-lemma-find-compact-subset}
$f : X \to Y$를 하우스도르프 준콤팩트 위상공간들의 연속 전사라 하자.
$f(E) = Y$이고 모든 진닫힌 부분집합 $E' \subset E$에 대하여
$f(E') \not = Y$인 준콤팩트 부분집합 $E \subset X$가 존재한다.
\end{lemma}

\begin{proof}
$X$의 준콤팩트 부분집합들이 정확히 닫힌 부분집합들임을(Lemma
\ref{topology-lemma-closed-in-compact}) 이제부터 별도로 언급하지 않고 사용하겠다.
$f(E) = Y$인 모든 준콤팩트 부분집합 $E \subset X$의 모임
$\mathcal{E}$를 포함관계로 순서화하자. Zorn의 보조정리를 사용하여
$\mathcal{E}$에 극소원소가 있음을 보이겠다. 이를 위해서는
$\mathcal{E}$의 원소들로 이루어진 전순서족 $E_\lambda$가 주어질 때
교집합 $\bigcap E_\lambda$가 $\mathcal{E}$의 원소임을 보이면 충분하다.
이는 닫혀 있으므로 준콤팩트이다. 모든 $y \in Y$에 대하여 집합들
$E_\lambda \cap f^{-1}(\{y\})$는 공집합이 아니고 닫혀 있으므로,
Lemma \ref{topology-lemma-intersection-closed-in-quasi-compact}에 의해 교집합
$\bigcap E_\lambda \cap f^{-1}(\{y\}) = \bigcap (E_\lambda \cap f^{-1}(\{y\}))$
은 공집합이 아니다. 이것으로 증명이 끝난다.
\end{proof}

\begin{proposition}
\label{topology-proposition-projective-in-category-hausdorff-qc}
$X$를 하우스도르프 준콤팩트 위상공간이라 하자. 다음 조건들은 서로
동치이다.
\begin{enumerate}
\item $X$는 극단적으로 비연결이다.
\item $Y$가 하우스도르프 준콤팩트인 임의의 연속 전사
$f : Y \to X$에 대하여 연속인 절단이 존재한다.
\item 준콤팩트 하우스도르프 공간들의 연속사상들로 이루어진 임의의
실선 가환도식
$$
\xymatrix{
& Y \ar[d] \\
X \ar@{..>}[ru] \ar[r] & Z
}
$$
에서 $Y \to Z$가 전사이면, 그 도식을 가환이 되게 하는 점선 화살표가
위상공간의 범주에 존재한다.
\end{enumerate}
\end{proposition}

\begin{proof}
(3)이 (2)를 함의함은 분명하다. 반대로 (2)가 성립하고 $X \to Z$와
$Y \to Z$가 (3)과 같다고 하자. 그러면 (2)에 의해 사영
$X \times_Z Y \to X$의 절단이 존재하고, 이는 적절한 점선 화살표가
존재함을 함의한다(세부사항은 생략한다). 따라서 (3)과 (2)는 동치이다.

\medskip\noindent
$X$가 극단적으로 비연결이라고 가정하고 $f : Y \to X$를 (2)와 같다고
하자. Lemma \ref{topology-lemma-find-compact-subset}에 의해 $f(E) = X$이지만 모든
진닫힌 부분집합 $E' \subset E$에 대하여 $f(E') \not = X$인 준콤팩트
부분집합 $E \subset Y$가 존재한다. Lemma \ref{topology-lemma-isomorphism}에 의해
$f|_E : E \to X$가 위상동형사상이고, 그 역이 원하는 절단을 준다.

\medskip\noindent
(2)를 가정하자. $U \subset X$를 열린집합이라 하고 그 여집합을 $Z$라
하자. 연속 전사 $f : \overline{U} \amalg Z \to X$를 생각하자.
$\sigma$를 절단이라 하자. 그러면
$\overline{U} = \sigma^{-1}(\overline{U})$는 열려 있다. 따라서 $X$는
극단적으로 비연결이다.
\end{proof}

\begin{lemma}
\label{topology-lemma-rainwater}
$f : X \to X$를 하우스도르프 위상공간의 연속 전사 자기 사상이라 하자.
$f$가 $\text{id}_X$가 아니면 $X = E \cup f(E)$인 진닫힌 부분집합
$E \subset X$가 존재한다.
\end{lemma}

\begin{proof}
$f(p) \not = p$인 $p \in X$를 택하자. 서로소인 열린근방
$p \in U$, $f(p) \in V$를 택하고
% P01-E0028: the frozen formula omits grouping around U intersect f^{-1}(V).
$E = X \setminus U \cap f^{-1}(V)$로 놓자. 그러면 $p \not \in E$이므로
$E$는 진닫힌 부분집합이다. $x \in X$이면 $x \in E$이거나, 그렇지 않으면
$x \in U \cap f^{-1}(V)$이다. $x = f(y)$로 쓰자($f$가 전사이므로
가능하다). 만일 $y \in U \cap f^{-1}(V)$이면
$x = f(y) \in V$이지만 이는 $x \in U$와 모순이다. 따라서
$y \in E$이고 $x \in f(E)$이다.
\end{proof}

\begin{example}
\label{topology-example-stone-Cech-discrete}
Proposition \ref{topology-proposition-projective-in-category-hausdorff-qc}을 사용하면
이산공간 $X$의 스톤--체흐 콤팩트화 $\beta(X)$가 극단적으로 비연결임을
알 수 있다. 실제로 $Y$가 준콤팩트이고 하우스도르프일 때
$f : Y \to \beta(X)$를 연속 전사라 하자. 그러면 $X$가 이산이므로 사상
$X \to \beta(X)$를 연속인(!) 사상 $X \to Y$로 들어 올릴 수 있다.
스톤--체흐 콤팩트화의 보편성질에 의해 분해
$X \to \beta(X) \to Y$를 얻는다. $\beta(X) \to Y \to \beta(X)$가
조밀한 부분집합 $X$ 위에서 항등사상과 같으므로 절단을 얻었다고 결론
내린다. 특히 이산공간의 스톤--체흐 콤팩트화는 완전불연결이고, 따라서
프로유한이라고 결론 내린다(Definition
\ref{topology-definition-extremally-disconnected} 다음의 논의와 Lemma
\ref{topology-lemma-profinite}를 보라).
\end{example}

\noindent
Example \ref{topology-example-stone-Cech-discrete}가 제공하는 극단적으로 비연결인
공간들을 사용하여, 모든 준콤팩트 하우스도르프 공간이 준콤팩트
하우스도르프 공간들의 범주에서 ``사영 덮개''를 가짐을 증명할 수 있다.

\begin{lemma}
\label{topology-lemma-existence-projective-cover}
\begin{slogan}
모든 준콤팩트 하우스도르프 공간에는 표준적인 극단적 비연결 덮개가 있다
\end{slogan}
$X$를 준콤팩트 하우스도르프 공간이라 하자. $X'$가 준콤팩트이고
하우스도르프이며 극단적으로 비연결이 되게 하는 연속 전사
$X' \to X$가 존재한다. $X'$의 모든 진닫힌 부분집합이 $X$ 위로 사상되지
않는다고 요구하면 $X'$는 동형을 제외하고 유일하다.
\end{lemma}

\begin{proof}
$Y = X$로 놓되 Y에는 이산위상을 주자. $X' = \beta(Y)$로 놓자.
연속사상 $Y \to X$는 $Y \to X' \to X$로 분해된다. Example
\ref{topology-example-stone-Cech-discrete}에 의해 이것은 보조정리의 첫 번째 주장을
증명한다.

\medskip\noindent
Lemma \ref{topology-lemma-find-compact-subset}에 의해 $X$ 위로 전사이면서
$E$의 어떤 진닫힌 부분집합도 $X$ 위로 전사이지 않는 준콤팩트 부분집합
$E \subset X'$를 찾을 수 있다. $X'$가 극단적으로 비연결이므로 $X$ 위의
연속사상 $f : X' \to E$가 존재한다(Proposition
\ref{topology-proposition-projective-in-category-hausdorff-qc}). $f$와 사상
$E \to X'$를 합성하면 연속 자기 사상 $f|_E : E \to E$를 얻는다.
$f|_E$가 전사가 아니면 그 상이 $X$ 위로 전사인 진닫힌 부분집합이 될
것이므로, 이 자기 사상은 전사이어야 한다. 또한 이것이 항등사상이
아니면 Lemma \ref{topology-lemma-rainwater}에 의해 $E$가 극소가 아니므로,
$f|_E$는 $\text{id}_E$이어야 한다. 따라서 $\text{id}_E$는 극단적으로
비연결인 공간 $X'$를 경유하여 분해된다. 형식적인 범주론적 논증
(Proposition \ref{topology-proposition-projective-in-category-hausdorff-qc}의
특징화를 사용한다)에 의해 $E$가 극단적으로 비연결임을 알 수 있다.

\medskip\noindent
% P01-E0029: the frozen uniqueness paragraph invokes minimality of X' although only E was made minimal.
유일성을 증명하기 위해 두 번째 극소 덮개 $X'' \to X$가 있다고 하자.
Proposition \ref{topology-proposition-projective-in-category-hausdorff-qc}에서 증명한
들어올림 성질에 의해 $X$ 위의 연속사상 $g : X' \to X''$를 찾을 수 있다.
$g$는 닫힌사상임에 유의하라(Lemma \ref{topology-lemma-closed-map}). 따라서
$g(X') \subset X''$는 $X$ 위로 전사인 닫힌 부분집합이고, $X''$의
극소성에 의해 $g(X') = X''$라고 결론 내린다. 한편
$E \subset X'$가 진닫힌 부분집합이면 $X'$의 극소성에 의해 $E$는
$X$ 위로 사상되지 않으므로 $g(E) \not = X''$이다. Lemma
\ref{topology-lemma-isomorphism}에 의해 $g$는 동형사상이다.
\end{proof}

\begin{remark}
\label{topology-remark-size-projective-cover}
$X$를 준콤팩트 하우스도르프 공간이라 하자. $\kappa$를 $X$의 기수보다
크거나 같은 무한기수라 하자. 그러면 극소 준콤팩트 하우스도르프
극단적 비연결 덮개 $X' \to X$(Lemma
\ref{topology-lemma-existence-projective-cover})의 기수는 $2^{2^\kappa}$ 이하이다.
실제로 $X$로 전단사로 사상되는 부분집합 $S \subset X'$를 택하자.
$X'$의 극소성에 의해 집합 $S$는 $X'$에서 조밀하다. 따라서 Lemma
\ref{topology-lemma-dense-image}에 의해 $|X'| \leq 2^{2^\kappa}$이다.
\end{remark}






\section{기타}
\label{topology-section-miscellany}


\noindent
다음 보조정리는 준분리 스킴에 수반되는 바탕 위상공간에 적용된다.

\begin{lemma}
\label{topology-lemma-topology-quasi-separated-scheme}
$X$를 다음 조건을 만족하는 위상공간이라 하자.
\begin{enumerate}
\item 준콤팩트 열린집합들로 이루어진 위상의 기저가 존재한다.
\item 임의의 두 준콤팩트 열린집합의 교집합이 준콤팩트이다.
\end{enumerate}
그러면 다음이 성립한다.
\begin{enumerate}
\item $X$는 국소 준콤팩트이다.
\item 준콤팩트 열린집합 $U \subset X$는 레트로콤팩트이다.
\item 임의의 준콤팩트 열린집합 $U \subset X$에는, $J$가 유한하고 모든
$U_j$ 및 $U_j \cap U_{j'}$가 준콤팩트인 열린덮개
$\mathcal{U} : U = \bigcup_{j\in J} U_j$들의 공종계가 존재한다.
% P01-E0030: the frozen reader-facing list contains the editorial placeholder below.
\item 여기에 더 추가하라.
\end{enumerate}
\end{lemma}

\begin{proof}
생략한다.
\end{proof}

\begin{definition}
\label{topology-definition-isolated-point}
$X$를 위상공간이라 하자. $\{x\}$가 $X$에서 열려 있으면 $x \in X$를
$X$의 {\it 고립점}이라고 한다.
\end{definition}



\section{분할과 층화}
\label{topology-section-stratifications}

\noindent
층화는 여러 가지 서로 다른 방식으로 정의할 수 있다. 이 절에서 선택한
정의들에 관한 의견을 환영한다.

\begin{definition}
\label{topology-definition-paritition}
$X$를 위상공간이라 하자. $X$의 {\it 분할}은 국소 닫힌 부분집합들
$X_i$로의 분해 $X = \coprod X_i$이다. $X_i$들을 이 분할의
{\it 조각}이라고 한다.
% P01-E1229: the frozen wording makes the direction of refinement ambiguous/reversed.
$X$의 두 분할이 주어졌을 때, 한쪽의 조각들이 다른 쪽의 조각들의
합집합이면 전자가 후자를 {\it 세분한다}고 한다.
\end{definition}

\noindent
모든 위상공간 $X$에는 연결성분들로 이루어진 분할이 있다. $X$의
기약성분들이 유한개 $Z_1, \ldots, Z_r$라면, 부분집합
$I \subset \{1, \ldots, r\}$를 지표로 하는 조각
$X_I = \bigcap_{i \in I} Z_i \setminus (\bigcup_{i \not \in I} Z_i)$들로
이루어진 분할이 존재하며, 이는 연결성분들로 이루어진 분할을 세분한다.

\begin{definition}
\label{topology-definition-good-stratification}
$X$를 위상공간이라 하자. $X$의 {\it 좋은 층화}는 모든
$i, j \in I$에 대하여
$$
X_i \cap \overline{X_j} \not = \emptyset
\Rightarrow
X_i \subset \overline{X_j}.
$$
를 만족하는 분할 $X = \coprod X_i$이다.
\end{definition}

\noindent
좋은 층화 $X = \coprod_{i \in I} X_i$가 주어지면,
$X_i \subset \overline{X_j}$일 필요충분조건이 $i \leq j$인 것으로
정하여 $I$ 위에 부분순서를 얻는다. 그러면
$$
\overline{X_j} = \bigcup\nolimits_{i \leq j} X_i
$$
임을 알 수 있다. 하지만 대수기하학에서 흔히 일어나는 일은 마지막
표시식의 왼쪽이 오른쪽의 부분집합이라는 것뿐이다. 이로부터 다음 정의가
나온다.

\begin{definition}
\label{topology-definition-stratification}
$X$를 위상공간이라 하자. $X$의 {\it 층화}는 분할
$X = \coprod_{i \in I} X_i$와 $I$ 위의 부분순서로 주어지며, 모든
$j \in I$에 대하여
$$
\overline{X_j} \subset \bigcup\nolimits_{i \leq j} X_i
$$
를 만족한다. 조각들 $X_i$를 이 층화의 {\it 층}이라고 한다.
\end{definition}

\noindent
흔히 층화에 추가 조건을 부과한다. 예를 들어 층화가
{\it 국소 유한}이면 특히 다루기 좋다. 이는 모든 점에 유한개의 층과만
만나는 근방이 있다는 뜻이다. 더 일반적으로 다음 정의를 도입한다.

\begin{definition}
\label{topology-definition-locally-finite}
$X$를 위상공간이라 하자. $I$를 집합이라 하고 각 $i \in I$에 대하여
$E_i \subset X$를 부분집합이라 하자. 모든 $x \in X$에 대하여
$\{i \in I | E_i \cap U \not = \emptyset\}$가 유한이 되게 하는 $x$의
열린근방 $U$가 존재하면 모임 $\{E_i\}_{i \in I}$를
{\it 국소 유한}이라고 한다.
\end{definition}


\begin{remark}
\label{topology-remark-locally-finite-stratification}
위상공간 $X$의 국소 유한 층화 $X = \coprod X_i$가 주어지면, $I$로
지표화된 $X$의 닫힌 부분집합들의 족 $Z_i = \bigcup_{j \leq i} X_j$를
얻고, 이들은
$$
Z_i \cap Z_j = \bigcup\nolimits_{k \leq i, j} Z_k
$$
를 만족한다. 반대로 부분순서집합 $I$로 지표화된 닫힌 부분집합들
$Z_i \subset X$가 주어져서 $X = \bigcup Z_i$이고, 모든 점에 유한개의
$Z_i$와만 만나는 근방이 있으며, 위 표시식이 성립한다고 하자. 그러면
$X_i = Z_i \setminus \bigcup_{j < i} Z_j$로 놓아 $X$의 국소 유한
층화를 얻는다.
\end{remark}

\begin{lemma}
\label{topology-lemma-partition-refined-by-stratification}
$X$를 위상공간이라 하자. $X = \coprod X_i$를 $X$의 유한 분할이라 하자.
그러면 이를 세분하는 $X$의 유한 층화가 존재한다.
\end{lemma}

\begin{proof}
$T_i = \overline{X_i}$ 및 $\Delta_i = T_i \setminus X_i$로 놓자.
$S$를 $T_i$들과 $\Delta_i$들의 모든 교집합으로 이루어진 집합이라 하자.
(예를 들어 $T_1 \cap T_2 \cap \Delta_4$는 $S$의 원소이다.) 그러면
$S = \{Z_s\}$는 $X$의 닫힌 부분집합들의 유한 모임이며, 모든
$s, s' \in S$에 대하여 $Z_s \cap Z_{s'} \in S$이다. $S$ 위에
포함관계로 부분순서를 정의하자. 이제
$Y_s = Z_s \setminus \bigcup_{s' < s} Z_{s'}$로 놓으면 원하는 층화를
얻는다.
\end{proof}

\begin{lemma}
\label{topology-lemma-constructible-partition-refined-by-stratification}
$X$를 위상공간이라 하자. $X = T_1 \cup \ldots \cup T_n$이 구성가능
부분집합들의 합집합으로 표현되었다고 하자. 각 $X_i$가 구성가능이고 각
$T_k$가 층들의 합집합이 되게 하는 유한 층화 $X = \coprod X_i$가
존재한다.
\end{lemma}

\begin{proof}
구성가능 부분집합의 정의에 의해 각 $T_i$를 $U, V \subset X$가
레트로콤팩트 열린집합인 $U \cap V^c$들의 유한 합집합으로 쓸 수 있다.
따라서 $U_i, V_i \subset X$가 레트로콤팩트 열린집합인
$T_i = U_i \cap V_i^c$의 경우만 가정해도 된다. $S$를
$\emptyset, X, U_i^c, V_i^c$ 및 이들의 유한 교집합들로 이루어진
$X$의 닫힌 부분집합들의 유한집합이라 하자. $Z \in S$이면 $Z$는
$X$에서 구성가능이다(Lemma \ref{topology-lemma-constructible}). 또한 모든
$Z, Z' \in S$에 대하여 $Z \cap Z' \in S$이다. $S$ 위에 포함관계로
부분순서를 정의하자. $Z \in S$에 대하여
$X_Z = Z \setminus \bigcup_{Z' < Z,\ Z' \in S} Z'$로 놓으면 보조정리에
명시된 성질들을 만족하는 층화 $X = \coprod_{Z \in S} X_Z$를 얻는다.
\end{proof}

\begin{lemma}
\label{topology-lemma-noetherian-partition-refined-by-stratification}
$X$를 뇌터 위상공간이라 하자. $X$의 임의의 유한 분할은 유한 좋은
층화로 세분될 수 있다.
\end{lemma}

\begin{proof}
$X = \coprod X_i$를 $X$의 유한 분할이라 하자. $Z$를 $X$의
기약성분이라 하자. 유한한 지표집합에 대하여
$X = \bigcup \overline{X_i}$이므로 $Z \subset \overline{X_i}$인 $i$가
존재한다. $X_i$가 국소 닫혀 있으므로 이는 $Z \cap X_i$가 $Z$의
열린집합을 포함함을 뜻한다. 따라서 $Z \cap X_i$는 $X$의 열린집합
$U$를 포함한다(Lemma \ref{topology-lemma-Noetherian}).
$X_i = U \amalg X_i^1 \amalg X_i^2$로 쓰되
$X_i^1 = (X_i \setminus U) \cap \overline{U}$이고
$X_i^2 = (X_i \setminus U) \cap \overline{U}^c$로 놓자.
$i' \not = i$에 대해서는
$X_{i'}^1 = X_{i'} \cap \overline{U}$ 및
$X_{i'}^2 = X_{i'} \cap \overline{U}^c$로 놓는다. 그러면
$$
X \setminus U = \coprod X^k_l
$$
는 $\overline{U} \setminus U = \bigcup X_l^1$인 분할이다.
$X \setminus U$는 닫혀 있고 $X$보다 진정으로 작음에 유의하라. 뇌터
귀납법에 의해 이 분할을 유한 좋은 층화
$X \setminus U = \coprod_{\alpha \in A} T_\alpha$로 세분할 수 있다.
그러면 $X = U \amalg \coprod_{\alpha \in A} T_\alpha$는 처음의 분할을
세분하는 $X$의 유한 좋은 층화이다.
\end{proof}







\section{공간들의 여극한}
\label{topology-section-colimits}

\noindent
위상공간의 범주는 쌍대곱을 가진다. 구체적으로, $I$가 집합이고 각
$i \in I$에 대하여 위상공간 $X_i$가 주어졌다면 집합
$\coprod_{i \in I} X_i$에 {\it 쌍대곱 위상}을 준다. 이 위상의 기저로는
$U_i \subset X_i$가 열린집합인 꼴의 집합 $U_i$들을 사용한다.

\medskip\noindent
위상공간의 범주는 쌍대등화자를 가진다. 구체적으로
$a, b : X \to Y$가 위상공간들의 사상이면, $a$와 $b$의 쌍대등화자는
집합의 범주에서의 쌍대등화자 $Y/\sim$에 몫위상을 준 것이다(Section
\ref{topology-section-submersive}).

\begin{lemma}
\label{topology-lemma-colimits}
위상공간의 범주는 여극한을 가지며 집합으로 가는 망각 함자는 이들과
교환한다.
\end{lemma}

\begin{proof}
이는 위의 논의와 Categories, Lemma
\ref{categories-lemma-colimits-coproducts-coequalizers}에서 따른다.
여극한의 존재에 대한 또 다른 증명의 개요는 Categories, Remark
\ref{categories-remark-how-to-use-it}에 있다. 위의 논의로부터 망각 함자가
여극한과 교환함이 따른다. 이를 보는 다른 방법은 Categories, Lemma
\ref{categories-lemma-adjoint-exact}를 사용하고, 망각 함자가 집합에
대응하는 무차별(또는 비이산) 위상공간을 대응시키는 함자를 오른쪽
수반함자로 가짐을 이용하는 것이다.
\end{proof}





\section{위상군, 위상환, 위상가군}
\label{topology-section-topological-groups}

\noindent
이 절은 정의와 초등적인 성질들만을 간단히 다룬다.

\begin{definition}
\label{topology-definition-topological-group}
{\it 위상군}은 곱셈 $G \times G \to G$, $(x, y) \mapsto xy$와 역원사상
$G \to G$, $x \mapsto x^{-1}$이 연속이 되게 하는 위상을 갖춘 군 $G$이다.
{\it 위상군의 준동형}은 연속인 군 준동형이다.
\end{definition}

\noindent
$G$가 위상군이고 $H \subset G$가 부분군이면 유도위상을 갖춘 $H$는
위상군이다. $G$가 위상군이고 $G \to H$가 군의 전사이면 몫위상을 갖춘
$H$는 위상군이다.

\begin{example}
\label{topology-example-automorphisms-of-a-set}
$E$를 집합이라 하자. 자기 사상들의 집합 $\text{Map}(E, E)$에
콤팩트-열린 위상, 즉 $f : E \to E$가 주어졌을 때 $f$의 근방기본계가
유한 부분집합 $S \subset E$에 대한 집합들
$U_S(f) = \{f' : E \to E \mid f'|_S = f|_S\}$로 주어지는 위상을 줄 수
있다. $E$에 이산위상을 주면 이 위상에 대하여 작용
$$
\text{Map}(E, E) \times E \longrightarrow E,\quad
(f, e) \longmapsto f(e)
$$
은 연속이다. $X$가 위상공간이고 $X \times E \to E$가 연속사상이면
사상 $X \to \text{Map}(E, E)$도 연속이다. 달리 말해 콤팩트-열린 위상은
위에 표시된 ``작용'' 사상이 연속이 되게 하는 가장 엉성한 위상이다.
합성사상
$$
\text{Map}(E, E) \times \text{Map}(E, E) \to \text{Map}(E, E)
$$
도 연속이다(위의 근방에 대한 기술을 사용하면 쉽게 확인할 수 있다).
마지막으로 $\text{Aut}(E) \subset \text{Map}(E, E)$가 가역사상들의
부분집합이면 역원사상 $i : \text{Aut}(E) \to \text{Aut}(E)$,
$f \mapsto f^{-1}$도 연속이다. 실제로 $S \subset E$가 유한이라고 하면
$i^{-1}(U_S(f^{-1})) = U_{f^{-1}(S)}(f)$이다. 따라서 $\text{Aut}(E)$는
Definition \ref{topology-definition-topological-group}의 의미에서 위상군이다.
\end{example}

\begin{lemma}
\label{topology-lemma-topological-group-limits}
위상군의 범주는 극한을 가지며, 극한은 (a) 위상공간의 범주와 (b) 군의
범주로 가는 망각 함자들과 교환한다.
\end{lemma}

\begin{proof}
곱과 등화자에 대하여 존재성과 교환성을 증명하면 충분하다. Categories,
Lemma \ref{categories-lemma-limits-products-equalizers}를 보라.
$G_i$, $i \in I$를 위상군들의 모임이라 하자. 통상적인 곱
$G = \prod G_i$에 곱위상을 주자. 위상공간으로서
$G \times G = \prod (G_i \times G_i)$이므로(어느 범주에서나 곱은 곱과
교환한다) $G$ 위의 곱셈이 연속임을 알 수 있다. 역원사상도 마찬가지이다.
$a, b : G \to H$를 위상군들의 두 준동형이라 하자. 그러면 등화자로는
위상공간의 사상들로서 $a$와 $b$의 등화자를 그대로 택할 수 있고, 이는
유도위상을 갖춘 군의 사상들로서의 등화자와 같다.
\end{proof}

\begin{lemma}
\label{topology-lemma-profinite-group}
$G$를 위상군이라 하자. 다음 조건들은 서로 동치이다.
\begin{enumerate}
\item 위상공간으로서 $G$는 프로유한이다.
\item $G$는 유한 이산 위상군들로 이루어진 도식의 극한이다.
\item $G$는 유한 이산 위상군들의 쌍대여과 극한이다.
\end{enumerate}
\end{lemma}

\begin{proof}
위상공간들에 대해서는 이에 해당하는 결과를 이미 얻었다. Lemma
\ref{topology-lemma-profinite}를 보라. 이를 Lemma
\ref{topology-lemma-topological-group-limits}와 결합하면 (1)이 (3)을 함의함을
증명하는 것으로 충분하다.

\medskip\noindent
먼저 항등원 $e$의 모든 근방 $E$가 열린 부분군을 포함함을 증명하자.
$G$가 유한 이산 위상공간들의 쌍대여과 극한이므로(Lemma
\ref{topology-lemma-profinite}), 유한 이산공간 $T$로 가는 연속사상
$f : G \to T$ 중 $f^{-1}(f(\{e\})) \subset E$를 만족하는 것을 선택할 수
있다. 다음을 생각하자.
$$
H = \{g \in G \mid f(gg') = f(g')\text{ for all }g' \in G\}
$$
이는 $G$의 부분군이고 $E$에 포함된다. 따라서 $H$가 열려 있음을 보이면
충분하다.
% P01-E0032: the frozen construction does not define U_t when t is outside f(G).
$t \in T$를 택하고 $W = f^{-1}(\{t\})$로 놓자. $W \subset G$는 열리고
닫혀 있으며, 특히 준콤팩트이다. 각 $w \in W$에 대하여
$U_wU'_w \subset W$인 열린근방들 $e \in U_w \subset G$와
$w \in U'_w \subset W$가 존재한다. 준콤팩트성에 의해
$W = \bigcup U'_{w_i}$인 $w_1, \ldots, w_n$을 찾을 수 있다. 그러면
% P01-E0033: the frozen conclusion leaves g unquantified.
$U_t = U_{w_1} \cap \ldots \cap U_{w_n}$은 $e$의 열린근방이고, 모든
$w \in W$에 대하여 $f(gw) = t$이다. $T$가 유한이므로
$\bigcap_{t \in T} U_t \subset H$는 $e$의 열린근방이다.
$H \subset G$가 부분군이므로 $H$는 열려 있다.

\medskip\noindent
$H \subset G$를 열린 부분군이라 하자. $G$가 준콤팩트이므로 $G$에서
$H$의 지수는 유한이다. 이를테면
% P01-E0034: the frozen left-coset representatives do not justify the written conjugates.
$G = Hg_1 \cup \ldots \cup Hg_n$이라 하자. 그러면
$N = \bigcap_{i = 1, \ldots, n} g_iHg_i^{-1}$은 $H$에 포함되는 열린
정규부분군이다. $N$도 유한 지수를 가지므로 $G \to G/N$은 유한 이산
위상군으로 가는 전사이다.

\medskip\noindent
사상
$$
G \longrightarrow \lim_{N \subset G\text{ open and normal}} G/N
$$
을 생각하자. 이 사상은 위상군의 동형이라고 주장한다. 오른쪽의 극한은
쌍대여과되어 있으므로(두 열린 정규부분군의 교집합도 열리고 정규이다),
이것으로 보조정리의 증명이 끝난다. 각 $G \to G/N$이 연속이므로 이 사상은
연속이다. $G$가 하우스도르프이고 위의 논증에 의해 $e$의 모든 근방이
어떤 $N$을 포함하므로 이 사상은 단사이다. Lemma
\ref{topology-lemma-intersection-closed-in-quasi-compact}에 의해 이 사상은
전사이다. Lemma \ref{topology-lemma-bijective-map}에 의해 이 사상은
위상동형사상이다.
\end{proof}

\begin{definition}
\label{topology-definition-profinite-group}
Lemma \ref{topology-lemma-profinite-group}의 동치인 조건들을 만족하는 위상군을
{\it 프로유한군}이라고 한다.
\end{definition}

\noindent
$G_1 \to G_2 \to G_3 \to \ldots$가 위상군들의 계이면, 위상군으로서의
여극한 $G = \colim G_n$(Lemma \ref{topology-lemma-topological-group-colimits})은
일반적으로 위상공간으로서의 여극한(Lemma \ref{topology-lemma-colimits})과 다르다.
비록 두 여극한의 바탕집합은 같지만 그렇다. Examples, Section
\ref{examples-section-colimit-topology}를 보라.

\begin{lemma}
\label{topology-lemma-topological-group-colimits}
위상군의 범주는 여극한을 가지며, 여극한은 군의 범주로 가는 망각 함자와
교환한다.
\end{lemma}

\begin{proof}
여극한의 존재를 증명하기 위해 Categories, Remark
\ref{categories-remark-how-to-use-it}의 논증을 사용하겠다. 즉,
% P01-E0031: the frozen display has Top although the prose and objects require TopGroup.
$\mathcal{I} \to \textit{Top}$, $i \mapsto G_i$가 위상군의 범주
$\textit{TopGroup}$으로 가는 함자라고 가정하자. 그러면
$$
F : \textit{TopGroup} \longrightarrow \textit{Sets},\quad
H \longmapsto \lim_\mathcal{I} \Mor_{\textit{TopGroup}}(G_i, H)
$$
를 생각할 수 있다. 이 함자는 극한과 교환한다. 더욱이 임의의 위상군 $H$와
$F(H)$의 원소 $(\varphi_i : G_i \to H)$가 주어지면, 준동형들
$\varphi_i$의 상이 $H'$에 포함되고 기수가 $|\coprod G_i|$ 이하인 부분군
$H' \subset H$가 존재한다(여기서 쌍대곱은 군의 범주에서의 쌍대곱,
즉 $G_i$들의 자유곱이다). 구체적으로 $\varphi_i$들의 상들로 생성된
부분군에 유도위상을 주면 된다. 따라서 Categories, Lemma
\ref{categories-lemma-a-version-of-brown}의 가정들이 만족됨은 분명하고,
함자 $F$를 표현하는 위상군 $G$를 얻는다. 이는 정확히 $G$가 도식
$i \mapsto G_i$의 여극한이라는 뜻이다.

\medskip\noindent
군으로 가는 망각 함자와의 교환에 관한 주장을 보기 위해 Categories,
Lemma \ref{categories-lemma-adjoint-exact}를 사용하겠다. 실제로 망각 함자는
오른쪽 수반함자, 즉 군에 대응하는 무차별(또는 비이산) 위상군을 대응시키는
함자를 가진다.
\end{proof}

\begin{definition}
\label{topology-definition-topological-ring}
{\it 위상환}은 덧셈 $R \times R \to R$, $(x, y) \mapsto x + y$와 곱셈
$R \times R \to R$, $(x, y) \mapsto xy$가 연속이 되게 하는 위상을 갖춘
환 $R$이다. {\it 위상환의 준동형}은 연속인 환 준동형이다.
\end{definition}

\noindent
Stacks project에서 환은 $1$을 갖는 가환환이다. $R$이 위상환이면
$x \mapsto -x$가 연속이므로 $(R, +)$는 위상군이다. $R$이 위상환이고
$R' \subset R$이 부분환이면 유도위상을 갖춘 $R'$은 위상환이다.
$R$이 위상환이고 $R \to R'$이 환의 전사이면 몫위상을 갖춘 $R'$은
위상환이다.

\begin{lemma}
\label{topology-lemma-topological-ring-limits}
위상환의 범주는 극한을 가지며, 극한은 (a) 위상공간의 범주와 (b) 환의
범주로 가는 망각 함자들과 교환한다.
\end{lemma}

\begin{proof}
곱과 등화자에 대하여 존재성과 교환성을 증명하면 충분하다. Categories,
Lemma \ref{categories-lemma-limits-products-equalizers}를 보라.
$R_i$, $i \in I$를 위상환들의 모임이라 하자. 통상적인 곱
$R = \prod R_i$에 곱위상을 주자. 위상공간으로서
$R \times R = \prod (R_i \times R_i)$이므로(어느 범주에서나 곱은 곱과
교환한다) $R$ 위의 덧셈과 곱셈이 연속임을 알 수 있다.
$a, b : R \to R'$를 위상환들의 두 준동형이라 하자. 그러면 등화자로는
위상공간의 사상들로서 $a$와 $b$의 등화자를 그대로 택할 수 있고, 이는
유도위상을 갖춘 환의 사상들로서의 등화자와 같다.
\end{proof}

\begin{lemma}
\label{topology-lemma-topological-ring-colimits}
위상환의 범주는 여극한을 가지며, 여극한은 환의 범주로 가는 망각 함자와
교환한다.
\end{lemma}

\begin{proof}
Lemma \ref{topology-lemma-topological-group-colimits}의 증명에서 사용한 것과 정확히
같은 논증으로 여극한의 존재를 알 수 있다. 환으로 가는 망각 함자와의
교환에 관한 주장을 보기 위해 Categories, Lemma
\ref{categories-lemma-adjoint-exact}를 사용한다. 실제로 망각 함자는 오른쪽
수반함자, 즉 환에 대응하는 무차별(또는 비이산) 위상환을 대응시키는 함자를
가진다.
\end{proof}

\begin{definition}
\label{topology-definition-topological-module}
$R$을 위상환이라 하자. {\it 위상가군}은 덧셈
$M \times M \to M$과 스칼라곱 $R \times M \to M$이 연속이 되게 하는
위상을 갖춘 $R$-가군 $M$이다. {\it 위상가군의 준동형}은 연속인 가군
준동형이다.
\end{definition}

\noindent
$R$이 위상환이고 $M$이 위상가군이면 $x \mapsto -x$가 연속이므로
$(M, +)$는 위상군이다. $R$이 위상환이고 $M$이 위상가군이며
$M' \subset M$이 부분가군이면 유도위상을 갖춘 $M'$은 위상가군이다.
$R$이 위상환이고 $M$이 위상가군이며 $M \to M'$이 가군의 전사이면
몫위상을 갖춘 $M'$은 위상가군이다.

\begin{lemma}
\label{topology-lemma-topological-module-limits}
$R$을 위상환이라 하자. $R$ 위의 위상가군들의 범주는 극한을 가지며,
극한은 (a) 위상공간의 범주와 (b) $R$-가군의 범주로 가는 망각 함자들과
교환한다.
\end{lemma}

\begin{proof}
곱과 등화자에 대하여 존재성과 교환성을 증명하면 충분하다. Categories,
Lemma \ref{categories-lemma-limits-products-equalizers}를 보라.
$M_i$, $i \in I$를 $R$ 위의 위상가군들의 모임이라 하자. 통상적인 곱
$M = \prod M_i$에 곱위상을 주자. 위상공간으로서
$M \times M = \prod (M_i \times M_i)$이므로(어느 범주에서나 곱은 곱과
교환한다) $M$ 위의 덧셈이 연속임을 알 수 있다.
% P01-E1232: the frozen source calls the module action multiplication, not scalar multiplication.
곱셈 $R \times M \to M$도 마찬가지이다.
$a, b : M \to M'$를 $R$ 위의 위상가군들의 두 준동형이라 하자. 그러면
등화자로는 위상공간의 사상들로서 $a$와 $b$의 등화자를 그대로 택할 수
있고, 이는 유도위상을 갖춘 가군의 사상들로서의 등화자와 같다.
\end{proof}

\begin{lemma}
\label{topology-lemma-topological-module-colimits}
$R$을 위상환이라 하자. $R$ 위의 위상가군들의 범주는 여극한을 가지며,
여극한은 $R$ 위의 가군들의 범주로 가는 망각 함자와 교환한다.
\end{lemma}

\begin{proof}
Lemma \ref{topology-lemma-topological-group-colimits}의 증명에서 사용한 것과 정확히
같은 논증으로 여극한의 존재를 알 수 있다. $R$-가군으로 가는 망각 함자와의
교환에 관한 주장을 보기 위해 Categories, Lemma
\ref{categories-lemma-adjoint-exact}를 사용한다. 실제로 망각 함자는 오른쪽
수반함자, 즉 가군에 대응하는 무차별(또는 비이산) 위상가군을 대응시키는
함자를 가진다.
\end{proof}
